\documentclass[a4paper,11pt]{memoir}

\usepackage[utf8]{inputenc}
\usepackage[T1]{fontenc}
\usepackage[svgnames]{xcolor}
\usepackage{fouriernc}
\usepackage[frenchb]{babel}
% TODO \usepackage[french]{babelbib}
\usepackage{mathpartir}
\usepackage{syntax}
\usepackage{amssymb}
\usepackage{tikz}
\usepackage{amsmath}
\usepackage{float}
\usepackage{url}
\usepackage{xspace}
\usepackage{graphicx}
\usepackage[shortlabels]{enumitem}
\usepackage[backref=page]{hyperref}
\usepackage{links}
\usepackage[a4paper]{geometry}
\usepackage{amsthm}
\usepackage{stmaryrd}
\usepackage[normalem]{ulem}
\usepackage{wasysym}
\usepackage{multirow}
\usepackage{thmtools}
\usepackage{gramster}
\usepackage{switch}
\usepackage{pelagiccod}

\newcommand{\irule}[3]{ \inferrule*[right=(#1)]{#2}{#3} }

% Fix '::' -> ': :' when using babel
\newenvironment{grammarf}{\shorthandoff{:}\begin{grammar}}{\end{grammar}\shorthandon{:}}

\newcommand{\gramisa}[0]{::=~}
\newcommand{\gramor}[0]{|\hspace{3mm}}

\newcommand{\gramlr}[2]{%
  \begin{minipage}{0.2\textwidth}\raggedleft #1\end{minipage}
  \begin{minipage}{0.8\textwidth}#2\end{minipage}
}

% Types
\newcommand{\tInt}{\textsc{Int}\xspace}
\newcommand{\tFloat}{\textsc{Float}\xspace}
\newcommand{\tVoid}{\textsc{Void}\xspace}

\newcommand{\tFV}[1]{\textrm{FreeVars}(#1)}

\def\opxor{\mathbin{\char`\^}}

% Constant
\newcommand{\cNil}{\textsc{Nil}}

% Program & language names
\newcommand{\ctonewspeak}{\textsc{c2newspeak}}
\newcommand{\cpp}{\textsc{cpp}}
\newcommand{\newspeak}{\textsc{Newspeak}}

% Newspeak
\newcommand{\npkDecl}[2]{\textsc{Decl}~(#1) \{ #2 \}}
\newcommand{\npkIf}[3]{\textsc{If}(#1) \{ #2 \} \textsc{Else} \{ #3 \}}
\newcommand{\npkDoWith}[2]{\textsc{Do}\{ #1 \}\textsc{With}~#2:}
\newcommand{\npkWhile}[1]{\textsc{While}(1)\{#1\}}
\newcommand{\npkGoto}[1]{\textsc{Goto}~#1}
\newcommand{\npkReturn}[1]{\textsc{Return}~#1}

% Continuations
\def\kPass{\textsc{Kpass}}
\def\kSeq#1#2{\textsc{Kseq}(#1, #2)}
\def\kBlock#1#2{\textsc{Kblock}(#1, #2)}
\def\kGoto#1{\textsc{Kgoto}(#1)}
\def\kCall#1#2#3#4{\textsc{Kcall}(#1, #2, #3, #4)}
\def\kRet#1{\textsc{Kret}(#1)}

\newcommand{\insertcode}[1]{\input{gen/#1.pyg.tex}}

% \chapter* + add to toc
%\newcommand{\chapterstar}[1]{\chapter*{#1} \addcontentsline{toc}{chapter}{#1}}

% Register names
\def\defreg#1{\expandafter\def\csname#1\endcsname{\textsc{#1}\xspace}}

\defreg{eax}
\defreg{ebx}
\defreg{ecx}
\defreg{edx}
\defreg{esi}
\defreg{edi}
\defreg{ebp}
\defreg{esp}
\defreg{eip}
\defreg{eflags}

\def\crtrois{\textsc{CR3}}

\newcommand{\ring}[1]{\emph{ring}~#1}

\newcommand{\asminstr}[1]{\texttt{#1}}

\newcommand{\wip}[0]{%
\begin{center}
    \includegraphics[width=5cm]{gen/travaux.pdf}
\end{center}
}

\newcommand{\listcons}[0]{%
\texttt{(::)}%
}

\newcommand{\hexa}[1]{%
\texttt{0x#1}%
}

% Debug tikz bounding box
\def\drawbb{\draw[thick] (current bounding box.south west) rectangle (current bounding box.north east);}

% comes from the memoir manual {{{

\newcommand*{\addtotextheight}[1]{%
\twocolumn
\addtolength{\textheight}{#1}%
\onecolumn}

%% for main body, bottom of text at 1in, footer below
%% top of header at 1in, first text line double spaced
%% below base of header
\newlength{\linespace}
\setlength{\linespace}{\baselineskip} % current equivalent of \onelineskip
\setlength{\headheight}{\onelineskip}
\setlength{\headsep}{\linespace}
\addtolength{\headsep}{-\topskip}
\setlength{\uppermargin}{1in}
\addtolength{\uppermargin}{\headheight}
\addtolength{\uppermargin}{\headsep}

%% and for the bottom
\setlength{\lowermargin}{1in}
\setlength{\textheight}{\paperheight}

\addtolength{\textheight}{-\uppermargin}
\addtolength{\textheight}{-\lowermargin}
%% footnote settings
\setlength{\footskip}{\onelineskip}
\setlength{\footnotesep}{\onelineskip}


%% the fiddle lengths (..ta.. for title/approval page, others for prelims)
\newlength{\toptafiddle} \setlength{\toptafiddle}{2\linespace}
\newlength{\bottafiddle} \setlength{\bottafiddle}{0pt}
\newlength{\topfiddle}
\setlength{\topfiddle}{\toptafiddle}
\newlength{\botfiddle}
\setlength{\botfiddle}{\onelineskip}

%%% general macro for Abstract, etc., headings
\newcommand*{\pretoctitle}[1]{{\clearpage\centering \vspace*{-\topfiddle}#1\par}}

%%% Start the ACKNOWLEDGEMENTS
\newcommand{\asuacknowledgements}{\pretoctitle{ACKNOWLEDGEMENTS}}

%%% Start the ABSTRACT
\newcommand{\asuabstract}{ \addtotextheight{-\botfiddle} \pretoctitle{ABSTRACT}}

\addtodef{\mainmatter}{\addtotextheight{\botfiddle}{}}

%% make it easy to center any dedication
\newcommand{\asudedication}[1]{%
{\clearpage\mbox{}\vfill\centering #1 \par\vfill\clearpage}}

% }}} /memoir

\deflink{cil}{CIL - C Intermediate Language}{http://kerneis.github.com/cil/}
\deflink{cmm}{The C~-~- language}{http://www.cminusminus.org/}
\deflink{haskell}{Haskell Programming Language -- Official Website}{http://www.haskell.org/}
\deflink{ocaml}{OCaml -- Home}{http://ocaml.org/}
\deflink{penjili}{Penjili project}{http://www.penjili.org/}
%\deflink{perl}{Perl Programming Language -- Official Website}{http://www.perl.org/}
\deflink{python}{Python Programming Language -- Official Website}{http://www.python.org/}
\deflink{rust}{The Rust Programming Language}{http://www.rust-lang.org/}
\deflink{sparse}{Sparse - a Semantic Parser for {C}}{https://sparse.wiki.kernel.org/index.php/Main_Page}


\input{chapstyle.tex}

\def\disprule#1{
  \switch{#1}
    {Cst-Int}{ \irule{Cst-Int}
                 { }
                 { Γ ⊢ i : \tInt}
    }
    {Cst-Float}{ \irule{Cst-Float}
                   { }
                   { Γ ⊢ d : \tFloat}
    }
    {Cst-Null}{ \irule{Cst-Null}
                  { }
                  { Γ ⊢ \eNull : t*}
    }
    {Cst-Unit}{ \irule{Cst-Unit}
                  { }
                  { Γ ⊢ \eUnit : \tUnit }
    }
    {Lv-Var}{ \irule{Lv-Var}
                { x : t ∈ Γ }
                { Γ ⊢ x : t }
    }
    {Lv-Deref}{ \irule{Lv-Deref}
                  { Γ ⊢ lv : t* }
                  { Γ ⊢ *lv : t }
    }
    {Lv-Index}{ \irule{Lv-Index}
                  { Γ ⊢ e : \tInt \\
                    Γ ⊢ lv : t[]
                  }
                  { Γ ⊢ lv[e] : t }
    }
    {Lv-Field}{ \irule{Lv-Field}
                  { (s, l, t_l) ∈ S \\
                    Γ ⊢ lv : \tstruct{s}
                  }
                  { Γ ⊢ lv.l : t_l }
    }
    {Op-Int}{ \irule{Op-Int}
                { \opbin ∈ \{+,-,\times,/,\&,|,\opxor,\&\&,||,\ll,\gg \} \\
                  Γ ⊢ e_1 : \tInt \\
                  Γ ⊢ e_2 : \tInt
                }{
                  Γ ⊢ e_1~\opbin~e_2 : \tInt
                }
    }
    {Op-Float}{ \irule{Op-Float}
                { \opbin ∈ \{+.,-.,\times.,/.\} \\
                  Γ ⊢ e_1 : \tFloat \\
                  Γ ⊢ e_2 : \tFloat
                }{
                  Γ ⊢ e_1~\opbin~e_2 : \tFloat
                }
    }
    {Eq-Num}{ \irule{Eq-Num}
                {t ∈ \{ \tInt, \tFloat \}}
                {\textsc{Eq}(t)}
    }
    {Eq-Ptr}{ \irule{Eq-Ptr}
                {\textsc{Eq}(t)}
                {\textsc{Eq}(t*)}
    }
    {Eq-Array}{ \irule{Eq-Array}
                  {\textsc{Eq}(t)}
                  {\textsc{Eq}(t[])}
    }
    {Op-Eq}{ \irule{Op-Eq}
             { \opbin ∈ \{=,≠\} \\
               Γ ⊢ e_1 : t \\
               Γ ⊢ e_2 : t \\
               \textsc{Eq}(t)
             }{
               Γ ⊢ e_1~\opbin~e_2 : \tInt
             }
    }
    {Op-Comparable}{ \irule{Op-Comparable}
                       { \opbin ∈ \{=,≠,≤,≥,<,>\} \\
                         Γ ⊢ e_1 : t \\
                         Γ ⊢ e_2 : t \\
                         t ∈ \{ \tInt, \tFloat \}
                       }{
                         Γ ⊢ e_1~\opbin~e_2 : \tInt
                       }
    }
    {Unop-Plus-Int}{ \irule{Unop-Plus-Int}
                       { Γ ⊢ e : \tInt }
                       { Γ ⊢ + e : \tInt }
    }
    {Unop-Plus-Float}{ \irule{Unop-Plus-Float}
                         { Γ ⊢ e : \tFloat }
                         { Γ ⊢ +. e : \tFloat }
    }
    {Unop-Minus-Int}{ \irule{Unop-Minus-Int}
                        { Γ ⊢ e : \tInt }
                        { Γ ⊢ - e : \tInt }
    }
    {Unop-Minus-Float}{ \irule{Unop-Minus-Float}
                          { Γ ⊢ e : \tFloat }
                          { Γ ⊢ -. e : \tFloat }
    }
    {Unop-Not}{ \irule{Unop-Not}
                { \opun ∈ \{ \sim, ! \} \\
                  Γ ⊢ e : \tInt
                }{
                  Γ ⊢ \opun~e : \tInt
                }
              }
    {Ptr-Arith}{ \irule{Ptr-Arith}
                   { \opbin ∈ \{ +_p, -_p \} \\
                     Γ ⊢ e_1 : t* \\
                     Γ ⊢ e_2 : \tInt \\
                   }
                   { Γ ⊢ e_1~\opbin~e_2 : t* }
               }
    {Addr}{ \irule{Addr}
              { Γ ⊢ lv : t }
              { Γ ⊢ \&lv : t* }
          }
    {Set}{ \irule{Set}
             { Γ ⊢ lv : t \\
               Γ ⊢ e : t
             }
             { Γ ⊢ lv ← e : t }
         }
    {Array}{ \irule{Array}
               { ∀ i ∈ [1;n] , Γ ⊢ e_i : t }
               { Γ ⊢ \{ e_1 ;
                       … ; e_n \}
                       : t[]
               }
           }
    {Struct}{ \irule{Struct}
                { ∀ i ∈ [1;n], Γ ⊢ e_i : t_i \\
                  ∀ i ∈ [1;n], (s, l_i, t_i) ∈ S
                }
                { Γ ⊢ \{ l_1 : e_1 ;
                     … ; l_n : e_n \}
                    : \tstruct{s}
                }
            }
    {Call}{ \irule{Call}
              { Γ ⊢ e : ( t_1 ,
                         … , t_n ) → t \\
                ∀ i ∈ [1;n], Γ ⊢ e_i : t_i
              }
              { Γ ⊢ e ( e_1
                         , …
                         , e_n
                         )
                      : t
              }
          }
    {Pass}{ \irule{Pass}
              { }
              {Γ ⊢ \iPass}
          }
    {Seq}{ \irule{Seq}
             { Γ ⊢ i_1 \\
               Γ ⊢ i_2
             }{
               Γ ⊢ i_1;i_2
             }
         }
    {Exp}{ \irule{Exp}
             { Γ ⊢ e : t }
             { Γ ⊢ e }
         }
    {If}{ \irule{If}
            { Γ ⊢ e : \tInt \\
              Γ ⊢ i_1 \\
              Γ ⊢ i_2
            }
            { Γ ⊢ \iIf{e}{i_1}{i_2} }
        }
    {While}{ \irule{While}
               { Γ ⊢ e : \tInt \\
                 Γ ⊢ i
               }
               { Γ ⊢ \iWhile{e}{i} }
           }
    {Return}{ \irule{Return}
                { Γ ⊢ \vRet ← e }
                { Γ ⊢ \iReturn{e} }
            }
    {Fun}{ \irule{Fun}
             { Γ' = (Γ - \vRet), a_1 : t_1, …, a_n : t_n,
                    l_1 : t'_1, …, l_n : t'_p,
                    \vRet : t \\
               ∀ i ∈ [1;p], Γ ⊢ e_i : t'_i \\
               Γ' ⊢ i
             }
             { Γ ⊢ \mathrm{fun} (a_1, …, a_n)
                   ((l_1, e_1), …, (l_p, e_p)) \{i\}
                    : (t_1, …, t_n) → t
             }
         }
    {Ph-Exp}{ \irule{Ph-Exp}
                { Γ ⊢ e : t }
                { \typh{Γ}{e}{Γ} }
            }
    {Ph-Var}{ \irule{Ph-Var}
                { Γ ⊢ e : t \\
                  Γ' = (x, t), Γ
                }
                { \typh{Γ}{p}{Γ'} }
            }
    {Ph-Struct}{ \irule{Ph-Struct}
                   { S' =
                     (x_1, t_1, s),
                     …
                     (x_n, t_n, s)
                   }
                   { \typh{Γ}{ \tstruct{s}
                                  \{ x_1 : t_1
                                   ; …
                                   ;  x_n : t_n
                                  \}
                              }{Γ}
                   }
               } % TODO
    {END}
}

\def\semrule#1{
  \switch{#1}
    {Ctx}{ \irule{Ctx}
             { \mm{m}{e}{m'}{e'} }
             { \mm{m}{\ctxSub{C}{e}}{m'}{\ctxSub{C}{e'}} }
    }
    {Ctx-Lv}{ \irule{Ctx-Lv}
             { \mmlv{m}{lv}{m'}{lv'} }
             { \mmlv{m}{\ctxSubL{C_L}{lv}}{m'}{\ctxSubL{C_L}{lv'}} }
    }
    {Ctx-Instr}{ \irule{Ctx-Instr}
             { \mmlv{m}{i}{m'}{i'} }
             { \mmlv{m}{\ctxSubI{C_I}{i}}{m'}{\ctxSubI{C_I}{i'}} }
    }
    {Phi-Var}{ \irule{Phi-Var}
                 { a = \mathrm{Lookup}(x, m) }
                 {
                   \mmlv{m}{x}{m}{a}
                 }
    }
    {Phi-Deref}{ \irule{Phi-Deref}
                   { }
                   { \mms{*φ}{\widehat{*}φ} }
    }
    {Phi-Struct}{ \irule{Phi-Struct}
                    {  }
                    { \mms{lv.l}{lv\widehat{.}l} }
    }
    {Phi-Array}{ \irule{Phi-Array}
                   { }
                   { \mms{φ[n]}{φ\widehat{[n]}} }
    }
    {Exp-Cst}{ \irule{Exp-Cst}
                 { }
                 {\mms{c}{\widehat{c}}}
             }
    {Exp-Fun}{ \irule{Exp-Fun}
                 { }
                 {\mms{f}{\widehat{f}}}
             }
    {Exp-Lv}{ \irule{Exp-Lv}
              { }
              {\mm{m}{φ}{m}{m[φ]_Φ}}
    }
    {Exp-UnOp}{ \irule{Exp-UnOp}
                { }
                {\mms{\opun~v}{\widehat{\opun}~v}}
    }
    {Exp-BinOp}{ \irule{Exp-BinOp}
                   { \opbin ∉ \{/, \%\} }
                   {\mms{v_1~\opbin~v_2}{v_1~\widehat{\opbin}~v_2}}
    }
    {Exp-Div}{ \irule{Exp-Div}
                 { \opbin ∈ \{/, \%\} \\
                   v_2 ≠ \widehat{0}
                 }
                 {\mms{v_1~\opbin~v_2}{v_1~\widehat{\opbin}~v_2}}
    }
    {Exp-Div-Zero}{ \irule{Exp-Div-Zero}
                 { \opbin ∈ \{/, \%\}
                 }
                 {\msi{m}{v_1~\opbin~0} → Ω_{div}}
    }
    {Exp-AddrOf}{ \irule{Exp-AddrOf}
                    { a = \mathrm{Lookup} (x, m)
                    }
                    {\mm{m}{\&x}{m}{a}}
    }
    {Exp-Set}{ \irule{Exp-Set}
                 { }
                 {\mm{m}{φ ← v}{m[φ ← v]}{v}}
    }
    {Exp-Struct}{ \irule{Exp-Struct}
                   { }
                   { \mms{ \{ l_1 : v_1 ; … ; l_n : v_n \} }
                         { \widehat{\{ l_1 : v_1 ; … ; l_n : v_n \}} }
                   }
    }
    {Exp-Array}{ \irule{Exp-Array}
                 { }
                 {
                   \mms{ [ v_1, … , v_n ] }
                       { \widehat{[ v_1, … , v_n ]} }
                 }
    }
    {Exp-Call}{ \irule{Exp-Call}
                  { f =
                    \mathrm{fun} (a_1, …, a_n)
                                 ((l'_1, e'_1), …, (l'_p, e'_p))
                                 \{ i \} \\
                    m_1 = \mathrm{Push}(m_0, ((a_1, v_1), …, (a_n, v_n))) \\
                    \evstar{m_1}{e'_1 \\ \vdots \\ e'_p}
                           {m_2}{v'_1 \\ \vdots \\ v'_p} \\
                    m_3 = \mathrm{Extend}(m_2, ((l_1, v_1), …, (l_n, v_n))) \\
                    \mm{m_3}{i}{m_4}{\iReturn{v}} \\
                    m_5 = \mathrm{Pop}(m_4) \\
                    m_6 = \mathrm{Cleanup}(m_5) \\
                  }
                  {\mm{m_0}{f (v_1, …, v_n)}{m_6}{v}}
    }
    {Seq}{ \irule{Seq}
             { \mmi{m}{i}{m'}{\iPass} }
             { \mmi{m}{(i;i')}{m'}{i'} }
    }
    {Pass}{ \irule{Pass}
              { }
              { \mmi{m}{(\iPass;i)}{m}{i} }
    }
    {Exp}{ \irule{Exp}
             { }
             { \mmi{m}{v}{m}{\iPass} }
    }
    {If-False}{ \irule{If-False}
                  { }
                  { \mmi{m}{\iIf{0}{i_t}{i_f}}{m}{i_f} }
    }
    {If-True}{ \irule{If-True}
                 { v ≠ 0 }
                 { \mmi{m}{\iIf{v}{i_t}{i_f}}{m}{i_t} }
    }
    {While}{ \irule{While}
               { }
               { \mmi{m}{\iWhile{e}{i}}{m}{\iThen{e}{i;\iWhile{e}{i}}}
               }
    }
    {Return}{ \irule{Return}
              { }
              { \mmi{m}{\iReturn{v};i}{m}{\iReturn{e}} }
    }
    {Ph-Struct}{ \irule{Ph-Struct}
                   { }
                   {\ph{m}{\tstruct{s} \{ … \}}{m}}
    }
    {Ph-Exp}{ \irule{Ph-Exp}
                {\mm{m}{e}{m'}{v}}
                {\ph{m}{e}{m'}}
            }
    {Ph-Var}{ \irule{Ph-Var}
                {\mm{m}{e}{m'}{v}}
                {\ph{(s,g)}{x = e}{(s, (x, v)::g)}}
    }
    {Ph*-Nil}{ \irule{Ph*-Nil}
                 { }
                 {\phstar{m}{[]}{m}}
    }
    {Ph*-Cons}{ \irule{Ph*-Cons}
                  {\ph{m}{p}{m'} \\
                   \phstar{m'}{ps}{m''}}
                  {\phstar{m}{p::ps}{m''}}
    }
    {END}
}

\def\figctx{
  \gramlr{Contextes}{
  \begin{align*}
  C \gramisa & C_L
  \\ \gramor & \ctxOp{C}{e}
  \\ \gramor & \ctxOp{v}{C}
  \\ \gramor & \ctxUnOp{C}
  \\ \gramor & \ctxSet{C}{e}
  \\ \gramor & \ctxSet{φ}{C}
  \\ \gramor & \{ l_1:v_1 ; … ; l_i:C ; … ; l_n:e_n \}
  \\ \gramor & [ v_1 ; … ; C ; … ; e_n ]
  \\ \gramor & C (e_1, …, e_n)
  \\ \gramor & f (v_1, …, C, …, e_n)
  \\
  \\
  C_L \gramisa & \ctxEmpty
  \\   \gramor & \ctxLvDeref{C_L}
  \\   \gramor & \ctxLvField{C_L}{l}
  \\   \gramor & \ctxLvIndex{C_L}{e}
  \\   \gramor & \ctxLvIndex{φ}{C}
  \\
  \\
  C_I \gramisa & C_I;i
  \\   \gramor & \iIf{C}{i_1}{i_2}
  \\   \gramor & \iReturn{C}
  \\   \gramor & C
  \end{align*}
  }
}


\usetikzlibrary{calc}
\usetikzlibrary{shapes}
\usetikzlibrary{arrows}
\usetikzlibrary{decorations.pathreplacing}

\usepackage{fancyvrb}
\input{gen/pygments-style.tex}

\floatstyle{boxed}
\restylefloat{figure}
\newsubfloat{figure}

\epigraphfontsize{\small\itshape}
\setlength\epigraphwidth{8cm}
\setlength\epigraphrule{0pt}
\def\epigraph#1#2{}

\setlist[itemize]{label=$\bullet$}

\maxtocdepth{subsection}    % What appears in toc
\maxsecnumdepth{subsection} % What get numbered

\def\partintro{}
\renewcommand{\afterpartskip}{\vfil \partintro \vfil \newpage}

%\theoremstyle{definition}
\newtheorem{definition}{Définition}[chapter]
\newtheorem{theorem}{Théorème}[chapter]
\newtheorem{lemma}{Lemme}[chapter]
\newtheorem{example}{Exemple}[chapter]

\setlength{\headsep}{5mm}
\addtolength{\textheight}{-1cm}



\title{Analyse de sécurité de logiciels système par typage statique}
\subtitle{Application au noyau Linux}
\author{Étienne Millon}
\date{xx-yy-2013}

\usepackage{url}
\ifpdf
\usepackage{pdfcolmk}
\fi
%% check if using xelatex rather than pdflatex
\ifxetex
\usepackage{fontspec}
\fi
\usepackage{graphicx}
%%\usepackage{hyperref}
%% drawing package
\usepackage{tikz}
%% for dingbats
\usepackage{pifont}
\providecommand{\HUGE}{\Huge}% if not using memoir
\newlength{\drop}% for my convenience
%% specify the Webomints family
\newcommand*{\wb}[2]{\fontsize{#1}{#2}\usefont{U}{webo}{xl}{n}}
%% select a (FontSite) font by its font family ID
\newcommand*{\FSfont}[1]{\fontencoding{T1}\fontfamily{#1}\selectfont}
%% if you don’t have the FontSite fonts either \renewcommand*{\FSfont}[1]{}
%% or use your own choice of family.
%% select a (TeX Font) font by its font family ID
\newcommand*{\TXfont}[1]{\fontencoding{T1}\fontfamily{#1}\selectfont}
%% Generic publisher’s logo
\newcommand*{\plogo}{\fbox{$\mathcal{PL}$}}

%% Some shades
\definecolor{Dark}{gray}{0.2}
\definecolor{MedDark}{gray}{0.4}
\definecolor{Medium}{gray}{0.6}
\definecolor{Light}{gray}{0.8}
%%%% Additional font series macros
\makeatletter
%%%% light series
%% e.g., kernel doc, section s: line 12 or thereabouts
\DeclareRobustCommand\ltseries
{\not@math@alphabet\ltseries\relax
\fontseries\ltdefault\selectfont}
%% e.g., kernel doc, section t: line 32 or thereabouts
\newcommand{\ltdefault}{l}
%% e.g., kernel doc, section v: line 19 or thereabouts
\DeclareTextFontCommand{\textlt}{\ltseries}
% heavy(bold) series
\DeclareRobustCommand\hbseries
{\not@math@alphabet\hbseries\relax
\fontseries\hbdefault\selectfont}
\newcommand{\hbdefault}{hb}
\DeclareTextFontCommand{\texthb}{\hbseries}
\makeatother


\usetikzlibrary{fadings}

\newcommand*{\titleUL}{
  \newgeometry{left=2cm, right=2cm, top=2cm, bottom=2cm}%
  \begingroup% University of Liege
\drop=0.1\textheight
\begin{center}
{\textsc{Université Pierre et Marie Curie\\
École Doctorale Informatique, Télécommunications et Électronique
}}\\
\vspace{1cm}
\begin{tikzpicture}
  \begin{scope}
    \node (A) {\includegraphics[height=\textwidth]{curies.pdf}};
    \path[opacity=0.2, fill=white]
      (A.north west) rectangle (A.south east);
    %\path[opacity=0.8, fill=white, path fading=north, fading transform={rotate=-45}]
    %(A.north west) rectangle (A.south east);
    %\clip (A.center) circle(8.5cm);
    %\shade[opacity=0.2, shading=axis,bottom color=black!10,top color=black,shading angle=-45]
    %(A.south west) rectangle (A.north east);
  \end{scope}
\end{tikzpicture}
\vspace{-13cm}
\\[\drop]
\begin{minipage}{13cm}
  \begin{center}
\rule{\textwidth}{1pt}\par
\vspace{0.5\baselineskip}
{\huge\bfseries \thetitle\\
\vspace{5mm}
\large --- \thesubtitle ---
}\\[0.5\baselineskip]
\rule{\textwidth}{1pt}
  \end{center}
\end{minipage}
\par
\vspace{1cm}
{{\Large\textsc{\theauthor}} \\
sous la direction d'Emmanuel Chailloux et de Sarah Zennou
}
\vspace{16mm}\\
{{\LARGE\textsc{Thèse}}\\
pour obtenir le titre de\\
{\Large Docteur en Sciences}\\
{\large mention Informatique}
}
\vspace{2cm}\\
{\large Soutenue le \thedate{} devant un jury composé de\\

aaa\\
aaa\\
aaa\\
aaa\\
aaa
}
\end{center}
\endgroup
\restoregeometry
}


\begin{document}

\frontmatter

\begin{titlingpage}
\titleUL{}
\end{titlingpage}

\cleartorecto

\asuacknowledgements

% TODO

% Directeurs

% Rapporteurs

% Autres membres

% Ecole doc

% Equipe enseignante

% APR

% Financement

% 325

% Potes

% Famille

% Anaïs


\clearpage

%\asudedication{Dédicace}

\tableofcontents

\mainmatter

\chapter{Introduction}

\label{cha:intro}
Communication,
audiovisuel,
transports,
médecine :
tout ces domaines se sont transformés dans les dernières décennies,
en particulier grâce à la révolution numérique.
En effet le plus petit appareil électrique contient maintenant des composants
matériels programmables.

% TODO faire une intro "la vie est numérique" avec du Berry @ CdF

En 2013, on pense bien sûr aux téléphones portables dont la fonctionnalité et la
complexité les rapprochent des ordinateurs de bureau. Par exemple, le système
d'exploitation Android de Google est basé sur le noyau Linux, destiné à la base
aux micro-ordinateurs.

Le noyau d'un système d'exploitation est chargé de faire l'intermédiaire entre
le matériel (processeur, mémoire, périphériques, …) et les applications
exécutées sur celui-ci (par exemple un navigateur web, une calculatrice ou un
carnet d'adresses). Il doit aussi garantir la sécurité et l'isolation de
celles-ci.

En tant qu'intermédiaire de confiance, le noyau a un certain nombre de
responsabilités et est le seul à avoir accès à certaines informations sensibles.
Il est capital de s'assurer qu'il est bien le seul à pouvoir y accéder. En
particulier, il faut pouvoir vérifier que les requêtes faites par l'utilisateur
au noyau ne peuvent pas détourner ce dernier et lui faire fuiter des
informations confidentielles.

Le problème est que comme tous les logiciels, les noyaux de système
d'exploitation sont écrits par des humains, qui ne sont pas parfaits. Loin de
là: on estime qu'avant relecture, 1000 lignes de code contiennent entre 5 et 100
erreurs de programmation en moyenne.
% TODO sourcer
Les activités de relecture et de
débogage ont beau prendre la majeure partie du temps de développement, il est
facile de laisser passer des défauts de programmation.

Une technique efficace est de réaliser des tests, c'est-à-dire exécuter le code
du programme sous un environnement contrôlé. On peut alors détecter des
comportements non désirés. Mais même avec des tests les plus exhaustifs il n'est
pas possible de couvrir tous les cas d'utilisation.

Parmi les techniques communes pour éliminer les ``bugs'', on peut ajouter à un
langage de programmation un \emph{système de types statiques}.

Celui-ci permet d'associer à chaque fragment d'un programme, une étiquette
décrivant quel genre de valeur sera produite par son évaluation. Cette analyse,
faite à la compilation, permet de rejeter les programmes incorrects du point de
vue de cette catégorisation.

Le but de cet thèse est de montrer que le typage statique peut être utilisé pour
détecter et interdire certains bugs présents dans les noyaux de systèmes
d'exploitation.

\section{Rôle d'un système d'exploitation}

On décrit ici la fonction générale d'un système d'exploitation. Ce sujet est
détaillé en profondeur, y compris avec des cas d'étude, dans~\cite{tanenbaum}.

Un ordinateur est constitué de nombreux composants matériels : microprocesseur,
mémoire, et divers périphériques. Pourtant, au niveau de l'utilisateur, des
dizaines de logiciels permettent d'effectuer toutes sortes de calculs et de
communications. Le système d'exploitation permet de faire l'interface entre ces
niveaux d'abstraction.

Au cours de l'histoire des systèmes informatiques, la manière de les programmer
a beaucoup évolué. Au départ, les programmeurs avaient accès au matériel dans
son intégralité : toute la mémoire pouvait être accédée, toutes les instructions
pouvaient être utilisées.

Néanmoins c'est un peu restrictif, puisque cela ne permet qu'à une personne
d'interagir avec le système. Dans la seconde moitié des années 60, sont apparus
les premiers systèmes ``à temps partagé'', permettant à plusieurs utilisateurs
de travailler en même temps.

Permettre l'exécution de plusieurs programmes en même temps est une idée
révolutionnaire, mais elle n'est pas sans difficultés techniques : en effet les
ressources de la machine doivent être aussi partagées entre les utilisateurs et
les programmes. Par exemple, plusieurs programmes vont utiliser le processeur
les uns à la suite des autres ; et chaque programme aura à sa disposition une
partie de la mémoire principale, ou du disque dur.

Si deux programmes (ou plus) s'exécutent de manière concurrente sur le même
matériel, il faut s'assurer que l'un ne puisse pas écrire dans la mémoire de
l'autre, ou que les deux n'utilisent pas la carte réseau en même temps. Ce sont
des rôles du système d'exploitation.

Ainsi, au lieu d'accéder directement au matériel via des instructions de bas
niveau, les programmes communiquent avec le noyau, qui centralise donc les
appels au matériel, et abstrait certaines opérations.

Par exemple, comparons ce qui se passe concrètement lors de la copie de données
depuis un cédérom ou une mémorette USB.

\begin{itemize}

  \item Dans le cas du cédérom, il faut interroger le bus SATA, interroger le
    lecteur sur la présence d'un disque dans le lecteur, activer le moteur,
    calculer le numéro de trame des données sur le disque, demander la lecture,
    puis déclencher une
    copie de la mémoire.

  \item Avec une mémorette USB, il faut interroger le bus USB, rechercher le bon
    numéro de périphérique, le bon numéro de canal dans celui-ci, lui appliquer
    une commande de lecture au bon numéro de bloc, puis copier la mémoire.

\end{itemize}

Ces deux opérations, bien qu'elles aient le même but (copier de la mémoire
depuis un périphérique amovible), ne sont pas effectuées de la même manière.
C'est pourquoi le système d'exploitation fournit les notions de fichier,
lecteur, etc : le programmeur n'a plus qu'à utiliser des commandes de haut
niveau (``monter un lecteur'', ``ouvrir un fichier'', ``lire dans un fichier'')
et selon le type de lecteur, le système d'exploitation effectuera les actions
appropriées.

En résumé, un système d'exploitation est l'intermédiaire entre le logiciel et
le matériel, et en particulier assure les rôles suivants :

% TODO À affiner ou supprimer

\paragraph{Gestion des processus}

Un système d'exploitation peut permettre d'exécuter plusieurs programmes à la
fois. Il faut alors orchestrer ces différents processus et les séparer en terme
de temps et de ressources partagées.

\paragraph{Gestion de la mémoire}

Chaque processus, en plus du noyau, doit disposer d'un espace mémoire différent.
C'est-à-dire qu'un processus ne doit pas pouvoir interférer avec un autre.

\paragraph{Gestion des fichiers}

Les processus peuvent accéder à une hiérarchie de fichiers, indépendamment de la
manière d'y accéder.

\paragraph{Gestion des périphériques}

Le noyau étant le seul code ayant des privilèges, c'est lui qui doit communiquer
avec les périphériques matériels.

\paragraph{Abstractions}

Le noyau fournit aux programmes une interface unifiée, permettant de stocker des
informations de la même manière sur un disque dur ou une clef USB (alors que
l'accès se déroulera de manière très différente en pratique).

\section{Noyaux de systèmes d'exploitation}

% TODO [E] ??? privilèges ?

Les programmes qui composent un système d'exploitation sont exécutés avec un
niveau de privilège différent, correspondant à un jeu d'instructions plus ou
moins important. Cette vérification est faite directement pas le processeur,
d'une manière matérielle. En plus de cette distinction, les programmes ont une
vue différente de la mémoire. Si un programme utilisateur avait accès aux
structures de données internes du noyau, il pourrait obtenir des informations
auquel il n'a pas normalement accès, ou obtenir des privilège supplémentaires
(par exemple en obtenant les privilèges de l'administrateur sur la machine).

C'est pourquoi les noyaux stockent leur code et données dans un endroit en
mémoire qui n'est pas accessible aux programmes utilisateur. Sur un système
Intel 32 bits, cela correspond aux adresses supérieures à \texttt{0xc0000000},
ou 3 Gio.

% TODO trad

Because user-mode programs cannot run low-level instructions, they cannot do
anything interesting on their own. They have to use the kernel's facilities to
perform input/output or manipulate the operating system's abstractions (such as
creating a new process). To this end, they can communicate with the kernel
via the narrow interface of \emph{system calls}.

For example, every time a user process calls the \texttt{read()} function, an
architecture-specific mechanism invokes the kernel-mode function
\texttt{sys\_read()}. The user passes a file descriptor, a pointer and an amount
of bytes to fill and the kernel will actually perform input/output and copy the
resulting bytes on the user's buffer.

But if the user supplies a pointer whose value is in the kernel's reserved area,
this zone will be overwritten, potentially leading to a security breach. In a
sense, the kernel has been abused because it accessed the memory with its own
privileges instead of the originator's. This is known as the \emph{confused
deputy problem}\cite{hardy88confused}.

When implementing a system call, it is thus necessary to forbid direct
dereference of a pointer whose value can be controlled by the user. In case of a
system call with a pointer argument, one has to dynamically verify that the
address provided by the user lies within the process's memory space.

% TODO [E] est-ce ici que tu veux le placer ? rien ne t'empêche de le dire
% plusieurs fois

The object of this work is to detect the places where it is necessary to insert
dynamic checks. To this end, we augment the type system of a subset of the C
programming language with \emph{type qualifiers}: instead of having a single
pointer type, we add syntactic rules to distinguish between safe pointers that
can be dereferenced with the \texttt{*}, \texttt{->} and \texttt{[]} operators ;
and unsafe pointers that have to go through a special API provided by the
kernel, consisting most notably of the two functions \texttt{copy\_from\_user}
and \texttt{copy\_to\_user}.

\section{Systèmes de types}

La plupart des langages de programmation incorporent la notion de type, qui
permet de détecter ou d'empêcher de manipuler des données incompatibles entre
elles.

Nous avons vu dans le chapitre~\ref{cha:os} qu'au niveau du langage machine, les
seules données qu'un ordinateur manipule sont des nombres. Selon les opérations
effectuées, ils seront interprétés comme des entiers, des adresses mémoires, ou
des caractères. Pourtant il est clair que certaines opérations n'ont pas de
sens: par exemple, ajouter deux adresses, ou déréférencer le résultat d'une
division sont des comportements qu'on voudrait pouvoir empêcher.

En un mot, le but du typage est de classifier les objets et de restreindre les
opérations possibles selon la classe d'un objet : ``ne pas ajouter des pommes et
des oranges''. Le modèle qui permet cette classification est appelé
\emph{système de types} et est en général constitué d'un ensemble de
\emph{règles de typage}, comme ``un entier plus un entier égale un entier''.

Il y a deux grandes familles de systèmes de types, selon quand se fait la
vérification de types. On peut en effet l'effectuer au moment de l'exécution, ou
au contraire prévenir les erreurs à l'exécution en la faisant au moment de la
compilation (ou avant l'interprétation).

\paragraph{Typage dynamique :} dans ce cas, chaque valeur manipulée par le
programme est décorée d'une étiquette définissant comment interpréter la valeur
en question. Les règles de typage sont alors réalisées à l'exécution. Par
exemple, l'opérateur "$+$" vérifie que ces deux opérandes ont une étiquette
"entier", et construit alors une valeur obtenue en faisant l'addition des deux
valeurs, avec une étiquette "entier". Symboliquement :

\begin{verbatim}
fonction addition(x, y) {
  if ( (etiquette_type(x) != type_entier)
    || (etiquette_type(y) != type_entier)
    ) {
      erreur ("L'addition est seulement définie pour deux entiers")
  } else {
    vx = valeur(x)
    vy = valeur(y)
    v = addition_machine(vx, vy)
    return Valeur(valeur=v, type=type_entier)
  }
}
\end{verbatim}

Le langage Python~\link{python} utilise cette stratégie, qui est illustrée dans
la session interactive de la figure~\ref{fig:typage-dynamique}.

\begin{figure}
  \insertcode{typage-dynamique.pycon}
  \caption{Session Python présentant le typage dynamique}
  \label{fig:typage-dynamique}
\end{figure}

L'analyse de taintage consiste à ajouter des étiquettes au données décrivant
leur provenance, de manière à ce que chaque accès dangereux soit vérifié à
l'exécution. Cela se marie bien avec les langages typés dynamiquement, comme le
montre l'exemple célèbre du mode ``souillé'' (\emph{tainted}) de
Perl~\cite{perlCamelBook}. Il est aussi possible, mais plus délicat, d'appliquer
cette technique à des langages compilés~\cite{clause-etal-issta07,oakland10}.
% TODO [E]: à préciser

\paragraph{Typage statique :} dans ce cas on fait les vérifications à la
compilation. En quelque sorte, l'approche dynamique est pessimiste, puisqu'elle
demande de traiter très souvent le cas où les types ne sont pas corrects.
Intuitivement, dans le cas où toutes les fonctions se comportent bien, faire la
vérification est inutile. Pour vérifier ceci, on donne à chaque fonction un
contrat comme "si deux entiers sont passés, le fonction renverra un entier sans
provoquer d'erreur". Cet ensemble de contrats peut être vérifié statiquement par
le compilateur, à l'aide d'un système de types statique.

% TODO [E]: division par zéro

Ainsi la fonction "$+$" est typée \texttt{(int * int)} $→$ \texttt{int}. Les
règles permettant de vérifier le typage sont par exemple les suivantes :

\begin{itemize}
\item
  une constante entière est toujours de type \texttt{int}.
\item
  si $f$ a pour type $(t_1, …, t_n) → t$, et que chaque $e_i$ a pour type $t_i$
  ($i ∈ [1;n]$), alors $f(e_1, …, e_n)$ a pour type $t$.
\item
  si en considérant que chaque $e_i$ a pour type $t_i$ ($i ∈ [1;n]$),
  on arrive à typer le corps de $f$ et que sa valeur de retour a alors
  pour type $t$,
  alors $f$ a pour type $(t_1, …, t_n) → t$.
\end{itemize}

Cet ensemble de règles, une fois formalisé et implanté, est en général assez
efficace pour vérifier la correction d'un programme.
% TODO [E] ?

\paragraph{Typage fort ou faible :}

Si un système de types statique permet d'éliminer totalement la nécessité de
réaliser des tests de typage, on dit qu'il est \emph{fort}. Mais ce n'est que
rarement le cas. En effet, il peut y avoir des constructions au sein du langage
qui permettent de contourner le système de types, comme un opérateur de
transtypage\ref{fig:javacast}. À l'exécution, une erreur de types est levée :

\begin{figure}
  \insertcode{cast.java}
  \caption{Transtypage en Java}
  \label{fig:javacast}
\end{figure}

\begin{Verbatim}
Exception in thread "main" java.lang.ClassCastException:
    java.lang.Integer cannot be cast to java.lang.Float
        at Cast.main(Cast.java:5)
\end{Verbatim}

\paragraph{Polymorphisme :} parfois, il est trop restrictif de donner un unique
contrat à une fonction. Quel doit être le type d'une fonction ajoutant un
élément à une liste ?

En première approximation, on peut imaginer fournir une version du code par type
de données à manipuler. C'est la solution retenue par les premières versions du
langage Pascal, ce qui rendait très difficile l'écriture de
bibliothèques\cite{PascalNoEscape}.

Une amélioration peut être de générer à la compilation une fonction qui aurait
le même code mais des annotations de typage différentes : on obtient alors
plusieurs copies de la fonction qui sont chacunes spécialisées pour un type en
particulier. C'est la solution retenue par le langage C++ et son système de
\emph{templates}.

En fait, si à l'exécution toutes les données sont représentées de manière
uniforme, il n'est pas nécessaire de créer plusieurs copies de la fonction.

\begin{figure}
  \insertcode{listappend.ml}
  \caption{Fonction de concaténation de listes en OCaml.}
  \label{fig:listappend}
\end{figure}

Par exemple, la fonction de la figure~\ref{fig:listappend} n'opère que sur la
structure du type liste (en utilisant ses constructeurs \texttt{{[}{]}} et
\listcons ainsi que le filtrage) : les éléments de \texttt{lx} et \texttt{ly} ne
sont pas manipulés à part pour les transférer dans le résultat.

% TODO [E] tu parles uniquement du polymorphisme paramétrique : attention il y a
% la classe de polymorphisme et son implantation qui peuvent diverger

Elle peut donc être typée avec n'importe quelle type 
$a~\textrm{list} → a~\textrm{list} → a~\textrm{list}$. Pour introduire cette
généricité, on modifie le système de types pour transformer :

\[ ∀ a. \textrm{append} :
   a~\textrm{list} → a~\textrm{list} → a~\textrm{list}
\]

en :

\[ \textrm{append} : ∀ a.
   a~\textrm{list} → a~\textrm{list} → a~\textrm{list}
\]

% TODO [E] que veux tu dire ?

% TODO [E] où la variable de types peut prendre n'importe quelle valeur de type

Au lieu d'associer à chaque expression un type, dans certains cas on lui associe
un schéma de types, instanciable en un type concret. En quelque sorte, on
déplace la vérification du site d'appel au site de définition. Cette technique a
été décrite en premier dans~\cite{Milner78}.
Un tour d'horizon de différents types statiques, avec en particulier du
polymorphisme, on pourra se référer à~\cite{TAPL}.

\paragraph{Expressivité :} les systèmes de types les plus simples expriment des
contrats esssentiellement liés à la sûreté d'exécution, pour ne pas utiliser des
valeurs de types incompatibles entre eux. Mais il est possible d'étendre le
langage avec des annotations plus riches : par exemple en vérifiant statiquement
que des listes ne sont pas vides\cite{lightweight-static-capabilities}, ou dans
le domaine de la sécurité, d'empêcher des fuites d'information~\cite{LZ06a}.

\paragraph{Qualificateurs de types :} Dans le cas particulier des vulnérabilités
liées à une mauvaise utilisation de la mémoire, les développeurs du noyau Linux
ont ajouté un système d'annotations au code source. Un pointeur peut être décoré
d'une annotation \texttt{\_\_kernel} ou \texttt{\_\_user} selon s'il est sûr ou
pas. Celle-ci sont ignorées par le compilateur, mais un outil d'analyse statique
ad-hoc nommé Sparse~\link{sparse}~\cite{TorvaldsSparse} % TODO citer un seul
peut être utilisé pour détecter les cas les plus simples d'erreurs.
Il demande aussi au programmeur d'ajouter beaucoup d'annotations dans le
programme. % TODO quantifier

Ce système d'annotations sur les types a été formalisé sous le nom de
\emph{qualificateurs de types} : chaque type peut être décoré d'un ensemble de
qualificateurs (à la manière de \texttt{const}), et des règles de typage
permettent d'établir des propriétés sur le programme. Ces analyses ont été
implantée dans l'outil CQual~\cite{toplas-quals}. Ce système peut servir à
inférer les annotations \texttt{const}~\cite{pldi99}, à l'analyse de souillure
pour les chaîne de format~\cite{usenix01} % TODO citer Newsham sur le format ?
et des propriétés dépendantes du flot de contrôle, comme des invariants sur les
verrous~\cite{pldi02}, à rapprocher du concept de \emph{typestates}
~\cite{tse12-typestate}. Il a également été appliqué à la classe de
vulnérabilités sur les pointeurs utilisateurs dont il est ici
l'objet~\cite{cquk-usenix04}. Puisqu'elle consiste à ajouter un qualificateur à
chaque étage de type, cette approche est plus générique mais plus complexe que
la nôtre. % TODO bof

% TODO [E] préciser pourquoi

\section{Langages d'implantation}

% TODO C

Le langage C~\cite{KandR} présente des difficultés pour être analysé
automatiquement.

\begin{itemize}
\item les nombres flottants
(types \texttt{float}, \texttt{double} et \texttt{long double}) ont une
sémantique particulière, et il n'est pas correct d'approcher leur sémantique
par une sémantique dans $ℝ$. Afin d'être correct, il faut établir des domaines
spécifiques au flottant, comme~\cite{floatpoly}. Un tour d'horizon des
difficultés liées aux flottants est effectué dans~\cite{floatpitfalls}.

\item les pointeurs sur fonction
rendent floue la limite qui est habituellement présente entre instructions et
données. En leur présence il est impossible de faire une analyse de flot de
contrôle indépendante du flot de données. Pour pouvoir les traiter, il faut que
le domaine abstrait en question soit assez précis pour qu'un pointeur abstrait
se concrétise en un ensemble réduit de fonctions. Dans le cas où le domaine ne
peut pas borner l'ensemble des fonctions possibles et renvoie $\top$, l'analyse
ne peut pas continuer.

\item l'allocation dynamique
de données, présente dans le langage C par le biais des fonctions
\texttt{malloc} et \texttt{free}, modifie le modèle mémoire
nécessaire. Sans celle-ci, l'ensemble des zones mémoire possibles peut être
décrit statiquement : ce sont les noms de variable. Ce qu'introduit
\texttt{malloc} au langage, c'est une zone mémoire qui n'a pas de nom, et sur
laquelle on n'a qu'un pointeur.

% TODO, bof + expand

\item le transtypage (\emph{casts}) entre entiers et pointeurs
est particulièrement délicat à traiter. Dans les modèles abstraits, les
pointeurs sur données ou sur fonctions n'ont pas de représentation numérique,
seulement une représentation symbolique. Même dans l'exécution concrète, la
représentation numérique d'un pointeur est lié à de nombreux choix faits par
l'environnement d'exécution (comme la randomisation de l'espace d'adressage) qui
ne peuvent pas facilement être modélisés.

\end{itemize}

% TODO Ada

% TODO ASM

% TODO Rust

\section{Langages d'analyse}

Le langage C \cite{KandR,AnsiC} a été conçu pour être une sorte d'assembleur
portable, permettant décrire du code indépendamment de l'architecture sur
laquelle il sera compilé. Historiquement, c'est il a permis de créer Unix, et
ainsi de nombreux logiciels bas niveau sont écrits en C. En particulier, il
existe des compilateurs de C vers les différents langages machine pour à peu
près toutes les architectures.

\begin{figure}
  \centering

  \begin{tikzpicture}
[>=latex]
  \node (me) {Langage intermédiaire};
  \node[ above left of=me, xshift=-4cm] (fe1) {Langage source 1};
  \node[ below left of=me, xshift=-4cm] (fe2) {Langage source 2};
  \node[above right of=me, xshift=4cm] (be1) {Langage destination 1};
  \node[below right of=me, xshift=4cm] (be2) {Langage destination 2};

  \draw[->] (fe1) -- (me);
  \draw[->] (fe2) -- (me);

  \draw[->] (me) -- (be1);
  \draw[->] (me) -- (be2);

  \path (me) edge [loop below] node {Optimisations, analyses} (me);

  \coordinate (legbase) at ($ (me) + (0,-2cm) $);

  \draw[bigbrace] ($ (me.south east |- legbase) + (-2mm, 0) $)
               to node [yshift=-5mm] {Middle-end}
                  ($ (me.south west |- legbase) + (2mm, 0) $);
  \draw[bigbrace] (me.south west |- legbase)
               to node [yshift=-5mm]{Front-end}
                  (fe2.south west |- legbase);
  \draw[bigbrace] (be2.south east |- legbase)
               to node [yshift=-5mm]{Back-end}
                  (me.south east |- legbase);
\end{tikzpicture}


  \caption{Décomposition d'un compilateur : front-ends, middle-end, back-ends}
  \label{fig:middle-end}
\end{figure}

Lors de l'écriture d'un compilateur, on a besoin d'un langage intermédiaire qui
fasse l'intermédiaire entre \emph{front-end} et \emph{back-end}
(figure~\ref{fig:middle-end}). Depuis ce langage on doit pouvoir exprimer des
transformations intermédiaires sur cette représentation (analyses sémantiques,
optimisations, etc), mais aussi compiler ce langage vers un langage machine.

L'idée de prendre C comme langage intermédiaire est très séduisante, mais
malheureusement sa sémantique est trop complexe et trop peu spécifiée. Il est
donc judicieux d'utiliser un langage plus simple à cet effet. Dans de nombreux
projets, des sous-ensembles de C ont été définis pour aller dans ce sens.

Les premiers candidats sont bien entendu les représentations intermédiaires
utilisées dans les compilateurs C. Elles ont l'avantage d'accepter en plus du C
standard, les diverses extensions (GNU, Microsoft, Plan9) utilisées par la
plupart des logiciels. En particulier, le noyau Linux repose fortement sur les
extensions GNU.

\paragraph{GCC} utilise une représentation interne nommée
GIMPLE\cite{gcc-gimple}. Il s'agit d'une structure d'arbre écrite en C, reposant
sur de nombreuses macros afin de cacher les détails d'implémentation pouvant
varier entre deux versions de GCC. Cette représentation étant réputée difficile
à manipuler, le projet MELT\cite{gcc-melt} permet de générer une passe de
compilation à partir d'un dialecte de Lisp.

\paragraph{LLVM}\cite{llvm-pres} est un compilateur développé par la communauté
puis sponsorisé Apple. À la différence de GCC, sa base de code est écrite en
C++. Il utilise une représentation intermédiaire qui peut être manipulée soit
sous forme d'une structure de données C++, soit d'un fichier de code-octet
compact, soit sous forme textuelle.

\paragraph{Objective Caml}\link{ocamlManual} utilise pour sa génération de code
une représentation interne nommée Cmm, disponible dans les sources du
compilateur sous le chemin \texttt{asmcomp/cmm.mli} (il s'agit donc d'une
structure de données OCaml). Ce langage a l'avantage d'être très restreint, mais
malheureusement il n'existe pas directement de traducteur permettant de compiler
C vers Cmm.

\paragraph{C-~-}\cite{spjcmm} \link{cmm}, dont le nom est inspiré du précédent,
est un projet qui visait à unifier les langages intermédiaires utilisés par les
compilateurs. L'idée est que si un front-end peut émettre du C-~- (sous forme de
texte), il est possible d'obtenir du code machine efficace. Le compilateur
Haskell GHC utilise une représentation intermédiaire très similaire à C-~-.

Comme le problème de construire une représentation intermédiaire adaptée à une
analyse statique n'est pas nouveau, plusieurs projets ont déjà essayé d'y
apporter une solution. Puisque qu'ils sont développés en parallèle des
compilateurs, le support des extensions est en général moins important dans ces
langages.

\paragraph{CIL}\cite{NeculaCil} \link{kerneisCil} est une représentation en
OCaml d'un programme C, développée depuis 2002. Grâce à un mécanisme de
greffons, elle permet de prototyper rapidement des analyses statiques de
programmes.

\paragraph{Compcert} est un projet qui vise à produire un compilateur certifié
pour C. C'est à dire que le fait que les transformations conservent la
sémantique est prouvé. Il utilise de nombreux langages intermédiaires, dont CIL.
Pour le front-end, le langage se nomme Clight\cite{cfront}. Les passes de
middle-end, quant à elles, utilisent Cminor\cite{cminorSL}.

\paragraph{Newspeak}\cite{newspeak} est un langage intermédiaire développé par
EADS Innovation Works, et qui est spécialisé dans l'analyse de valeurs par
interprétation abstraite. Il sera décrit plus en détails dans la
section~\ref{sec:npk}.

\section{Plan de la thèse}

Cette thèse comporte trois parties.

\paragraph{La partie~\ref{part:ctx}} présente le contexte de ces travaux. Le
fonctionnement général d'un système d'exploitation y est détaillé, et les
problèmes de manipulation de pointeurs contrôlés par l'utilisateur y sont
introduits. On fait ensuite un tour d'horizon des techniques existantes
permettant de traiter ce problème.

% TODO [E] + proposition

\paragraph{La partie~\ref{part:lang}} décrit notre solution : \langname, un
langage impératif. On y décrit sa syntaxe, sa sémantique ainsi qu'un système de
types statiques. On l'étend ensuite pour capturer les problèmes d'adressage
mémoire présents dans les systèmes d'exploitation en ajoutant des pointeurs
contrôlés par l'utilisateur. Le système de types est également étendu. Pour
chacune de ces variantes, on établit la propriété de sûreté de typage reliant la
sémantique dynamique aux types statiques.

% TODO [E] dire un mot de Newspeak dès l'intro ou l'état de l'art

\paragraph{La partie~\ref{part:xp}} documente la démarche expérimentale associée
à ces travaux. L'implantation du système de types est décrite, afin que la
manière de transformer automatiquement du code C en \langname. Un cas d'étude
est déroulé, consistant d'un bug ayant touché le noyau Linux. Il est démontré
que le système de type capture précisément ce genre d'erreur de programmation.
Enfin, les possibilités d'extension tant théoriques qu'expérimentales sont
présentées.

%  - omniprésence
%  - complexité
%  - vérification
%  - taux de bugs après relecture
%  - vérification automatique obligatoire
%  - historique des techniques
%  - analyse non intrusive
%  - méthodes de développement du noyau
%  - syntaxe, sémantique, typage

% TODO [E] dire au début que c'est tendance Unix si c'est bien le cas

% TODO [E] + besoins EADS

% vim: spelllang=fr



\long\def\partintro{
Après avoir décrit le contexte général de ces travaux, nous on décrivons les
enjeux.

Le chapitre~\ref{cha:os} explore plus en détail le fonctionnement d'un système
d'exploitation, y compris la séparation du code en plusieurs niveaux de
privilèges. L'architecture Intel 32 bits est prise comme support. En
particulier, le mécanisme des appels système est décrit et on montre qu'une
implantation naïve de la communication entre espaces utilisateur et noyau casse
toute isolation.

Le chapitre~\ref{cha:etatdelart} consiste en un tour d'horizon des techniques
existantes en analyses de programmes. Ces analyses se centrent autour des
problèmes liés à la vérification de code système ou embarqué, y compris le
problème de manipulation mémoire évoqué dans le chapitre~\ref{cha:os}.

On conclut en présentant notre technique: \langname, un langage permettant de
typer des programmes impératifs, plus précisément en ajoutant des types
pointeurs abstraits.
}

\part{Méthodes formelles pour la sécurité}
\label{part:ctx}


\chapter{Systèmes d'exploitation}

\label{cha:os}
Le système d'exploitation est le programme qui permet à un système informatique
d'exécuter d'autres programmes. Son rôle est donc capital et ses responsabilités
multiples. Dans ce chapitre, nous allons voir à quoi il sert, et comment il peut
être implanté. Pour ce faire, nous présentons ici de quoi est constitué un
sustème Intel 32 bits et ce dont on se sert pour y implanter un système
d'exploitation.

\section{Architecture physique}

Un système informatique est principalement constitué d'un processeur (ou CPU
pour \emph{Central Processing Unit}), de mémoire principale (ou RAM pour
\emph{Random Access Memory}), et de divers périphériques.

Le processeur est constitué de divers registres internes qui permettent
d'encoder l'état dans lequel il se trouve : quelle est l'instruction courante
(registre \eip), quelle est la hauteur de la pile système (registre \esp), etc.
Son fonctionnement peut être vu de la manière la plus simple qui soit comme le
cycle :

\begin{itemize}
\item
  charger depuis la mémoire la prochaine instruction
\item
  \emph{(optionnel)} charger depuis la mémoire les données référencées
  par l'instruction
\item
  effectuer l'instruction
\item
  \emph{(optionnel)} stocker en la mémoire les données modifiées
\end{itemize}

Les instructions sont constituées d'un mnémonique (indicant quelle opération
faire) et d'un ensemble d'opérandes. La signification des opérandes dépend du
mnémonique, mais en général elles permettent de désigner les sources et la
destination. Les opérandes peuvent avoir plusieurs formes : une valeur
immédiate (\$4), un nom de registre (\texttt{\%eax}) ou une référence à la
mémoire (directement : \texttt{addr} ou indirectement : \texttt{(\%ecx)}
\footnote{
  Cela consiste à interpréter le contenu du regitre \ecx comme une adresse
  mémoire.
}
). On décrit les mnémoniques les plus utilisés, permettant de compiler un cœur
de langage impératif:

\begin{itemize}

\item \texttt{mov src, dst} copie le contenu de \texttt{src} dans \texttt{dst}.

\item \texttt{add src, dst} calcule la somme des contenus de \texttt{src} et
  \texttt{dst} et place ce résultat dans \texttt{dst}.

\item \texttt{push src} place \texttt{src} sur la pile, c'est-à-dire que cette
  instruction enlève au pointeur de pile \esp la taille de \texttt{src}, puis
  place \texttt{src} à l'adresse mémoire de la nouvelle valeur \esp.

\item \texttt{pop src} réalise l'opération inverse : elle charge le contenu de
  la mémoire à l'adresse \esp dans \texttt{src} puis incrémente \esp de la
  taille correspondante.

\item \texttt{jmp addr} saute à l'adresse \texttt{addr} : c'est l'équivalent de
  \texttt{mov addr, \%eip}.

\item \texttt{call addr} sert aux appels de fonction : cela revient à
  \texttt{push \%eip} puis \texttt{jmp addr}.

\item \texttt{ret} sert à revenir d'une fonction : c'est l'équivalent de
  \texttt{pop \%eip}.

\end{itemize}

Certaines de ces instructions font référence à la pile par le biais du registre
\esp. Cette zone mémoire n'est pas gérée de manière particulière. Elle permet de
gérer la pile des appels de fonction en cours grâce à la manière dont
\texttt{jmp} et \texttt{ret} fonctionnent. Elle sert aussi à stocker les
variables locales des fonctions.

À l'aide de ces quelques instructions on peut implanter des algorithmes
impératifs. Mais pour faire quelque chose de visible, comme afficher à l'écran
ou envoyer un paquet sur le réseau, cela ne suffit pas : il faut parler au reste
du matériel.

Pour ceci, il y a deux techniques principales. D'une part, certains
périphériques sont dits \emph{memory-mapped} : ils sont associés à un espace
mémoire particulier, qui ne permet pas de stocker des informations mais de lire
ou d'écrire des données dans le périphérique. Par exemple, écrire à l'adresse
$0xB8000$ permet d'écrire des caractères à l'écran. L'autre système principal
est l'utilisation des ports d'entrée/sortie. Cela correspond à des instructions
spéciales \texttt{in \%ax, port} et \texttt{out port, \%ax} où \texttt{port}
est un numéro qui correspond à un périphérique particulier. Par exemple, en
écrivant dans le port $0x60$ on peut contrôler l'état des indicateurs lumineux
du clavier PS/2.

\section{Tâches et niveaux de privilèges}
\label{sec:taches}

\subsection*{Alternance des tâches}

Sans mécanisme particulier, le processeur exécuterait uniquement une suite
d'instruction à la fois. Pour lui permettre d'exécuter plusieurs tâches, un
système de partage du temps existe.

À des intervalles de temps réguliers, le système est programmé pour recevoir une
interruption. C'est une condition exceptionnelle (au même titre qu'une division
par zéro) qui fait sauter automatiquement le processeur dans une routine de
traitement d'interruption. À cet endroit le code peut sauvegarder les registres
et restaurer un autre ensemble de registres, ce qui permet d'exécuter plusieurs
tâches de manière entrelacée. Si l'alternance est assez rapide, cela peut donner
l'illusion que les programmes s'exécutent en parallèle. Comme l'interruption
peut survenir à tout moment, on parle de multitâche préemptif.

En plus de cet ordonnancement de processus, l'architecture Intel permet
d'affecter des niveaux de privilège à ces tâches, en restreignant le type
d'instructions exécutables, ou en donnant un accès limité à la mémoire aux
tâches de niveaux moins élevés.

Il y a 4 niveaux de privilèges (nommés aussi \emph{rings}): le \ring{0} est le
plus privilégié, le \ring{3} le moins privilégié. Dans l'exemple précédent, on
pourrait isoler l'ordonnanceur de processus en le faisant s'exécuter en \ring{0}
alors que les autres tâches seraient en \ring{3}.

\subsection*{Mémoire virtuelle}

À partir du moment où plusieurs processus s'exécutent de manière concurrente, un
problème d'isolation se pose: si un processus peut lire dans la mémoire d'un
autre, des informations peuvent fuiter; et s'il peut y écrire, il peut en
détourner l'exécution.

Le mécanisme de mémoire virtuelle permet de donner à deux tâches une vue
différente de la mémoire: c'est-à-dire que vue de tâches différentes, une
adresse contiendra une valeur différente (figure~\ref{fig:memoire-virtuelle}).

\begin{figure}[h] % fig:memoire-virtuelle {{{
\centering
\begin{tikzpicture}
  [memlayout/.style={draw
                    ,shape=rectangle
                    ,minimum height=5cm
                    ,minimum width=1cm
                    }
  ,node distance=2.5cm
  ]

  \node[memlayout]            (A) {};
  \node[memlayout,right of=A] (B) {};
  \node[memlayout,right of=B] (C) {};

  \def\putlabel#1#2{
    \path (#1.north) -- ++(0,3mm) node {#2};
  }

  \putlabel{A}{Processus 1}
  \putlabel{B}{Mémoire}
  \putlabel{C}{Processus 2}

  \def\mmap#1#2#3#4#5{
    \begin{scope}
      [mzone/.style={draw
                    ,shape=rectangle
                    ,fill=black!20
                    ,minimum height=#5 mm
                    ,minimum width=1cm
                    }
      ,corr/.style={dashed}
      ]
      \node[mzone] (L) at ($ (#1.north) - (0, #2) $) {};
      \node[mzone] (R) at ($ (#3.north) - (0, #4) $) {};

      \draw[corr] (L.north east) -- (R.north west);
      \draw[corr] (L.south east) -- (R.south west);
    \end{scope}
  }

  \mmap{A}{1.5}{B}{2.5}{4}
  \mmap{A}{4.5}{B}{4}{3}
  \mmap{B}{1}{C}{1}{2}
  \mmap{B}{3.2}{C}{2.7}{4}
\end{tikzpicture}
\caption{Mécanisme de mémoire virtuelle.}
\label{fig:memoire-virtuelle}
\end{figure} % }}}

\label{page:mem-virt}
Ce mécanisme est contrôlé par valeur du registre \crtrois: les 10 premiers bits
d'une adresse virtuelle sont un index dans le répertoire de pages qui commence à
l'adresse contenue dans \crtrois. À cet index, se trouve l'adresse d'une table
de pages. Les 10 bits suivants de l'adresse sont un index dans cette page,
donnant l'adresse d'une page de 4 kio (figure~\ref{fig:pagetables}).
Plus de détails sur l'utilisation de ce mécanisme seront donnés dans la
section~\ref{sec:linux-sys}.

\begin{figure}[h]
\begin{tikzpicture}
[ bit/.style={ draw
             , shape=rectangle
             , minimum width=4mm
             , minimum height=4mm
             , node distance=4mm
             }
, pstuff/.style={ draw
                , shape=rectangle
                , node distance=1cm
                , text width=2cm
                , text centered
                }
, >=latex
]

\node [bit] (b31) {};

% generate this ?
\node [bit, right of=b31] (b30) {}; \node [bit, right of=b30] (b29) {}; \node [bit, right of=b29] (b28) {};
\node [bit, right of=b28] (b27) {}; \node [bit, right of=b27] (b26) {}; \node [bit, right of=b26] (b25) {};
\node [bit, right of=b25] (b24) {}; \node [bit, right of=b24] (b23) {}; \node [bit, right of=b23] (b22) {};
\node [bit, right of=b22] (b21) {}; \node [bit, right of=b21] (b20) {}; \node [bit, right of=b20] (b19) {};
\node [bit, right of=b19] (b18) {}; \node [bit, right of=b18] (b17) {}; \node [bit, right of=b17] (b16) {};
\node [bit, right of=b16] (b15) {}; \node [bit, right of=b15] (b14) {}; \node [bit, right of=b14] (b13) {};
\node [bit, right of=b13] (b12) {}; \node [bit, right of=b12] (b11) {}; \node [bit, right of=b11] (b10) {};
\node [bit, right of=b10] (b9) {}; \node [bit, right of=b9] (b8) {}; \node [bit, right of=b8] (b7) {};
\node [bit, right of=b7] (b6) {}; \node [bit, right of=b6] (b5) {}; \node [bit, right of=b5] (b4) {};
\node [bit, right of=b4] (b3) {}; \node [bit, right of=b3] (b2) {}; \node [bit, right of=b2] (b1) {};
\node [bit, right of=b1] (b0) {};

\foreach \b in {0,11,12,21,22,31}
    \node [above of=b\b, node distance=5mm] {\small \b};

\newcommand{\mbrace}[3]{
    \draw [bigbrace]                      ($ (b#1.south east) + (-0.1,-0.3) $)
    -- node [midway, yshift=-2mm] (#3) {} ($ (b#2.south west) +  (0.1,-0.3) $);
} % mbrace

\mbrace{0}{11}{bpp};
\mbrace{12}{21}{bpt};
\mbrace{22}{31}{bpd};

\node[pstuff, below of=bpd] (pd) {Répertoire de pages};
\node[pstuff, below of=bpt] (pt) {Table de pages};
\node[pstuff, below of=bpp] (pp) {Page de 4 kio};

\node[left of=pd, node distance=2cm] (cr) {\crtrois};

\draw [->] (bpp) -- (pp);
\draw [->] (bpt) -- (pt);
\draw [->] (bpd) -- (pd);

\draw [->] (cr) -- (pd);
\draw [->] (pd) -- (pt);
\draw [->] (pt) -- (pp);
\draw [->] (pp) -- ++ (3,0);

\end{tikzpicture}

\caption{Implantation de la mémoire virtuelle}
\label{fig:pagetables}
\end{figure}

\section{Appels système}

Avec une telle isolation, tout le code qui est exécuté en \ring{3} a une
expressivité limitée. Il ne peut pas contenir d'instructions privilégiées comme
\asminstr{in} ou \asminstr{out}, ni faire référence à des périphériques mappés
en mémoire. C'est en effet au noyau d'accéder au matériel, et pas au code
utilisateur.

Il est donc nécessaire d'appeller une routine du noyau depuis le code
utilisateur. C'est le but des appels système. Cela consiste à coupler une
fonction du \ring{3} à une fonction du \ring{0}: en appellant la fonction non
privilégiée, le flot d'exécution se retrouve dans le noyau avec les bons
privilèges.

Bien sûr, il n'est pas possible de faire directement un \asminstr{call} puisque
cela consisterait à faire un saut vers une zone plus privilégiée. Il y a
plusieurs manières d'implanter ce mécanisme. Nous décrivons ici la technique
historique à l'aide d'interruptions.

Le processeur peut répondre à des interruptions, qui sont des évenements
extérieurs. Cela permet d'écrire du code asynchrone. Par exemple, une fois qu'un
long transfert mémoire est terminé, une interruption est reçue. D'autres
interruptions dites logicielles peuvent arriver lorsqu'une erreur se produit.
Par exemple, diviser par zéro provoque l'interruption 0, et tenter d'exécuter
une instruction privilégiée provoque l'interruption 14. On peut aussi provoquer
une interruption par une instruction \asminstr{int} dédiée.

Une table globale définit, pour chaque numéro d'interruption : quelle est la
routine à appeller pour la traiter, avec quel niveau de privilège, ainsi que le
niveau de privilège requis pour pouvoir déclencher celle-ci avec l'instruction
\asminstr{int}.

Il est donc possible de créer une interruption purement logicielle (on utilise
en général le numéro 128 soit \hexa{80}), déclenchable en \ring{3} et traitée en
\ring{0}. Les registres sont préservés, donc on peut les utiliser pour passer un
numéro d'appel système (par exemple 3 pour \texttt{read()} et 5 pour
\texttt{open()}) et leurs arguments.

\section{Le \emph{Confused Deputy Problem}}
\label{sec:secu-syscalls}

On a vu que les appels systèmes permettent aux programmes utilisateur d'accéder
au services du noyau. Ils forment donc une interface particulièrement sensible aux
problèmes de sécurité.

Comme pour toutes les interfaces, on peut être plus ou moins fin. D'un
côté, une interface pas assez fine serait trop restrictive et ne permettrait pas
d'implanter tout type de logiciel. De l'autre, une interface trop laxiste
(«écrire dans tel registre matériel») empêche toute isolation. Il faut
donc trouver la bonne granularité.

Nous allons présenter ici une difficulté liée à la manipulation de mémoire au
sein de certains types d'appels système.

Il y a deux grands types d'appels systèmes: d'une part, ceux qui renvoient un
simple nombre, comme \texttt{getpid} qui renvoie le numéro du processus
appelant.

\insertcode{appel-syscall-getpid.c}

Ici, pas de difficulté particulière: la communication entre le \ring{0} et le
\ring{3} est faite uniquement à travers les registres, comme décrit dans la
section~\ref{sec:impl-syscall}.

Mais la plupart des appels systèmes communiquent de l'information de manière
indirecte, à travers un pointeur. L'appellant alloue une zone mémoire dans son
espace d'adressage et passe un pointeur à l'appel système. Ce mécanisme est
utilisé par exemple par la fonction \texttt{gettimeofday}
(figure~\ref{fig:appel-gettimeofday}).

\begin{figure}[h]
\insertcode{appel-syscall-gettimeofday.c}
\caption{Appel de gettimeofday}
\label{fig:appel-gettimeofday}
\end{figure}

Considérons une implantation naïve de cet appel système qui écrirait
directement à l'adresse pointée. Si le pointeur fourni est dans l'espace
d'adressage du processus, on est dans le cas d'utilisation normal et l'écriture
est donc possible.

Si l'utilisateur passe un pointeur dont la valeur correspond à la mémoire
réservée au noyau, que se passe-t-il? Comme le déréférencement est fait dans le
code du noyau, il est également fait en \ring{0}, et va pouvoir être réalisé
sans erreur: l'écriture se fait et potentiellement une structure importante du
noyau est écrasée.

Un utilisateur malveillant peut donc utiliser cet appel système pour écrire à
n'importe quelle adresse dans l'espace d'adressage du noyau. Ce problème vient
du fait que l'appel système utilise les privilèges du noyau au lieu de celui qui
contrôle la valeur des paramètres sensibles. Cela s'appelle le \emph{Confused
Deputy Problem}\cite{hardy88confused}.

La bonne solution est de tester dynamiquement la valeur du pointeur: s'il pointe
en espace noyau, il faut indiquer une erreur plutôt que d'écrire. Sinon, il peut
toujours y avoir une erreur, mais au moins le noyau est protégé.

\begin{figure}
\insertcode{implem-gettime.c}
\caption{Implantation de l'appel système \texttt{gettimeofday}}
\label{fig:implem-gettime}
\end{figure}

Dans le noyau, un ensemble de fonctions permet d'effectuer des copies sûres. La
fonction \texttt{access\_ok} réalise le test décrit précédemment. Les fonctions
\texttt{copy\_from\_user} et \nolinkurl{copy_to_user} réalisent une copie de la
mémoire après avoir fait ce test. Ainsi, l'implantation correcte de l'appel
système \texttt{gettimeofday} fait appel à celle-ci
(figure~\ref{fig:implem-gettime}).

Pour préserver la sécurité du noyau, il est donc nécessaire de vérifier la
valeur de tous les pointeurs dont la valeur est contrôlée par l'utilisateur.
Cette conclusion est assez contraignante, puisqu'il existe de nombreux endroits
dans le noyau où des données proviennent de l'utilisateur. Il est donc
raisonnable de vouloir vérifier automatiquement et statiquement l'absence de
tels défauts.

%\clearpage

%\section{Architecture et assembleur Intel}

%L'implantation d'un système d'exploitation est très proche du matériel sur
%lequel il s'exécute. Pour étudier une implantation en particulier, voyons ce que
%permet le matériel lui-même.

%Dans cette section nous décrivons le fonctionnement d'un processeur utilisant
%une architecture Intel 32 bits. Les exemples de code seront écrits en syntaxe
%AT\&T, celle que comprend l'assembleur GNU.\@

%La référence pour la description de l'assembleur Intel est la documentation du
%constructeur~\cite{intelsys}; une bonne explication de l'agencement dans la
%pile peut aussi être trouvée dans~\cite{SmashingTheStack}.

%Pour faire des calculs, le processeur est composé de registres, qui sont des
%petites zones de mémoire interne, et peut accéder à la mémoire principale.

%La mémoire principale contient divers types des données:

%\begin{itemize}
%\item le code des programmes à exécuter
%\item les données à disposition des programmes
%\item la pile d'appels
%\end{itemize}

%La pile d'appels est une zone de mémoire qui est notamment utilisée pour tenir
%une trace des calculs en cours. Par exemple, c'est ici que seront stockées les
%données propres à chaque fonction appelée: paramètres, adresse de retour et
%variables locales. La pile est manipulée par un pointeur de pile (\emph{stack
%pointer}), qui est l'adresse du «haut de la pile». On peut la manipuler en
%empilant des données (les placer au niveau du pointeur de pile et déplacer celui
%si) ou dépilant des données (déplacer le pointeur de pile dans l'autre sens et
%retourner la valeur présente à cet endroit).

%L'état du processeur est défini par la valeur de ses registres, qui sont des
%petites zones de mémoire interne (quelques bits chacun). Par exemple, la valeur
%du pointeur de pile est stockée dans \esp{}. Le registre \ebp{}, couplé à \esp{}
%sert à adresser les variables locales et paramètres d'une fonction, comme ce
%sera expliqué dans la section~\ref{sec:convappel}.

%L'adresse de l'instruction courante est accessible dans le registre \eip.

%En plus de ces registres spéciaux, le processeur possède de nombreux registres
%génériques, qui peuvent être utilisés pour réaliser des calculs intermédiaires.
%Ils sont nommés \eax{}, \ebx{}, \ecx{}, \edx{}, \esi{} et \edi{}. Ils peuvent
%être utilisés pour n'importe quel type d'opération, mais certains sont
%spécialisés: par exemple il est plus efficace d'utiliser \eax{} en accumulateur,
%ou \ecx{} en compteur.

%Les calculs sont décrits sous forme d'une suite d'instructions. Chaque
%instruction est composée d'un mnémonique et d'une liste d'opérandes. Les
%mnémoniques (\asminstr{mov}, \asminstr{call}, \asminstr{sub}, etc) définissent
%un type d'opération à appliquer sur les opérandes. L'instruction peut aussi être
%précédée d'une étiquette, qui correspondra à son adresse.

%\begin{center}
  %\shorthandoff{!}
  %\begin{tikzpicture}
    %[bbr/.style={bigbrace, decoration={mirror, amplitude=2mm}}
    %,asmbox/.style={minimum height=1cm}
    %,dway/.style={dashed}
    %]

    %\node[asmbox]              (lbl) {\asminstr{lbl:}};
    %\node[asmbox,right of=lbl] (mov) {\asminstr{mov}};
    %\node[asmbox,right of=mov] (op1) {\asminstr{\$4,}};
    %\node[asmbox,right of=op1] (op2) {\asminstr{\%ebx}};

    %\draw[bbr] (lbl.south west) -- (lbl.south east);
    %\draw[bbr] (mov.south west) -- (mov.south east);
    %\draw[bbr] (op1.south west) -- (op2.south east);

    %\coordinate (ops) at ($ (op1.south west)!.5!(op2.south east) $);

    %\node[below left of=lbl, xshift=-1cm, yshift=-5mm]   (tlbl) {Étiquette};
    %\node[below of=tlbl, node distance=7mm] (tmov) {Mnémonique};
    %\node[below of=tmov, node distance=7mm] (tops) {Opérandes};

    %\draw[dway] (tlbl) -| ($ (lbl.south) + (0, -0.3) $);
    %\draw[dway] (tmov) -| ($ (mov.south) + (0, -0.3) $);
    %\draw[dway] (tops) -| ($ (ops.south) + (0, -0.3) $);

  %\end{tikzpicture}
  %\shorthandon{!}
%\end{center}

%Ces opérandes peuvent être de plusieurs types:

%\begin{itemize}
%\item
%\item
  %une opérande mémoire, c'est-à-dire le contenu de la mémoire à une
  %adresse effective. Cette adresse effective peut être exprimée de
  %plusieurs manières:
  %\begin{itemize}
  %\item
    %directement: \texttt{addr}
  %\item
    %indirectement: \texttt{(\%ecx)}. L'adresse effective est le contenu
    %du registre.
  %\item
    %«base + déplacement»: \texttt{4(\%ecx)}. L'adresse effective est
    %le contenu du registre plus le déplacement (4 ici).
  %\end{itemize}
%\end{itemize}

%En pratique il y a des modes d'adressage plus complexes, et toutes les
%combinaisons ne sont pas possibles, mais ceux-ci suffiront à décrire les
%exemples suivants:

%\section{Fonctions et conventions d'appel}
%\label{sec:convappel}

%\begin{figure} % {{{ fig:stackframe
%\begin{tikzpicture}
[ stack/.style={draw,shape=rectangle,minimum height=1cm,minimum width=5cm}
, bigbrace/.style={decorate,decoration={brace,amplitude=3mm}}
, >=latex
]

  \path (0,0) node [stack] (gloc) {Locales de g}
  -- ++(0,-1) node [stack] (gebp) {}
  -- ++(0,-1) node [stack] (gret) {}
  -- ++(0,-1) node [stack] (gpar) {Arguments de g}
  -- ++(0,-1) node [stack] (floc) {Locales de f}
  -- ++(0,-1) node [stack] (febp) {}
  -- ++(0,-1) node [stack] (fret) {}
  -- ++(0,-1) node [stack] (fpar) {Arguments de f};

  \draw [->] (gebp) -- ++(4, 0) |- (febp.5);
  \draw [->] (febp.355) -- ++(1.5, 0) -- ++(0,-4);

  \draw [dashed] (fpar.south west) -- ++ (0,-2);
  \draw [dashed] (fpar.south east) -- ++ (0,-2);
  \draw [dashed] (gloc.north east) -- ++ (0,2);
  \draw [dashed] (gloc.north west) -- ++ (0,2);

  \node at (5, 1) (ox) {};

  \node [text width=3cm] at (-0.8, 2)    {Adresses hautes (sens de \asminstr{push})};
  \node [text width=3cm] at (-0.8, -8.5) {Adresses basses (sens de \asminstr{pop})};

  \node [text width=3cm] at (-6, -3) (asmtext) {\tt
    g:  ... \\
    leave \\
    ret \\
    f:  ... \\
    pushl \$2 \\
    pushl \$4 \\
    call g \\
    addl \$8, \%esp \\
    leave \\
    ret \\
    ... \\
    call f\\
    ...
    };

  \draw [-*] (gret) -- ++ (-3, 0) |- ($ (asmtext.north west) + (2, -3.5) $);
  \draw [-*] (fret) -- ++ (-3, 0) |- ($ (asmtext.north west) + (2, -6) $);

  \draw [<-] (gebp.west) -- ++ (-0.5,0) node [xshift=-4mm] {\ebp};
  \draw [<-] (gloc.west) -- ++ (-0.5,0) node [xshift=-4mm] {\esp};

  \draw [<-] ($ (asmtext.north west) + (2, 0) $) -- ++ (0.5,0) node [xshift=3mm] {\eip};

  \draw [bigbrace]
        ($ (gloc.north east) + (3, 0) $)
     -- node [auto,midway,xshift=3mm] {Cadre de g}
        ($ (gpar.south east) + (3, 0.1) $);

  \draw [bigbrace]
        ($ (floc.north east) + (3, -0.1) $)
     -- node [auto,midway,xshift=3mm] {Cadre de f}
        ($ (fpar.south east) + (3, 0) $);

\end{tikzpicture}

%\caption[Cadres de pile]{ Cadres de pile. }
%\label{fig:stackframe}
%\end{figure} % }}}

%Dans le langage d'assemblage, il n'y a pas de notion de fonction; mais
%\asminstr{call} et \asminstr{ret} permettent de sauvegarder et de restaurer une
%adresse de retour, ce qui permet de faire un saut et revenir à l'adresse
%initiale. Ce système permet déjà de créer des procédures, c'est-à-dire des
%fonctions sans arguments ni valeur de retour.

%Pour gérer ceux-ci, il faut que les deux morceaux (appelant et appelé) se
%mettent d'accord sur une convention d'appel commune. La convention utilisée sous
%GNU/Linux est appelée \emph{cdecl} et possède les caractéristiques suivantes:

%\begin{itemize}
%\item la valeur de retour d'une fonction est stockée dans \eax{}
%\item \eax{}, \ecx{} et \edx{} peuvent être écrasés sans avoir à les sauvegarder
%\item les arguments sont placés sur la pile (et enlevés) par l'appelant. Les
  %paramètres sont empilés de droite à gauche.
%\end{itemize}

%Pour accéder à ses paramètres, une fonction peut donc utiliser un adressage
%relatif à \esp. Cela peut fonctionner, mais cela complique les choses si elle
%contient aussi des variables locales. En effet, les variables locales sont
%placées sur la pile, au dessus des (c'est-à-dire, empilées après) paramètres,
%augmentant le décalage.

%Pour simplifier, la pile est organisée en cadres logiques: chaque cadre
%correspond à un niveau dans la pile d'appels de fonctions. Si $f$ appelle $g$,
%qui appelle $h$, il y aura dans l'ordre sur la pile le cadre de $f$, celui de
%$g$ puis celui de $h$.

%Ces cadres sont chainés à l'aide du registre \ebp{}: à tout moment, \ebp{}
%contient l'adresse du cadre de l'appelant.

%Prenons exemple sur la figure~\ref{fig:stackframe}: pour appeler
%\verb!g(4,2)!, \texttt{f} empile les arguments de droite à gauche.
%L'instruction \asminstr{call g} empile l'adresse de l'instruction suivante sur
%la pile puis saute au début de \texttt{g}.

%Au début de la fonction, les trois instructions permettent de sauvegarder
%l'ancienne valeur de \ebp{}, faire pointer \ebp{} à une endroit fixe dans le
%cadre de pile, puis allouer 8 octets de mémoire pour les variables locales.

%Dans le corps de la fonction \texttt{g}, on peut donc se référer aux variables
%locales par \texttt{-4(\%ebp)}, \texttt{-8(\%ebp)}, etc, et aux arguments par
  %\texttt{8(\%ebp)}, \texttt{12(\%ebp)}, etc.

%À la fin de la fonction, l'instruction \asminstr{leave} est équivalente à
%\texttt{mov \%ebp, \%esp} puis \texttt{pop \%ebp} et permet de défaire le cadre
%de pile, laissant l'adresse de retour en haut de pile. Le \asminstr{ret} final
%la dépile et y saute.


% vim: spelllang=fr


%\chapter{Typage}

%\input{chap/03-typage.tex}

\chapter{État de l'art}

\label{cha:etatdelart}
Dans ce chapitre, nous présentons un tour d'horizon des techniques existantes
permettant d'analyser des programmes. Un accent est mis sur la propriété de
sûreté décrite dans le chapitre~\ref{cha:os}, mais on ne se limite pas à
celle-ci.

L'analyse statique de programmes est un sujet de recherche actif depuis
l'apparition de la science informatique.

\section{Taxonomie}

\paragraph{Techniques statiques et dynamiques :} l'analyse peut être faite au
moment de la compilation, ou au moment de l'exécution. En général on peut
obtenir des informations plus précises de manière dynamique, mais cela ne prend
en compte que les parties du programme qui seront vraiment exécutées. Un autre
problème des techniques dynamiques est qu'il est souvent nécessaire
d'instrumenter l'environnement d'exécution (ce qui --- dans le cas où cela est
possible --- peut se traduire par un impact en performances). L'approche
statique, en revanche, nécessite de construire à l'arrêt une carte mentale du
programme, ce qui n'est pas toujours possible dans certains langages.

Dans la suite, nous considérerons essentiellement des techniques statiques,
précisant le contraire lorsque c'est nécessaire.

\paragraph{Cohérence et complétude :} le but d'une analyse statique est de
catégoriser les programmes selon leurs caractéristiques à l'exécution. Or,

\begin{theorem}[de Rice]
  Toute propriété non triviale sur le comportement dynamique des programmes est
  indécidable.\cite{rice}
\end{theorem}
% TODO formaliser un peu plus?

Autrement dit, il n'est pas possible d'écrire un analyseur statique parfait,
c'est à dire ne se trompant jamais. Toute technique statique va donc de se
retrouver dans au moins un des cas suivants :

\begin{itemize}
\item
  un programme valide est rejeté : on parle de \emph{faux positif}.
\item
  un programme invalide n'est pas détecté : on parle de
  \emph{faux négatif}.
\end{itemize}

En général on préfère s'assurer que les programmes acceptés possèdent la
propriété recherchée, quitte à en rejeter certains.

\section{Méthodes syntaxiques}

L'analyse la plus simple consiste à traiter un programme comme du texte, et à y
rechercher des motifs dangereux. Ainsi, utiliser des outils comme \texttt{grep}
permet parfois de trouver un grand nombre de vulnérabilités~\cite{SpenderGrep}.

On peut continuer cette approche en recherchant des motifs mais en étant
sensible à la syntaxe et au flot de contrôle du programme. Cette notion de
\emph{semantic grep} est présente dans l'outil Coccinelle
\cite{coccinelle09,coccinelle11} : on peut définir des
\emph{patches sémantiques} pour détecter ou modifier des constructions
particulières.

%TODO lire coccinelle09

\section{Interprétation abstraite}

L'interprétation abstraite est une technique d'analyse générique qui permet de
simuler statiquement tous les comportements d'un programme Cousot
\cite{Cousot77,Cousot92-1}. Un exemple d'application est de calculer les bornes
de variations des variables pour s'assurer qu'aucun débordement de tableau n'est
possible~\cite{AllamigeonHymansSSTIC07}. Cette technique est très puissante mais
possède plusieurs inconvénients. D'une part, pour réaliser une analyse
interprocédurale il faut partir d'un point en particulier du programme (comme la
fonction \texttt{main}). Cette hypothèse n'est pas facilement satisfaite dans un
noyau de système d'exploitation, qui possède de nombreux points d'entrée.

% TODO virer non?

\begin{figure}%{{{
\centering
\subbottom[
  L'ensemble des états d'erreur est calculable (zone hachurée)…
]{
\label{fig:ia-f1}
% 1
\begin{tikzpicture}[scale=0.7]%{{{
\clip  (-1,-1) rectangle (6,5);
%
% Hachures
\foreach \y in {-10,-9.3,...,11} \draw [gray, very thin, rotate=30] (-4,\y) to (18,\y);
%
% Ensemble sûr
\path [fill=white]
                 (0, 0) -- ++(0, 2)
            -- ++(1, 2) -- ++(3,-1)
            -- ++(1,-3) --    cycle;
%
% Ensemble sûr (redraw)
\draw            (0, 0) -- ++(0, 2)
            -- ++(1, 2) -- ++(3,-1)
            -- ++(1,-3) --    cycle;
%
\end{tikzpicture}%}}}
}
\hspace{5mm}
\subbottom[
  …mais pas celui des états atteints.
]{
\label{fig:ia-f2}
% 2
\begin{tikzpicture}[scale=0.7]%{{{
\clip  (-1,-1) rectangle (6,5);
%
% Hachures
\foreach \y in {-10,-9.3,...,11} \draw [gray, very thin, rotate=30] (-4,\y) to (18,\y);
%
% Ensemble sûr
\path [fill=white]
                 (0, 0) -- ++(0, 2)
            -- ++(1, 2) -- ++(3,-1)
            -- ++(1,-3) --    cycle;
%
% Ensemble sûr (redraw)
\draw            (0, 0) -- ++(0, 2)
            -- ++(1, 2) -- ++(3,-1)
            -- ++(1,-3) --    cycle;
%
% Comportement réel (sans erreurs)
\draw [fill=green!50,scale=0.8] plot[smooth cycle]
      coordinates{(1,1) (1,2) (2,3) (3,3) (3,1)};
%
\end{tikzpicture}%}}}
}

\subbottom[
  On construit donc une surapproximation calculable de la sémantique. Si
  celle-ci n'a pas d'erreur, on est assuré que le programme n'en aura pas non
  plus.
]{
\label{fig:ia-f3}
% 4
\begin{tikzpicture}[scale=0.7]%{{{
\clip  (-1,-1) rectangle (6,5);
%
% Hachures
\foreach \y in {-10,-9.3,...,11} \draw [gray, very thin, rotate=30] (-4,\y) to (18,\y);
%
% Ensemble sûr
\path [fill=white]
                 (0, 0) -- ++(0, 2)
            -- ++(1, 2) -- ++(3,-1)
            -- ++(1,-3) --    cycle;
%
% Ensemble sûr (redraw)
\draw            (0, 0) -- ++(0, 2)
            -- ++(1, 2) -- ++(3,-1)
            -- ++(1,-3) --    cycle;
%
% Approximation (précise)
\draw [fill=blue!50,scale around={1.3:(3,3)}]
           (1,1) -- (1,2)
        -- (2,3) -- (3,3)
        -- (3,1) -- cycle;
%
% Comportement réel (sans erreurs)
\draw [fill=green!50,scale=0.8] plot[smooth cycle]
      coordinates{(1,1) (1,2) (2,3) (3,3) (3,1)};
%
\end{tikzpicture}%}}}
}
\hspace{5mm}
\subbottom[
  Si on détecte une intersection, il n'est pas possible de savoir si elle est
  due à une erreur atteignable ou à une approximation trop laxe.
]{
\label{fig:ia-f4}
% 6
\begin{tikzpicture}[scale=0.7]%{{{
\clip  (-1,-1) rectangle (6,5);
%
% Hachures
\foreach \y in {-10,-9.3,...,11} \draw [gray, very thin, rotate=30] (-4,\y) to (18,\y);
%
% Ensemble sûr
\path [fill=white]
                 (0, 0) -- ++(0, 2)
            -- ++(1, 2) -- ++(3,-1)
            -- ++(1,-3) --    cycle;
%
% Ensemble sûr (redraw)
\draw            (0, 0) -- ++(0, 2)
            -- ++(1, 2) -- ++(3,-1)
            -- ++(1,-3) --    cycle;
%
% Approximation (imprécise)
\draw [fill=blue!50,scale around={1.3:(3,3)}]
            (1,1) -- (1,2)
         -- (3,4)
         -- (3,1) -- cycle;
%
% Comportement réel (sans erreurs)
\draw [fill=green!50,scale=0.8] plot[smooth cycle]
      coordinates{(1,1) (1,2) (2,3) (3,3) (3,1)};
%
\end{tikzpicture}%}}}
}

\caption{
  Surapproximation en interprétation abstraite.
  Il n'est pas possible de déterminer si l'ensemble des états atteignables
  est inclus dans l'ensemble des états sûrs (figure~\ref{fig:ia-f2}). En
  revanche, en construisant une surapproximation on peut parfois conclure
  (figures~\ref{fig:ia-f3} et~\ref{fig:ia-f4}).
}
\label{fig:ia}

\end{figure}%}}}

Les domaines les plus simples ne capturent aucune relation entre variables. Ce
sont des domaines non relationnels. Par exemple le domaine des signes capture
uniquement le signe des variables (figure~\ref{fig:dom-sig}), et le domaine des
intervalles retient les bornes de variations extrémales des variables
(figure~\ref{fig:dom-intvl}).

\begin{figure}%{{{
\centering

\begin{minipage}{0.4\textwidth}
  \begin{tikzpicture}
  \node at (1, 2) (t) {$\top$};
  \node at (0, 1) (m) {$-$};
  \node at (1, 1) (z) {$0$};
  \node at (2, 1) (p) {$+$};
  \node at (1, 0) (b) {$\bot$};
  \draw (t) -- (p) -- (b);
  \draw (t) -- (m) -- (b);
  \draw (t) -- (z) -- (b);
  \end{tikzpicture}
\end{minipage}
\begin{minipage}{0.4\textwidth}
  \begin{align*}
  γ~(-) = & ℝ^- \\
  γ~(0) = & \{0\} \\
  γ~(+) = & ℝ^+ \\
  \end{align*}
\end{minipage}

\caption{Domaine des signes}
\label{fig:dom-sig}
\end{figure}%}}}

\begin{figure} % {{{
\centering
\begin{tikzpicture}
\node at (0,0)   (b)  {$\bot$};

\node at (-4,2) (cm2)   {$\{-2\}$};
\node at (-2,2) (cm1)   {$\{-1\}$};
\node at  (0,2)  (c0)   {$\{0\}$};
\node at  (2,2)  (c1)   {$\{1\}$};
\node at  (4,2)  (c2)   {$\{2\}$};

\node at (-3,3) (rm2m1) {[-2;-1]};
\node at (-1,3) (rm10)  {[-1;0]};
\node at  (1,3) (r01)   {[0;1]};
\node at  (3,3) (r12)   {[1;2]};

\node at (-2,4) (rm20)  {[-2;0]};
\node at  (0,4) (rm11)  {[-1;1]};
\node at  (2,4) (r02)   {[0;2]};

\node at (-1,5) (rm21)  {[-2;1]};
\node at  (1,5) (rm12)  {[-1;2]};

\node at  (0,6) (rm22)  {[-2;2]};

\node at  (0,9) (t)  {$\top$};

\draw (b) -- (cm2) -- (rm2m1) -- (rm20) -- (rm21) -- (rm22);
\draw (b) -- (cm1) -- (rm10) -- (rm11) -- (rm12);
\draw (b) -- (c0) -- (r01) -- (r02);
\draw (b) -- (c1) -- (r12);
\draw (b) -- (c2);

\draw (c2) -- (r12) -- (r02) -- (rm12) -- (rm22);
\draw (c1) -- (r01) -- (rm11) -- (rm21);
\draw (c0) -- (rm10) -- (rm20);
\draw (cm1) -- (rm2m1);

\draw[loosely dashed] (cm2)   to [bend left=30] (t);
\draw[loosely dashed] (rm2m1) to [bend left=30] (t);
\draw[loosely dashed] (rm20)  to [bend left=30] (t);
\draw[loosely dashed] (rm21)  to [bend left=30] (t);

\draw[loosely dashed] (rm22)  to (t);

\draw[loosely dashed] (rm12)   to [bend right=30] (t);
\draw[loosely dashed] (r02)   to [bend right=30] (t);
\draw[loosely dashed] (r12)   to [bend right=30] (t);
\draw[loosely dashed] (c2)    to [bend right=30] (t);

\draw[loosely dashed] (cm2) -- +(-1,0);
\draw[loosely dashed] (c2)  -- +(1,0);

\end{tikzpicture}

\caption{Domaine des intervalles}
\label{fig:dom-intvl}
\end{figure} % }}}

Lorsque plusieurs variables sont analysées en même temps, utiliser de tels
domaines consiste à considérer un produit cartésien d'ensembles abstraits
(figure~\ref{fig:dom-cartesien})

Cela revient à oublier les relations entre les variables. Des domaines abstraits
plus précis permettent de retenir celles-ci. Pour ce faire, il faut modéliser
l'ensemble des valeurs des variables comme un tout. Parmi les domaines
relationnels courants on peut citer :

\begin{itemize}

\item Le domaine des polyèdres, historiquement l'un des premiers
domaines relationnels. Il permet de retenir tous les invariants affines entre
fonctions (figure~\ref{fig:dom-poly}).

\item Le domaine des zones permet de représenter des relations affines de
forme $v_i - v_j \le c$ (figure~\ref{fig:dom-zones}).

\item Le domaine des octogones est un compromis entre les polyèdres et les
zones. Il permet de représenter les relations $\pm v_i \pm v_j \le c$
(figure~\ref{fig:dom-octo}).

\end{itemize}

\begin{figure}%{{{

  \centering

  %{{{
  \subbottom[Domaine non relationnel]{
    \label{fig:dom-cartesien}
      \begin{tikzpicture}
      \draw[->] (0,-0.5) -- (0,3);
      \draw[->] (-0.5,0) -- (4,0);
%
      \draw[fill=red!30] (1,1) rectangle (3,2);
%
      \draw[dashed] (3,1) -- (3,0);
      \draw[dashed] (1,1) -- (1,0);
%
      \draw[dashed] (1,2) -- (0,2);
      \draw[dashed] (1,1) -- (0,1);
%
      \draw[line width=3pt] (0,1) -- (0,2);
      \draw[line width=3pt] (3,0) -- (1,0);
%
      \end{tikzpicture}
  }%}}}
  \subbottom[Domaine des polyèdres]{%{{{
    \label{fig:dom-poly}
%
    \begin{tikzpicture}[scale=0.5]
    \path[use as bounding box] (0,0) rectangle (8, 6);
%
    \draw[fill=red!30] (1,2) -- (2,4) -- (4,5) -- (6,5) -- (7,3) -- (5,1) -- cycle;
%
    \end{tikzpicture}
  }%}}}

  \subbottom[Domaine des zones]{%{{{
    \label{fig:dom-zones}
%
    \begin{tikzpicture}
    \path[use as bounding box] (0,0) rectangle (4,4);
%
    \draw[fill=red!30] (1,2) -- (2,3) -- (3,3) -- (3,2) -- (2,1) -- (1,1) -- cycle;
%
    \end{tikzpicture}
  }%}}}
  \subbottom[Domaine des octaèdres]{%{{{
    \label{fig:dom-octo}
%
    \begin{tikzpicture}
    \path[use as bounding box] (0,0) rectangle (4,4);
%
    \draw[fill=red!30] (1,1.5) -- (1,2) -- (2,3) -- (2.5,3) -- (3,2.5) -- (3,2) -- (2,1) -- (1.5,1) -- cycle;
%
    \end{tikzpicture}
  }%}}}

  \caption{Quelques domaines abstraits}
  \label{fig:dom-abstraits}
\end{figure}%}}}

En plus des domaines numériques, il est nécessaire d'employer des domaines
spécialisés dans la modélisation de la mémoire. Cela est nécessaire pour pouvoir
"suivre" les pointeurs. Par exemple, on peut représenter un pointeur par un
ensemble de variables possiblement pointées et une valeur abstraite représentant
le décalage (\emph{offset}) du pointeur par rapport au début de la zone mémoire.
Cette valeur peut elle-même être abstraite par un domaine numérique.

Au delà des domaines eux-mêmes, l'analyse se fait sous forme d'un calcul de
point fixe. La manière la plus simple est d'utiliser un algorithme de
\emph{liste de travail}, décrit par exemple dans~\cite{tapsoft95}. Les
raffinement en revanche sont nombreux.

Dès~\cite{Cousot77} il est remarqué que la terminaison de l'analyse n'est
assurée que si le treillis des valeurs abstraites est de hauteur finie, ou qu'un
opérateur d'élargissement (\emph{widening}) $\nabla$ est employé. L'idée est
qu'une fois qu'on a calculé quelques termes d'une suite croissante, on peut
réaliser une projection de celle ci. Par exemple, dans le domaine des
intervalles, $[0;2]~\nabla~[0;3] = [0;+\infty[$. On atteint alors un point fixe
mais qui est plus grand que celui qu'on aurait obtenu sans cette accélération :
on perd en précision. Pour en gagner, on peut redescendre sur le treillis des
points fixe avec une suite d'itérations décroissantes \cite{granger}. Dans
l'itération de point fixe, il est possible d'obtenir les résultats de manière
plus efficace en choisissant un ordre particulier dans les calculs des
sous-itérations\cite{policy}.

% TODO pas clair

En termes d'ingéniérie logicielle, implanter un analyseur statique est un défi
en soi. En plus des domaines abstraits, d'un itérateur, il faut traduire le code
source à analyser dans un langage, et traduire les résultats de l'analyse en un
ensemble d'"alarmes" à présenter à l'utilisateur.

Pour des retours d'expérience, on peut se référer aux descriptions d'
Astrée~\cite{Astree04,Astree05,AstreeScale},
CGS~\cite{cgs},
ou Coverity~\cite{coverityBillion},

% TODO pas top

Une interprétation abstraite est par construction sûre et incomplète, donc ce
qui sépare une bon analyseur d'un mauvais est sa précision. Dans le cas du
langage C, de nombreuses constructions rendent imprécises les analyses :

% TODO c'est déplacé en c1

L'interprétation abstraite a été utilisée pour analyser le flot de données entre
objets~\cite{liang2012taint}.

Elle a l'inconvénient de remonter de nombreuses fausses alarmes, puisqu'elle
nécessite d'avoir une vue précise du programme. Elle est aussi plus adaptée aux
programmes qui n'ont qu'un point d'entrée.

\section{Typage}

La plupart des langages de programmation incorporent la notion de type, qui
permet de détecter ou d'empêcher de manipuler des données incompatibles entre
elles.

Nous avons vu dans le chapitre~\ref{cha:os} qu'au niveau du langage machine, les
seules données qu'un ordinateur manipule sont des nombres. Selon les opérations
effectuées, ils seront interprétés comme des entiers, des adresses mémoires, ou
des caractères. Pourtant il est clair que certaines opérations n'ont pas de
sens: par exemple, ajouter deux adresses, ou déréférencer le résultat d'une
division sont des comportements qu'on voudrait pouvoir empêcher.

En un mot, le but du typage est de classifier les objets et de restreindre les
opérations possibles selon la classe d'un objet : ``ne pas ajouter des pommes et
des oranges''. Le modèle qui permet cette classification est appelé
\emph{système de types} et est en général constitué d'un ensemble de
\emph{règles de typage}, comme ``un entier plus un entier égale un entier''.

Il y a deux grandes familles de systèmes de types, selon quand se fait la
vérification de types. On peut en effet l'effectuer au moment de l'exécution, ou
au contraire prévenir les erreurs à l'exécution en la faisant au moment de la
compilation (ou avant l'interprétation).

\paragraph{Typage dynamique :} dans ce cas, chaque valeur manipulée par le
programme est décorée d'une étiquette définissant comment interpréter la valeur
en question. Les règles de typage sont alors réalisées à l'exécution. Par
exemple, l'opérateur "$+$" vérifie que ces deux opérandes ont une étiquette
"entier", et construit alors une valeur obtenue en faisant l'addition des deux
valeurs, avec une étiquette "entier". Symboliquement :

\begin{verbatim}
fonction addition(x, y) {
  if ( (etiquette_type(x) != type_entier)
    || (etiquette_type(y) != type_entier)
    ) {
      erreur ("L'addition est seulement définie pour deux entiers")
  } else {
    vx = valeur(x)
    vy = valeur(y)
    v = addition_machine(vx, vy)
    return Valeur(valeur=v, type=type_entier)
  }
}
\end{verbatim}

Le langage Python~\link{python} utilise cette stratégie, qui est illustrée dans
la session interactive de la figure~\ref{fig:typage-dynamique}.

\begin{figure}
  \insertcode{typage-dynamique.pycon}
  \caption{Session Python présentant le typage dynamique}
  \label{fig:typage-dynamique}
\end{figure}

L'analyse de taintage consiste à ajouter des étiquettes au données décrivant
leur provenance, de manière à ce que chaque accès dangereux soit vérifié à
l'exécution. Cela se marie bien avec les langages typés dynamiquement, comme le
montre l'exemple célèbre du mode ``souillé'' (\emph{tainted}) de
Perl~\cite{perlCamelBook}. Il est aussi possible, mais plus délicat, d'appliquer
cette technique à des langages compilés~\cite{clause-etal-issta07,oakland10}.

\paragraph{Typage statique :} dans ce cas on fait les vérifications à la
compilation. En quelque sorte, l'approche dynamique est pessimiste, puisqu'elle
demande de traiter très souvent le cas où les types ne sont pas corrects.
Intuitivement, dans le cas où toutes les fonctions se comportent bien, faire la
vérification est inutile. Pour vérifier ceci, on donne à chaque fonction un
contrat comme "si deux entiers sont passés, le fonction renverra un entier sans
provoquer d'erreur". Cet ensemble de contrats peut être vérifié statiquement par
le compilateur, à l'aide d'un système de types statique.

Ainsi la fonction "$+$" est typée \texttt{(int * int)} $→$ \texttt{int}. Les
règles permettant de vérifier le typage sont par exemple les suivantes :

\begin{itemize}
\item
  une constante entière est toujours de type \texttt{int}.
\item
  si $f$ a pour type $(t_1, …, t_n) → t$, que $e_1$ a pour type $t_1$,
  \ldots{}, que $e_n$ a pour type $t_n$, alors $f(e_1, …, e_n)$ a pour
  type $t_n$
\item
  si en considérant que $e_1$ a pour type $t_n$, …, et que $e_n$ a pour type
  $t_n$, on arrive à typer le corps de $f$ et que sa valeur de retour a alors
  pour type $t$,
  alors $f$ a pour type $(t_1, …, t_n) → t$
\end{itemize}

Cet ensemble de règles, une fois formalisé et implanté, est en général assez
efficace pour vérifier la correction d'un programme.

\paragraph{Typage fort ou faible :}

Si un système de types statique permet d'éliminer totalement la nécessité de
réaliser des tests de typage, on dit qu'il est \emph{fort}. Mais ce n'est que
rarement le cas. En effet, il peut y avoir des constructions au sein du langage
qui permettent de contourner le système de types, comme un opérateur de
transtypage\ref{fig:javacast}. À l'exécution, une erreur de types est levée :

\begin{figure}
  \insertcode{cast.java}
  \caption{Transtypage en Java}
  \label{fig:javacast}
\end{figure}

\begin{Verbatim}
Exception in thread "main" java.lang.ClassCastException:
    java.lang.Integer cannot be cast to java.lang.Float
        at Cast.main(Cast.java:5)
\end{Verbatim}

\paragraph{Polymorphisme :} parfois, il est trop restrictif de donner un unique
contrat à une fonction. Quel doit être le type d'une fonction ajoutant un
élément à une liste ?

En première approximation, on peut imaginer fournir une version du code par type
de données à manipuler. C'est la solution retenue par les premières versions du
langage Pascal, ce qui rendait très difficile l'écriture de
bibliothèques\cite{PascalNoEscape}.

Une amélioration peut être de générer à la compilation une fonction qui aurait
le même code mais des annotations de typage différentes : on obtient alors
plusieurs copies de la fonction qui sont chacunes spécialisées pour un type en
particulier. C'est la solution retenue par le langage C++ et son système de
\emph{templates}.

En fait, si à l'exécution toutes les données sont représentées de manière
uniforme, il n'est pas nécessaire de créer plusieurs copies de la fonction.

\begin{figure}
  \insertcode{listappend.ml}
  \caption{Fonction de concaténation de listes en OCaml.}
  \label{fig:listappend}
\end{figure}

Par exemple, la fonction de la figure~\ref{fig:listappend} n'opère que sur la
structure du type liste (en utilisant ses constructeurs \texttt{{[}{]}} et
\listcons ainsi que le filtrage) : les éléments de \texttt{lx} et \texttt{ly} ne
sont pas manipulés à part pour les transférer dans le résultat.

Elle peut donc être typée avec n'importe quelle type 
$a~\textrm{list} → a~\textrm{list} → a~\textrm{list}$. Pour introduire cette
généricité, on modifie le système de types pour transformer :

\[ ∀ a. \textrm{append} :
   a~\textrm{list} → a~\textrm{list} → a~\textrm{list}
\]

en :

\[ \textrm{append} : ∀ a.
   a~\textrm{list} → a~\textrm{list} → a~\textrm{list}
\]

Au lieu d'associer à chaque expression un type, dans certains cas on lui associe
un schéma de types, instanciable en un type concret. En quelque sorte, on
déplace la vérification du site d'appel au site de définition. Cette technique a
été décrite en premier dans~\cite{Milner78}.
Un tour d'horizon de différents types statiques, avec en particulier du
polymorphisme, on pourra se référer à~\cite{TAPL}.

\paragraph{Expressivité :} les systèmes de types les plus simples expriment des
contrats esssentiellement liés à la sûreté d'exécution, pour ne pas utiliser des
valeurs de types incompatibles entre eux. Mais il est possible d'étendre le
langage avec des annotations plus riches : par exemple en vérifiant statiquement
que des listes ne sont pas vides\cite{lightweight-static-capabilities}, ou dans
le domaine de la sécurité, d'empêcher des fuites d'information~\cite{LZ06a}.

\paragraph{Qualificateurs de types :} Dans le cas particulier des vulnérabilités
liées à une mauvaise utilisation de la mémoire, les développeurs du noyau Linux
ont ajouté un système d'annotations au code source. Un pointeur peut être décoré
d'une annotation \texttt{\_\_kernel} ou \texttt{\_\_user} selon s'il est sûr ou
pas. Celle-ci sont ignorées par le compilateur, mais un outil d'analyse statique
ad-hoc nommé Sparse~\link{sparse}~\cite{TorvaldsSparse} % TODO citer un seul
peut être utilisé pour détecter les cas les plus simples d'erreurs.
Il demande aussi au programmeur d'ajouter beaucoup d'annotations dans le
programme. % TODO quantifier

Ce système d'annotations sur les types a été formalisé sous le nom de
\emph{qualificateurs de types} : chaque type peut être décoré d'un ensemble de
qualificateurs (à la manière de \texttt{const}), et des règles de typage
permettent d'établir des propriétés sur le programme. Ces analyses ont été
implantée dans l'outil CQual~\cite{toplas-quals}. Ce système peut servir à
inférer les annotations \texttt{const}~\cite{pldi99}, à l'analyse de souillure
pour les chaîne de format~\cite{usenix01} % TODO citer Newsham sur le format ?
et des propriétés dépendantes du flot de contrôle, comme des invariants sur les
verrous~\cite{pldi02}, à rapprocher du concept de \emph{typestates}
~\cite{tse12-typestate}. Il a également été appliqué à la classe de
vulnérabilités sur les pointeurs utilisateurs dont il est ici
l'objet~\cite{cquk-usenix04}. Puisqu'elle consiste à ajouter un qualificateur à
chaque étage de type, cette approche est plus générique mais plus complexe que
la notre. % TODO bof

\section{Analyse de code système}

Les logiciels système demandant des garanties de sécurité et de fiabilité, de
nombreuses analyses \emph{ad-hoc} ciblent les noyaux de systèmes d'exploitation.

Ajouter un système de types forts à C est l'idée centrale de
CCured~\cite{ccured-toplas}. Dans les cas où il n'est pas possible de conclure,
des vérifications à l'exécution sont ajoutées. Cependant, cela nécessite une
instrumentation dynamique qui est faite pour rester active, ce qui se paye en
performances.

Saturn~\cite{paste07} est un système pour analyser du code système écrit en C.
Il traite le problème des pointeurs utilisateur en utilisant une analyse de
forme ``pointe-sur''~\cite{oakland08}. Comme l'interprétation abstraite, son but
est d'être très précis, au détriment d'un temps de calcul important dans
certains cas.

\section{Logique de Hoare}

Une technique pour vérifier statiquement des propriétés sur la sémantique d'un
programme a été formalisée par Robert Floyd~\cite{FloydMeaning} et Tony
Hoare~\cite{hoare}.

Elle consiste à écrire les invariants qui sont maintenus à un point donné du
programme. Ces propositions sont écrites dans une logique $\mathcal{L}$.
Chaque instruction $i$ est annotée d'une pré-condition $P$
et d'une post-condition $Q$, ce que l'on note $\hoare{P}{i}{Q}$. Cela signifie
que si $P$ est vérifiée et que l'exécution de $i$ se termine
\footnote{
  Comme dans la plupart des cas, la vérification de la terminaison d'un
  algorithme est réalisée de manière séparée.
  % TODO faire un chapeau général façon Devie?
  % 1/ Terminaison
  % 2/ Correction
  % 3/ Complexité
}, alors $Q$ sera vérifiée.

En plus des règles de $\mathcal{L}$, des règles d'inférence traduisent la
sémantique du programme ; par exemple la règle de composition est :

\begin{mathpar}
  \irule{Hoare-Seq}
    { \hoare{P}{i_1}{Q} \\
      \hoare{Q}{i_2}{R}
    }{
      \hoare{P}{i_1;i_2}{R}
    }
\end{mathpar}

Les pré-conditions peuvent être renforcées et les post-conditions relâchées :

\begin{mathpar}
  \irule{Hoare-Consequence}
    { ⊢_{\mathcal{L}} P  ⇒ P' \\
      \hoare{P}{i}{Q} \\
      ⊢_{\mathcal{L}} Q' ⇒ Q
    }
    { \hoare{P'}{i}{Q'} }
\end{mathpar}

Il est alors possible d'annoter le programme avec ses invariants formalisés de
manière explicite dans $\mathcal{L}$. Ceux-ci seront vérifiés à la compilation.
% TODO à la compil?

La règle de conséquence permet de découpler les propriétés du programme lui-même
: plusieurs niveaux d'annotations sont possibles, du moins précis au plus
précis. En fait, il est même possible d'annoter chaque point de contrôle par
l'ensemble d'annotations vide : \hoare{T}{i}{T} est toujours vrai.

Augmenter graduellement les pré- et post-conditions est néanmoins assez
difficile, puisqu'il peut être nécessaire de modifier l'ensemble des conditions
à la fois.

Cette difficulté est mentionnée dans \cite{cssv}, où un système de programmation
par contrats est utilisé pour vérifier la correction de routines de manipulation
de chaînes en C.

Ce type d'annotations a été implanté par exemple pour le langage Java dans le
système JML\cite{jmlkluwer} ou pour le langage C\# dans Spec\#\cite{krml136}.

\section{Proposition}

% TODO

Le but de cette thèse est de montrer que les techniques de typage statique
permettent de vérifier que les pointeurs utilisateurs sont manipulés d'une
manière qui préserve l'intégrité du noyau d'un système d'exploitation.

% vim: spelllang=fr


\partccl{part:ctx}

Nous avons montré que l'écriture de noyaux de systèmes d'exploitation nécessite
de manipuler des données provenant d'une zone non sûre, l'espace utilisateur.
Parmi ces données, il arrive de récupérer des pointeurs qui servent à passer des
données par référence à l'appelant, dans certains appels système. Si on
déréférence ces pointeurs sans vérifier qu'ils pointent bien vers une zone
mémoire également contrôlée par l'appelant, on risque de lire ou d'écrire dans
des zones mémoires réservées au noyau seul.

Nous proposons une technique de typage pour détecter ces cas dangereux. Pour ce
faire, il faudra tout d'abord définir un langage impératif bien typé que nous
appellerons \langname.
Celui-ci s'appuie sur le langage \newspeak, qui est un langage intermédiaire
développé par EADS dans le but de vérifier la sûreté de programmes C embarqués.
À ce titre, il existe un compilateur qui est capable de traduire du code C vers
\newspeak.

Une fois \langname défini, nous lui ajoutons des constructions permettant
d'écrire du code noyau, et en particulier on lui ajoute des pointeurs
utilisateur: il s'agit de pointeurs dont la valeur est contrôlée par un
utilisateur interagissant avec le programme via un appel système. Ces pointeurs
ont un type distinct des pointeurs habituels.

Cela revient à introduire une séparation entre le type des données et sa
représentation, c'est-à-dire définir un type abstrait. Les pointeurs
utilisateur sont en fait des pointeurs classiques déguisés, pour lesquels on
interdit l'opérateur de déréférencement.

Cette technique est en fait générique : on peut également l'appliquer à certains
types d'entiers. En C, il est commun d'utiliser des \texttt{int} pour tout et
n'importe quoi : pour des entiers bien sûr (au sens de $ℤ$), mais aussi comme
identificateurs pour lesquels les opérations usuelles comme l'addition n'ont pas
de sens. Par exemple, sous Linux, l'opération d'ouverture de fichier renvoie un
entier, dit \emph{descripteur de fichier}, qui identifie ce fichier pour ce
processus.
Le langage autorise donc par exemple de multiplier entre
eux deux descripteurs de fichiers, mais le résultat n'a pas de raison \emph{a
priori} d'être un descripteur de fichier valide.

En n'offrant pas cette distinction, le langage C permet d'écrire du code qui
peut s'exécuter mais dont la sémantique n'est, quelque part, pas bien fondée. En
effet, le système de types de C est trop primitif pour pouvoir garantir une
véritable isolation entre deux types de même représentation : il n'y a pas de
types abstraits. Certes, \texttt{typedef} permet d'introduire un nouveau nom
pour un type, mais ce n'est qu'un raccourci syntaxique. Le compilateur ne peut
en effet pas considérer un programme sans avoir la définition quasi-complète des
types qui y apparaissent. La seule exception concerne les pointeurs sur
structures : si on ne fait que les affecter, il n'est pas nécessaire de
connaître la taille ni la disposition de la structure, donc la définition peut
ne pas être visible. Cette technique, connue sous le nom de \emph{pointeurs
opaques}, n'est pas applicable aux autres types.

En ajoutant une couche de typage, on interdit ces opérations à la compilation.
Cela permet deux choses : pour le code déjà écrit, de détecter et corriger les
manipulations dangereuses ; et, pour le nouveau code, de s'assurer qu'il est
correct. Par exemple, si on écrit un éditeur de texte, on peut éviter de
nombreuses erreurs de programmation en définissant un type \enquote{indice de
ligne} et un type \enquote{indice de colonne} incompatibles entre eux.

Le but de cette thèse est donc de définir un langage intermédiaire proche de C,
mais bien typable, et de lui adjoindre un système de types tel que les
programmes bien typés manipulent les pointeurs utilisateur sans causer de
problèmes de sécurité.

% vim: spelllang=fr


\begin{headingpage}
Dans cette partie, nous allons présenter un langage impératif modélisant le
langage C. Le chapitre~\ref{cha:lang} décrit sa syntaxe, ainsi que sa
sémantique. À ce point, de nombreux programmes sont acceptés mais qui provoquent
des erreurs à l'exécution.

Afin de rejeter ces programmes incorrects, on définit ensuite dans le
chapitre~\ref{cha:typbase} une sémantique statique s'appuyant sur un système de
types simples. Des propriétés de sûreté de typage sont ensuite établies,
permettant de catégoriser l'ensemble des erreurs à l'exécution possibles.

Le chapitre~\ref{cha:qualifs} commence par étendre notre langage avec une
nouvelle classe d'erreurs à l'exécution, modélisant les accès à la mémoire
utilisateur catégorisé comme dangereux dans le chapitre~\ref{cha:os}. Une
extension au système de types du chapitre~\ref{cha:typbase} est ensuite établie,
et on prouve que les programmes ainsi typés ne peuvent pas atteindre ces cas
d'erreur.

Trois types d'erreurs à l'exécution sont possibles:

\begin{itemize}
\item
  les erreurs de typage (dynamique), lorsqu'on tente d'appliquer à une
  opération des valeurs incompatibles (additionner un entier et une
  fonction par exemple).
\item
  les erreurs de sécurité, qui consistent en le déréférencement d'un
  pointeur dont la valeur est contrôlée par l'espace utilisateur.
  Celles-ci sont uniquement possibles en contexte noyau.
\item
  les erreurs mémoire, qui résultent d'un débordement de tableau, du
  déréférencement d'un pointeur invalide ou d'arithmétique de pointeur
  invalide.
\end{itemize}

En résumé, l'introduction des types simples enlève la possibilité de rencontrer
des erreurs de typage dynamique, et l'ajout des qualificateurs interdit les
erreurs de sécurité.

\begin{center}
\begin{tabular}{ll@{\hskip 15mm}ccc}
\toprule
\multirow{2}{*}{Langage} & \multirow{2}{*}{Types}  & \multicolumn{3}{c}{Erreurs possibles}   \\
                         &  & Typage      & Sécurité    & Mémoire     \\
\midrule
\langname{}       & sans      & \CheckedBox{} & N/A           & \CheckedBox{} \\
\langname{}       & simples   & \Square{}     & N/A           & \CheckedBox{} \\
\langname{} noyau & simples   & \Square{}     & \CheckedBox{} & \CheckedBox{} \\
\langname{} noyau & qualifiés & \Square{}     & \Square{}     & \CheckedBox{} \\
\bottomrule
\end{tabular}

% TODO ligne 3 s/simples/naifs/ ?

\end{center}
\end{headingpage}

\part{Un langage pour l'analyse de code système : \langname}
\label{part:lang}


\chapter{Syntaxe et sémantique}

\label{cha:lang}
Dans ce chapitre, on décrit le support de notre travail: un langage impératif
nommé \langname, sur lequel s'appuieront les analyses de typage des
chapitres~\ref{cha:typbase} et~\ref{cha:qualifs}.

Le langage C~\cite{KandR} est un langage impératif, conçu pour être un
\enquote{assembleur portable}. Ses types de données et les opérations associées
sont donc naturellement de très bas niveau.

Ses types de données sont établis pour représenter les mots mémoire manipulables
par les processeurs : essentiellement des entiers et flottants de plusieurs
tailles. Les types composés correspondent à des zones de mémoire contigües,
homogènes (dans le cas des tableaux) ou hétérogènes (dans le cas des
structures).

Une des spécificités de C est qu'il expose au programmeur la notion de pointeur,
c'est-à-dire des variables qui représentent directement une adresse en mémoire.
Les pointeurs peuvent être typés (on garde une indication sur le type de l'objet
stocké à cette adresse) ou \enquote{non typés}. Dans ce dernier cas, ils ont en
fait le type \texttt{void *} qui est compatible avec n'importe quel type
pointeur.

\added{Dans ce dernier cas...}

Son système de types rudimentaire ne permet pas d'avoir beaucoup de garanties
sur la sûreté du programme. En effet, aucune vérification n'est effectuée en
dehors de celles faites par le programmeur.

Le but ici est de définir \langname, un langage plus simple mais qui permettra
de raisonner sur une certaine classe de programmes C.

Tout d'abord, on commence par présenter les notations qui accompagneront le
reste des chapitres. Cela inclut la notion de lentille, qui est utilisée pour
définir les accès profonds à la mémoire. Cela permet de résoudre le problème de
mettre à jour une sous-valeur (par exemple un champ de structure) d'une
variable, par exemple. Les lentilles permettent de définir de manière
déclarative que pour faire cette opération, il faut obtenir l'ancienne valeur de
la variable, puis calculer une nouvelle valeur en remplaçant une sous-valeur,
avant de replacer cette nouvelle valeur à sa place en mémoire. En pratique, on
définira deux lentilles: une qui relie un état mémoire à la valeur d'une
variable, et une qui relie une valeur à une de ses sous-valeurs. Avec cette
technique, on peut définir en une seule fois les opérations de lecture et
d'écriture de sous-valeurs imbriquées.

Ensuite, on présente \langname en soi, c'est-à-dire sa syntaxe, ainsi qu'une
description de ses caractéristiques principales. En particulier, le modèle
mémoire est détaillé, ainsi que les valeurs manipulées par le langage.

Enfin, on décrit une sémantique opérationnelle pour ce langage. Cela permet de
définir précisément l'exécution d'un programme \langname au niveau de la
mémoire.

L'implantation de ces analyses est faite dans le chapitre~\ref{cha:implem}.
Puisque \langname n'est qu'un modèle, il s'agira d'adapter ces règles de typage
sur \newspeak, qui possède un modèle mémoire plus bas niveau.

\section{Notations}

\subsection*{Inférence}

La sémantique opérationnelle consiste en la définition d'une relation de
transition $\cdot\rightarrow\cdot$ entre états de l'interprète\footnote{Dans le
chapitre~\ref{cha:typbase}, la relation de typage $\cdot ⊢ \cdot : \cdot$ sera
définie par la même technique.}.

Cette relation est définie inductivement sur la syntaxe du programme. Plutôt que
de présenter l'induction explicitement, elle est représentée par des jugements
logiques et des règles d'inférence, de la forme:

\[
\irule{Nom}{P_1 \\ … \\ P_n}{C}
\]

Les $P_i$ sont les prémisses, et $C$ la conclusion. Cette règle s'interprète de
la manière suivante: si les $P_i$ sont vraies, alors $C$ est vraie.

Certaines règles n'ont pas de prémisse, ce sont des axiomes:

\[
\iaxiom{Ax}{A}
\]

Compte-tenu de la structure des règles, la dérivation d'une preuve (l'ordre dans
lequel les règles sont appliquées) pourra donc être vue sous la forme d'un arbre
où les axiomes sont les feuilles, en haut, et la conclusion est la racine, en
bas.

\[
  \irule{r1}{
    \irule{r2}
          {
            \iaxiom{r3}{A_1}
              \\
            \iaxiom{r4}{A_2}
          }
          {B_1}
    \\
    \irule{r5}
      {
        \iaxiom{r6} {A_3}
      }{B_2}
      }{C}
\]

\subsection*{Listes}
\label{page:def-listes}

$\sList{X}$ est l'ensemble des suites finies de $X$, indexées à partir de 1. Si
$u ∈ \sList{X}$, on note $|u|$ le nombre d'éléments de $u$ (le cardinal de son
domaine de définition). Pour $i ∈ [1 ; |u|]$, on note $u_i = u(i)$ le i-ème
élément de la suite.

On peut aussi voir les suites comme des listes: on note $[]$ la suite vide,
telle que $|[]| = 0$. On définit en outre la construction de suite de la manière
suivante: si $x ∈ X$ et $u ∈ \sList{X}$, la liste $x::u ∈ \sList{X}$ est la
liste $v$ telle que:

\begin{align*}
                       v_1 & = x \\
  ∀ i ∈ [1; |u|] , v_{i+1} & = u_i
\end{align*}

Cela signifie que la tête de liste ($x$ dans la liste $x :: u$) est toujours
accessible à l'indice $1$.

\subsection*{Lentilles}

Dans la définition de la sémantique de \langname, on utilise des \emph{lentilles
bidirectionnelles}. Cette notion n'est pas propre à la sémantique des
programmes. Il s'agit d'une technique permettant de relier la modification d'un
objet à la modification d'un de ses sous-composants. Cela a plusieurs
applications possibles. En programmation fonctionnelle pure (sans mutation), on
ne peut pas mettre à jour partiellement les valeurs composées comme des
enregistrements (\emph{records}). Pour simuler cette opération, on a en général
une opération qui permet de définir un nouvel enregistrement dans lequel seul un
champ a été mis à jour. C'est ce qui se passe avec le langage
Haskell~\cite{rwh}: \texttt{r \{ x = 5 \}} représente une valeur enregistrement
égale à \texttt{r} sur tous les champs, sauf pour le champ \texttt{x} où elle
vaut \texttt{5}. Utiliser des lentilles revient à ajouter dans le langage la
notion de champ en tant que valeur de première classe. Elles ont l'avantage de
pouvoir se composer, c'est-à-dire que si on a un champ nommé $x$ qui contient un
champ nommé $y$, alors on peut modifier le champ du champ automatiquement.

Dans ce cadre, les lentilles ont été popularisées par Van
Laarhoven~\cite{LaarhovenLenses}. Puisque cela sert à manipuler des données
arborescentes, on peut aussi appliquer cet outil aux systèmes de bases de
données ou aux documents structurés comme par exemple en
XML~\cite{PierceLenses}.

Dans notre cas, cela permettra par exemple de modifier un élément d'un tableau
qui est un champ de structure de la variable nommée $x$ dans le 3\ieme cadre de
pile.

\begin{definition}[Lentille]

Étant donnés deux ensembles $R$ et $A$, une \emph{lentille} $ℒ ∈ \setLens{R}{A}$
(ou \emph{accesseur}) est un moyen d'accéder en lecture ou en écriture à une
sous-valeur appartenant à $A$ au sein d'une valeur appartenant à $R$ (pour
\emph{record}). Elle est consistuée des opérations suivantes:

\begin{itemize}
\item
  une fonction de lecture\enspace$\mathrm{get}_ℒ : R → A$
\item
  une fonction de mise à jour\enspace$\mathrm{put}_ℒ : (A × R) → R$
\end{itemize}

telles que pour tous $a∈A, a'∈A, r∈R$:

\begin{align*}
\tag{\textsc{GetPut}}
\mathrm{put}_ℒ(\mathrm{get}_ℒ(r), r) &= r \\
\tag{\textsc{PutGet}}
\mathrm{get}_ℒ(\mathrm{put}_ℒ(a, r)) &= a \\
\tag{\textsc{PutPut}}
\mathrm{put}_ℒ(a', \mathrm{put}_ℒ(a, r)) &= \mathrm{put}_ℒ(a', r) \\
\end{align*}

On note $ℒ =
\mkLens{\mathrm{get}_{ℒ}}{\mathrm{put}_{ℒ}}$.

\textsc{GetPut} signifie que si on lit une valeur puis qu'on la réécrit, l'objet
n'est pas modifié; \textsc{PutGet} décrit l'opération inverse: si on écrit
une valeur dans le champ, c'est la valeur qui sera lue; enfin, \textsc{PutPut}
évoque le fait que chaque écriture est totale: quand deux écritures se suivent,
seule la seconde compte.

\end{definition}

Une illustration se trouve dans la figure~\ref{fig:lens-howto}.

\begin{figure}[h]

\tikzstyle{lenspatA}=[fill=white,pattern=crosshatch dots]
\tikzstyle{lenspatB}=[fill=blue!50]
\tikzstyle{lenspatC}=[fill=red!20]

  \begin{align*}
  \lensGet{ℒ}{\lensNodeBig{lenspatA}{lenspatB}} &= \lensInner{lenspatB} \\
  \lensPut{ℒ}{\lensInner{lenspatC}}{\lensNodeBig{lenspatA}{lenspatB}}
        &= \lensNodeBig{lenspatA}{lenspatC} \\
  \end{align*}

\caption{Fonctionnement d'une lentille}
\label{fig:lens-howto}
\end{figure}

\begin{example}[Lentilles de tête et de queue de liste]

Soit $E$ un ensemble. On rappelle que $\sList{E}$ désigne l'ensemble des listes
d'éléments de $E$.

\changed{$L(E)$ devient $\sList{E}$ pour rester cohérent}

On définit les fonctions suivantes. Notons qu'elles ne sont pas définies sur la
liste vide $[]$, qui pourra être traitée comme un cas d'erreur.

\begin{align*}
  \mathrm{get}_T     (t::q) &= t   & \mathrm{put}_T (t', t::q) &= t'::q  \\
  \mathrm{get}_Q     (t::q) &= q   & \mathrm{put}_Q (q', t::q) &= t::q'
\end{align*}

Alors
$T = \mkLens{\mathrm{get}_T}{\mathrm{put}_T} ∈ \setLens{\sList{E}}{E}$
et
$Q = \mkLens{\mathrm{get}_Q}{\mathrm{put}_Q} ∈ \setLens{\sList{E}}{\sList{E}}$.

On a par exemple:

\begin{align*}
\mathrm{get}_T (1::6::1::8::[]) &= 1   &   \mathrm{put}_Q (4::2::[], 3::6::1::5::[]) = 3::4::2::[]
\end{align*}

\end{example}

\begin{definition}[Lentille indexée]

Les objets de certains ensembles $R$ sont composés de plusieurs sous-objets
accessibles à travers un indice $i ∈ I$. Une lentille indexée est une fonction
$Δ$ qui associe à un indice $i$ une lentille entre $R$ et un de ses champs
$A_i$:

\[
  ∀ i ∈ I, ∃ A_i, Δ(i) ∈ \setLens{R}{A_i}
\]

On note alors:

\begin{align*}
    r {[ i ]}_Δ     & \eqdef \mathrm{get}_{Δ(i)}(r) \\
    r {[ i ← a ]}_Δ & \eqdef \mathrm{put}_{Δ(i)}(a, r) \\
\end{align*}

\end{definition}

Un exemple est illustré dans la figure~\ref{fig:lens-idx-ex}.

\begin{figure}[h]

% N&B
%\tikzstyle{lenspatA}=[fill=white,pattern=crosshatch dots]
%\tikzstyle{lenspatB}=[fill=black!50]
%\tikzstyle{lenspatC}=[fill=black!20]
\tikzstyle{lenspatA}=[fill=white,pattern=crosshatch dots]
\tikzstyle{lenspatB}=[fill=blue!50]
\tikzstyle{lenspatC}=[fill=red!20]

  \begin{align*}
    \lensGet{Δ(b)}{\lensNodeBigIdx{lenspatA}{lenspatB}{lenspatB}{lenspatB}{lenspatB}} &= \lensInnerStar{lenspatB} \\
    \lensGet{Δ(c)}{\lensNodeBigIdx{lenspatA}{lenspatB}{lenspatB}{lenspatB}{lenspatB}} &= \lensInner{lenspatB} \\
    \lensPut{Δ(b)}{\lensInnerStar{lenspatC}}{\lensNodeBigIdx{lenspatA}{lenspatB}{lenspatB}{lenspatB}{lenspatB}} &=
      \lensNodeBigIdx{lenspatA}{lenspatB}{lenspatC}{lenspatB}{lenspatB} \\
    \lensPut{Δ(c)}{\lensInner{lenspatC}}{\lensNodeBigIdx{lenspatA}{lenspatB}{lenspatB}{lenspatB}{lenspatB}} &=
      \lensNodeBigIdx{lenspatA}{lenspatB}{lenspatB}{lenspatC}{lenspatB}
  \end{align*}

\caption{Fonctionnement d'une lentille indexée}
\label{fig:lens-idx-ex}
\end{figure}

\begin{example}[Lentille \enquote{$n$\ieme élément d'un tuple}]

Soient $n ∈ ℕ$, et $n$ ensembles $E_1, …, E_n$.

Pour tout $i ∈ [1; n]$, on définit:

\begin{align*}
   g_i((x_1, …, x_n)) &= x_i \\
p_i(y, (x_1, …, x_n)) &= (x_1, …, x_{i-1}, y, x_{i+1}, …, x_n)\\
\end{align*}

Définissons $T(i) = \mkLens{g_i}{p_i}$. Alors $T(i) ∈ \setLens{(E_1×…×E_n)}{E_i}$.

Donc $T$ est une lentille indexée, et on a par exemple:

\begin{align*}
(3,1,4,1,5) {[2]}_T &= \mathrm{get}_{T(2)} ((3, 1, 4, 1, 5)) \\
                    &= 1 \\
\\
(9,2,6,5,3) {[3 ← 1]}_T &= \mathrm{put}_{T(3)} (1, (9,2,6,5,3)) \\
                        &= (9,2,1,5,3)
\end{align*}
\end{example}

La notation $3 ← 1$ peut surprendre, mais elle est à interpréter comme
\enquote{en remplaçant l'élément d'indice $3$ par $1$}.

\begin{definition}[Composition de lentilles]
\label{def:lens-comp}

  Soient $ℒ_1 ∈ \setLens{A}{B}$ et $ℒ_2 ∈ \setLens{B}{C}$.

  La composition de $ℒ_1$ et $ℒ_2$ est la lentille
  $ℒ ∈ \setLens{A}{C}$ définie de la manière suivante:

  \begin{align*}
    \mathrm{get}_{ℒ} (r) &=
        \mathrm{get}_{ℒ_2}
        (\mathrm{get}_{ℒ_1} r) \\
    \mathrm{put}_{ℒ} (a, r) &=
        \mathrm{put}_{ℒ_1} (\mathrm{put}_{ℒ_2}
        (a, \mathrm{get}_{ℒ_1} r), r) \\
  \end{align*}

  On notera alors $ℒ = ℒ_1 \ggg ℒ_2$ (\enquote{$ℒ_1$ flèche $ℒ_2$}).

\end{definition}

\begin{figure}[h]

\centering

%N&B
%\tikzstyle{lensInnerBefore}=[fill=black!85]
%\tikzstyle{lensInnerAfter}=[fill=black!10]
%\tikzstyle{lensMid}=[fill=black!50]
%\tikzstyle{lensExt}=[fill=white,pattern=crosshatch dots]

\tikzstyle{lensInnerBefore}=[fill=green!85]
\tikzstyle{lensInnerAfter}=[fill=orange!30]
\tikzstyle{lensMid}=[fill=blue!50]
\tikzstyle{lensExt}=[fill=white,pattern=crosshatch dots]

\shorthandoff{!}
\vspace{5mm}
\begin{tikzpicture}
[node distance=2cm
,bignode/.style={draw,shape=rectangle,minimum size=1cm,lensExt}
,smallnode/.style={draw,shape=circle,minimum size=8mm,lensMid}
,trinode/.style={draw,shape=regular polygon,regular polygon sides=3}
]

\node[bignode] (A) {};
\node[right of=A,smallnode] (B) {};
\node[right of=B,trinode,lensInnerAfter] (C) {};
\node[below of=B,smallnode] (D) {};
\node[below of=A,node distance=4cm,bignode] (E) {};

\draw[->] (A) to node[auto] {$\mathrm{get}_{ℒ_1}$} (B);
\draw[->] (B) to (D);
\draw[->] (A) to (E);
\draw (C) |- node[auto,near end,swap] {$\mathrm{put}_{ℒ_2}$} ($(B)!0.5!(D)$);
\draw (D) |- node[auto,near end,swap] {$\mathrm{put}_{ℒ_1}$} ($(A)!0.75!(E)$);

\node[smallnode] at (A) {};
\node[trinode,lensInnerBefore] at (A) {};

\node[trinode,lensInnerBefore] at (B) {};

\node[trinode,lensInnerAfter] at (D) {};

\node[smallnode] at (E) {};
\node[trinode,lensInnerAfter] at (E) {};

\node[left of=A, node distance=4cm,bignode] (GA) {};
\node[below of=GA,smallnode] (GB) {};
\node[below of=GB,trinode,lensInnerBefore] (GC) {};

\node[smallnode] at (GA) {};

\node[trinode, lensInnerBefore] at (GA) {};
\node[trinode, lensInnerBefore] at (GB) {};

\draw [->] (GA) to node[auto] {$\mathrm{get}_{ℒ_1}$} (GB);
\draw [->] (GB) to node[auto] {$\mathrm{get}_{ℒ_2}$} (GC);

\end{tikzpicture}
\vspace{5mm}
\shorthandoff{!}

\caption{Composition de lentilles}
\label{fig:compo-lens}
\end{figure}

Cette définition est illustrée dans la figure~\ref{fig:compo-lens}. Une preuve
que la composition est une lentille est donnée en annexe~\ref{proof:compo-lens}.

\section{Syntaxe}

Les figures~\ref{fig:stx-data} et~\ref{fig:stx} présentent
notre langage intermédiaire. Il contient la plupart des fonctionnalités
présentes dans les langages impératifs comme C.

\begin{figure}%{{{

  \figstxdata{}

  \caption{Syntaxe des expressions}
\label{fig:stx-data}
\end{figure}%}}}

Parmi les expressions, les constantes comportent les entiers et flottants, ainsi
que le pointeur \eNull qui correspond à une valeur par défaut pour les
pointeurs, et la valeur unité \eUnit qui pourra être retournée par les fonctions
travaillant par effets de bord uniquement.

\changed{Partout: lvalues devient valeurs gauches}
Les accès mémoire en lecture et écriture se font au travers de valeurs gauches
(\emph{left values} ou \emph{lvalues}): comme en C, elles tiennent leur nom du
fait que ce sont ces constructions qui sont à gauche du signe d'affectation. En plus
des variables, on obtient une valeur gauche en accédant par nom à un champ ou par
indice à un élément d'une valeur gauche, ou encore en appliquant l'opérateur
\texttt{*} de déréférencement à une expression. Pour assister le typage, l'accès
à un champ doit être décoré du type complet $S$, mais cette annotation est
ignorée lors de l'évaluation. Les valeurs gauches correspondent aussi à l'unité
d'adressage: c'est-à-dire que les pointeurs sont construits en prenant l'adresse
d'une valeur gauche avec l'opérateur~\texttt{\&}.

Les fonctions sont des expressions comme les autres, contrairement à C où elles
sont forcément déclarées globalement. Cela veut dire qu'on peut affecter une
fonction $f$ à une variable $x$ et l'appeller avec $x(a_1, a_2)$. Il est aussi
possible de déclarer une fonction au sein d'une fonction. Cependant cela ne
respecte pas l'imbrication lexicale: dans la fonction interne il n'est pas
possible de faire référence à des variables locales de la fonction externe,
seulement à des variables globales. En mémoire les fonctions sont donc
uniquement représentées par leur code: il n'y a pas de fermetures.

Enfin, on trouve aussi des expressions permettant de construire des valeurs
composées: les structures et les tableaux.

\begin{figure}%{{{

  \figstxinstr{}

  \caption{Syntaxe des instructions}
\label{fig:stx}
\end{figure}%}}}

Les instructions sont typiques de la programmation impérative. \langname
comporte bien sûr l'instruction vide qui ne fait rien et la séquence qui chaîne
deux instructions.

Une expression peut être évaluée dans un contexte d'instruction, pour ses effets
de bord. Remarquons que l'affectation est une expression, qui renvoie la valeur
affectée. Cela permet d'écrire $x ← (y ← z)$, comme dans un programme C où on
écrirait \texttt{x = y = z}.

Il est également possible de déclarer une variable locale avec
$\iDecl{x}{v}{i}$. $x$ est alors une nouvelle variable visible dans $i$ avec
pour valeur initiale $v$.

L'alternative et la conditionnelle sont classiques; en revanche, on ne fournit
qu'un seul type de boucle et pas de saut (instruction \texttt{goto}).

Les opérateurs sont donnés dans la figure~\ref{fig:stx-ops}. Ils correspondent à
ceux du langage C. La différence principale est que les opérations sur les
entiers, flottants et pointeurs sont annotés avec le type de données sur lequel
ils travaillent. Par exemple \enquote{$+$} désigne l'addition sur les entiers et
\enquote{$+.$} l'addition sur les flottants. Les opérations de test d'égalité,
en revanche, sont possibles pour les types numériques, les pointeurs, ainsi que
les types composés de types comparables.

\begin{figure}[h]%{{{

  \figstxops{}

  \caption{Syntaxe des opérateurs}
\label{fig:stx-ops}
\end{figure}%}}}

\section{Mémoire et valeurs}

L'interprète que nous nous apprêtons à définir manipule des valeurs qui sont
associées aux variables du programme.

La mémoire est constituée de variables (toutes mutables), qui contiennent des
valeurs. Ces variables sont organisées, d'une part en un ensemble de variables
globales, et d'autre part en une pile de contextes d'appel (qu'on appellera donc
aussi cadres de pile, ou \emph{stack frames} en anglais). Cette structure
empilée permet de représenter les différents contextes à chaque appel de
fonction: par exemple, si une fonction s'appelle récursivement, plusieurs
instances de ses variables locales sont présentes dans le programme. Le modèle
mémoire présenté ici de permet pas l'allocation dynamique sur un tas. Cette
limitation sera détaillée dans le chapitre~\ref{cha:conclusion}.

\added{Le modèle...}

La structure de pile des variables locales permet de les organiser en niveaux
indépendants: à \linebreak chaque appel de fonction, un nouveau cadre de pile
est créé, comprenant ses paramètres et ses variables locales. Au contraire, pour
les variables globales il n'y a pas de système d'empilement, puisque ces
variables sont accessibles depuis tout point du programme.

Pour identifier de manière non ambigüe une variable, on note simplement $x$ la
variable globale nommée $x$, et $(n, x)$ la variable locale nommée $x$ dans le
$n$\ieme cadre de pile\footnote{Les paramètres de fonction sont traités comme
des variables locales et se retrouvent dans le cadre correspondant.}.

Les affectations peuvent avoir la forme $x ← e$ où $x$ est une variable et $e$
est une expression, mais pas seulement. En effet, à gauche de $←$ on trouve en
général non pas une variable mais une valeur gauche (par définition). Pour
représenter quelle partie de la mémoire doit être accédée par cette valeur
gauche, on introduit la notion de chemin $φ$. Un chemin est une valeur gauche
évaluée: les cas sont similaires, sauf que tous les indices sont évalués. Par
exemple, $φ = (5, x).p$ représente le champ \enquote{$p$} de la variable $x$
dans le 5\ieme cadre de pile. C'est à ce moment qu'on évalue les
déréférencements qui peuvent apparaître dans une valeur gauche.

Les valeurs, quant à elles, peuvent avoir les formes suivantes (résumées sur la
figure~\ref{fig:interp-val}):

\begin{itemize}
\item

$\widehat{c}$: une constante. La notation circonflexe permet de distinguer
les constructions syntaxique des constructions sémantiques. Par exemple, à la
syntaxe $3$ correspond la valeur $\widehat{3}$.

\label{page:entiers-bits}
Les valeurs entières sont les entiers signés sur 32 bits, c'est-à-dire entre
$-2^{31}$ à $2^{31}-1$. Mais ce choix est arbitraire: on
aurait pu choisir des nombres à 64 bits, par exemple.
Les flottants sont les flottants IEEE 754 de 32 bits~\cite{ieee754}.

Il n'y a pas de distinction entre procédures et fonctions; toutes les fonctions
doivent renvoyer une valeur. Celles qui ne retournent pas de valeur
\enquote{intéressante} renvoient alors une valeur d'un type à un seul élément
noté \eUnit, et donc le type sera noté \tUnit. Cette notation évoque un
$n$-uplet à 0 composante.

\item

$\widehat{\&}~φ$: une référence mémoire. Ce chemin correspond à un pointeur sur
une valeur gauche. Par exemple, l'expression $\&x$ s'évalue en $\widehat{\&}~φ =
\widehat{\&}~(5, x)$ si $x$ désigne lexicalement une variable dans le 5\ieme
cadre de pile.

\item

$\widehat{ \eArray{v_1 ;…; v_n} }$: un tableau. C'est une valeur composée qui
contient un certain nombre (connu à la compilation) de valeurs d'un même type,
par exemple 100 entiers. On accède à ces valeurs par un indice entier. C'est une
erreur ($\serr{array}$) d'accéder à un tableau en dehors de ses bornes,
c'est-à-dire en dehors de $[0;n-1]$ pour un tableau à $n$ éléments.
Pareillement, $\widehat{\eArray{\cdot}}$ permet de désigner les valeurs tableau.
Par exemple, si $x$ vaut $2$ et $y$ vaut $3$, l'expression $\eArray{ x, y }$
s'évaluera en la valeur $\widehat{\eArray{ 2, 3 }}$

\item

$\widehat{ \eStruct{ l_1 : v_1 ; … ; l_n : v_n } }$: une structure. C'est une
valeur composée mais hétérogène. Les différents éléments (appelés \emph{champs})
sont désignés par leurs noms $l_i$ (pour \emph{label}). Dans le programme, le
nom de champ $l_i$ est décoré de la définition complète de la structure $S$.
Celle-ci n'est pas utilisée dans l'évaluation et sera décrite au
chapitre~\ref{cha:typbase}. Comme précédemment, on note
$\widehat{\eStruct{\cdot}}$ pour dénoter les valeurs.

\item

$\widehat{f}$: une fonction. On garde en mémoire l'intégralité de la définition
de la fonction (liste de paramètres, de variables locales et corps). Même si les
fonctions locales sont possibles, il n'est pas possible d'accéder aux variables
de la portée entourante depuis la fonction intérieure (il n'y a pas de
fermetures). Contrairement à C, les fonctions ne sont pas des cas spéciaux. Par
exemple, les fonctions globales sont simplement des variables globales de type
fonctionnel, et les \enquote{pointeurs sur fonction} de C sont remplacés par des
variables de type fonction.

\item $Ω$: une erreur. Par exemple le résultat l'évaluation de $5 / 0$ est
$\serr{div}$.

\end{itemize}

\begin{figure}%{{{

  \begin{align*}
  \gramdef{Valeurs}{v}
      { \widehat{c}     }{Constante}
      { \widehat{\&}~φ  }{Référence mémoire}
      { \widehat{
         \eStruct{ l_1 : v_1 ;
       … ; l_n : v_n }
       }                }{Structure}
      { \widehat{
        \eArray{v_1 ;…; v_n}
        }               }{Tableau}
      { \widehat{f}     }{Fonction}
      { Ω               }{Erreur}
      {END}
  \\
  \\
  \gramdef{Chemins}{φ}
     { a    }{Adresse}
     { φ\widehat{.}l  }{Accès à un champ}
     { φ\widehat{[n]} }{Accès à un élément}
     {END}
  \\
  \\
  \gramdef{Adresses}{a}
     { (n, x) }{Variable locale}
     { (x)    }{Variable globale}
     {END}
  \\
  \\
  \gramdef{Erreur}{Ω}
    { \serr{array} }{Débordement de tableau}
    { \serr{ptr}   }{Erreur de pointeur}
    { \serr{div}   }{Division par zéro}
    { \serr{field} }{Erreur de champ}
    { \serr{var}   }{Variable inconnue}
    { \serr{typ}   }{Données incompatibles}
    {END}
  \end{align*}

  \caption{Valeurs}
\label{fig:interp-val}
\end{figure}%}}}

Les erreurs peuvent être classifiées en deux grand groupes: d'une part,
$\serr{field}$, $\serr{var}$ et $\serr{typ}$ sont des erreurs de typage
dynamique, qui arrivent lorsqu'on accède dynamiquement à des données qui
n'existent pas ou qu'on manipule des types de données incompatibles. D'autre
part, $\serr{div}$, $\serr{array}$ et $\serr{ptr}$ correspondent à des valeurs
mal utilisées. Le but du système de types du chapitre~\ref{cha:typbase} sera
d'éliminer complètement les erreurs du premier groupe.

\section{Interprète}

La figure~\ref{fig:interp-stack} résume comment ces valeurs sont organisées. Une
pile est une liste de cadres de piles, et un cadre de pile est une liste de
couples (nom, valeur). Un état mémoire $m$ est un couple $(s, g)$ où $s$ est une
pile et $g$ un cadre de pile (qui représente les variables globales). On note
$|m| = |s|$ la hauteur de la pile (en nombre de cadres).

\begin{figure}%{{{

  \begin{align*}
  \gramdef{Pile}{s}
    { [] }{Pile vide}
    { \{ x_1 ↦v_1; … ; x_n ↦v_n\} :: s }{Ajout d'un cadre}
    {END}
  \\
  \\
  \gramdef{État mémoire}{m}
    { (s, \{ x_1 ↦v_1; … ; x_n ↦v_n\})}{Pile, globales}
    {END}
  \\
  \\
  \gramdef[4cm]{État d'interprète}{Ξ}
    { \msi{m}{e} }{Expression, mémoire}
    { \msi{m}{i} }{Instruction, mémoire}
    { Ω          }{Erreur}
    {END}
  \end{align*}

  \caption{Composantes d'un état mémoire}
\label{fig:interp-stack}
\end{figure}%}}}


Enfin, l'interprétation est définie comme une relation $\cdot → \cdot$ entre
états $Ξ$; ces états sont d'une des formes suivantes:

\begin{itemize}
\item
  un couple $\msi{m}{e}$ où $e$ est une expression et $m$ un état
  mémoire. $m$ est l'état mémoire sous lequel l'évaluation sera
  réalisée. Par exemple $\mm{([], [x↦3])}{3}{([], [x↦3])}{\widehat{3}}$
  L'évaluation des expressions est détaillée dans la
  section~\ref{sec:eval-exp}.
\item
  un couple $\msi{m}{i}$ où $i$ est une instruction et $m$ un état
  mémoire. La réduction des instructions est traitée dans la
  section~\ref{sec:eval-instr}.
  Par exemple, $ \msi{m}{(x ← 3 ; y ← x)} → \msi{m [x ↦ \widehat{3}]}{y ← x}
                           → \msi{m [x ↦ \widehat{3}][y ↦
    \widehat{3}]}{\iPass}$.

  Dans le cas général, utiliser des instructions pour représenter l'état des
  calculs ne suffit pas; il faut utiliser une continuation.
  C'est ce qui est fait par exemple dans la sémantique de
  CMinor~\cite{cminorSLfull}. \changed{Référence vers la version complète}
  Ici, le flot de contrôle est plus simple et on peut se contenter de retenir
  une simple instruction, ce qui simplifie la présentation.

\item
  un couple $\msi{m}{lv}$ où $lv$ est une valeur gauche et $m$ un état mémoire.
  L'évaluation des valeurs gauches est décrite en section~\ref{sec:eval-lv}.
\item
  une erreur $Ω$. La propagation des erreurs est détaillée dans la
  section~\ref{sec:eval-errors}.
\end{itemize}

L'évaluation des expressions, valeurs gauches et instructions se fait à petits
pas. C'est-à-dire qu'on simplifie d'étape en étape leur forme, jusqu'à arriver à
un cas de base:

\begin{itemize}
\item pour les expressions, une valeur $v$;
\item pour les instructions, l'instruction $\iPass$ ou $\iReturn{v}$ où $v$ est
    une valeur;
\item pour les valeurs gauches un chemin $φ$.
\end{itemize}

On considère en fait la clôture transitive de cette relation. Cela revient à
ajouter une règle :

\begin{mathpar}
    \irule{Trans}
        { Ξ_1 → Ξ_2
       \\ Ξ_2 → Ξ_3
        }
        { Ξ_1 → Ξ_3}
\end{mathpar}

\section{Opérations sur les valeurs}
\label{sec:sem-ops}

Un certain nombre d'opérations est possible sur les valeurs
(figure~\ref{fig:stx-ops}):

\begin{itemize}
\item
  les opérations arithmétiques $+$, $-$, $×$, $/$ et $\%$ sur les entiers.
  L'opérateur $\%$ correspond au modulo (reste de la division euclidienne).
  En cas de division par zéro, l'erreur $\serr{div}$ est levée.
\item
  les versions \enquote{pointées}
  $\floatop{+}$,
  $\floatop{-}$,
  $\floatop{×}$
  et
  $\floatop{/}$
  sur les flottants.
\item
  les opérations d'arithmétique de pointeur $+_p$ et $-_p$ qui à un chemin
  mémoire et un entier associent un chemin mémoire.
\item
  les opérations d'égalité $=$ et $≠$. L'égalité entre entiers ou entre
  flottants est immédiate. Deux valeurs composées (tableaux ou structures) sont
  égales si elles ont la même forme (même taille pour les tableaux, ou même
  champs pour les structures) et que toutes leurs sous-valeurs sont égales deux
  à deux. Deux références mémoire sont égales \linebreak lorsque les chemins
  qu'elles décrivent sont syntaxiquement égaux.
\item
  les opérations de comparaison $≤,≥,<,>$ sont définies avec leur sémantique
  habituelle sur les entiers et les flottants. Sur les références mémoires,
  elles sont définies dans le cas où les deux opérandes sont de la forme
  $φ[\cdot]$ par: $φ[n]~\opbin~φ[m] \eqdef n~\opbin~m$. Dans les autres cas,
  l'erreur $\serr{ptr}$ est renvoyée. Notamment, il n'est pas possible de
  comparer deux fonctions, deux tableaux ou deux structures.
  \added{Notamment...}
\item
  les opérateurs bit à bit sont définis sur les entiers. $\&$, $|$ et $\opxor$
  représentent respectivement la conjonction, la disjonction et la disjonction
  exclusive (XOR).
\item
  des versions logiques de la conjonction ($\&\&$) et de la disjonction ($||$)
  sont également présentes. Leur sémantique est donnée par le tableau suivant:

  \begin{center}
    \begin{tabular}{rr@{\hskip 1cm}cc}
      \toprule
       $n$ &  $m$ & $n~\&\&~m$ & $n~||~m$ \\
      \midrule
       $0$ &  $0$ & $0$        & $0$      \\
       $0$ & $≠0$ & $0$        & $1$      \\
      $≠0$ &  $0$ & $0$        & $1$      \\
      $≠0$ & $≠0$ & $1$        & $1$      \\
      \bottomrule
    \end{tabular}
  \end{center}

\item
  des opérateurs de décalage à gauche ($\ll$) et à droite ($\gg$) sont présents.
  Eux aussi ne s'appliquent qu'aux entiers.

\item
  les opérateurs arithmétiques unaires $+$, $-$, $\floatop{+}$ et $\floatop{-}$
  sont équivalents par définition à l'opération binaire correspondante.
  Par exemple $-4 \eqdef 0-4$.

\item
  $\sim$ inverse tous les bits de son opérande.
  $!$ est une version logique, % chktex 40
  c'est-à-dire que $!0 = 1$ et si $n ≠ 0$, $!n = 0$.
\end{itemize}

Si ces opérateurs sémantiques reçoivent des données incompatibles (par exemple
si on tente d'ajouter une fonction et un entier), l'erreur spéciale $\serr{typ}$
est renvoyée.

\section{Opérations sur les états mémoire}
\label{sec:memops}

\begin{definition}[Recherche de variable]

  La recherche de variable permet d'associer à une variable $x$ une adresse $a$.

  Chaque fonction peut accéder aux variables locales de la fonction en cours,
  ainsi qu'aux variables globales.

  Remarque: le cadre de variables locales le plus récent a toujours l'indice $1$.

\[
  \cLookup{(s, g)}{x} =
        \begin{cases}
            (|s|, x) & \mbox{ si } |s| > 0
                       \mbox{ et } (x ↦ v) ∈ s_1 \\
            x        & \mbox{ si } (x ↦ v) ∈ g \\
            \serr{var} & \mbox{ sinon }
        \end{cases}
\]

\end{definition}

En entrant dans une fonction, on rajoutera un cadre de pile qui contient les
paramètres de la fonction ainsi que ses variables locales. En retournant à
l'appelant, il faudra supprimer ce cadre de pile.

\begin{definition}[Manipulations de pile]

  On définit l'empilement d'un cadre de pile $c = ((x_1↦ v_1), …, (x_n↦ v_n))$
  sur un état mémoire $m = (s, g)$ (figure~\ref{fig:stackops-push}):

  \[
      \cPush{(s, g)}{c} = (c::s, g)
  \]

  On définit aussi l'extension du dernier cadre de pile, qui sert aux
  déclarations de variables locales (figure~\ref{fig:stackops-extend}):

  \[
    \cExtend{(c::s, g)}{x}{v} = ((x↦v :: c :: s), g)
  \]

  L'opération inverse de $\phxxx\cExtend$ sera simplement notée \enquote{$-$} :
  $m - x$, par exemple.

  De même on définit le dépilement (figure~\ref{fig:stackops-pop}):

  \[
      \cPop{(c::s, g)} = (s, g)
  \]

\end{definition}

\begin{definition}[Hauteur d'une valeur]
\label{def:hauteur-val}
Une valeur peut contenir une référence vers une variable de la pile.
La hauteur d'une valeur est l'indice du plus haut cadre qu'elle référence, ou
$-1$ sinon.

\begin{center}
\begin{minipage}{0.4\textwidth}
\begin{align*}
    \cH{\widehat{c}} &= -1 \\
    \cH{\widehat{f}} &= -1 \\
    \cH{\widehat{\&}~φ} &= \cHphi{φ} \\
    \cH{\widehat{\eStruct{l_1:v_1;…;l_n:v_n}}} &= \max_{i∈[1;n]} \cH{v_i} \\
    \cH{\widehat{\eArray{v_1,…,v_n}}} &= \max_{i∈[1;n]} \cH{v_i}
\end{align*}
\end{minipage}
\hspace{1.2em}
\mbox{où:}
\hspace{0.8em}
\begin{minipage}{0.3\textwidth}
\begin{align*}
    \cHphi{(x)} &= -1 \\
    \cHphi{(n,x)} &= n \\
    \cHphi{φ\widehat{.}l} &= \cHphi{φ} \\
    \cHphi{φ\widehat{[n]}} &= \cHphi{φ}
\end{align*}
\end{minipage}
\end{center}

\end{definition}


Les opérations $\mathrm{Extend}$ et $\mathrm{Pop}$ ne sont définies que pour une
pile non vide. Néanmoins cela ne pose pas de problème, puisque, lors de
l'exécution, la pile n'est vide que lors de l'évaluation d'expressions dans les
phrases de programme. À cet endroit, seules des expressions peuvent apparaître,
et leur évaluation ne manipule jamais la pile avec ces opérations.

\begin{figure}%{{{

  \centering

  \subbottom[Empilement]{
\label{fig:stackops-push}
  \begin{tikzpicture}
    [ stack/.style={draw,shape=rectangle,minimum height=5mm,minimum width=2cm}
    , every node/.style={font=\footnotesize}
    , faded/.style={fill=black!20}
    ]
    \path (8,0)   node [stack, faded] {}
    -- ++ (0,0.5) node [stack, faded] {}
    -- ++ (0,0.5) node [stack, faded] {}
    ;
    \path (12.4,0)  node [stack, faded] {}
    -- ++ (0,0.5) node [stack, faded] {}
    -- ++ (0,0.5) node [stack, faded] {}
    -- ++ (0,0.5) node [stack] {$x↦0$}
    ;
    \draw [->] (9.3, 1) -- node[auto] {$\cPush{\cdot}{(x ↦ 0)}$} ++(1.7, 0);
  \end{tikzpicture}
  }

  \subbottom[Extension de cadre]{
\label{fig:stackops-extend}
  \begin{tikzpicture}
    [ stack/.style={draw,shape=rectangle,minimum height=5mm,minimum width=2cm}
    , every node/.style={font=\footnotesize}
    , faded/.style={fill=black!20}
    ]
    \path (4,-3)  node [stack, faded] {}
    -- ++ (0,0.5) node [stack, faded] {}
    -- ++ (0,0.5) node [stack, faded] {}
    -- ++ (0,0.5) node [stack] {$x↦0$}
    ;
    \path (8.4,-3)  node [stack, faded] {}
    -- ++ (0,0.5) node [stack, faded] {}
    -- ++ (0,0.5) node [stack, faded] {}
    -- ++ (0,0.5) node [stack] {$x↦0, y↦3$}
    ;
    \draw [->] (5.3, -2) -- node[auto] {$\cExtend{\cdot}{y}{3}$} ++(1.7, 0);
  \end{tikzpicture}
  }

  \subbottom[Dépilement]{
\label{fig:stackops-pop}
  \begin{tikzpicture}
    [ stack/.style={draw,shape=rectangle,minimum height=5mm,minimum width=2cm}
    , every node/.style={font=\footnotesize}
    , faded/.style={fill=black!20}
    ]
    \path (0,0)   node [stack, faded] {}
    -- ++ (0,0.5) node [stack, faded] {}
    -- ++ (0,0.5) node [stack, faded] {}
    -- ++ (0,0.5) node [stack] {$…$}
    ;
    \path (4.2,0)   node [stack, faded] {}
    -- ++ (0,0.5) node [stack, faded] {}
    -- ++ (0,0.5) node [stack, faded] {}
    ;
    \draw [->] (1.3, 1) -- node[auto] {Pop} ++(1.7, 0);
  \end{tikzpicture}
  }

  \caption{Opérations de pile}
\label{fig:stackops}
\end{figure}%}}}

On définit aussi une opération de nettoyage de pile, qui sera utile pour les
retours de fonction.

En effet, si une référence au dernier cadre est toujours présente après
le retour d'une fonction, cela peut casser le typage.

Par exemple, dans la figure~\ref{fig:prog-cleanup}, l'exécution de $h(~)$ donne à
$p$ la valeur $(1, x)$. Puis en arrivant dans $g$, le déréférencement de $p$ va
modifier $x$ qui va avoir la valeur $1$. $x$, variable flottante, contient donc
un entier.
Dans la ligne marquée \texttt{(*)}, on réalise donc l'addition d'un entier
(contenu dans $x$ malgré le type de la variable) et d'un flottant. Cette
opération est bien typée dans le programme mais provoquera une erreur
$\serr{typ}$ à l'exécution.

Pour empêcher cela, on instrumente donc le retour de la fonction
\texttt{f} pour que \texttt{p} soit remplacé par $\eNull$. Alors dans
\texttt{h}, le déréférencement provoquera une erreur et empêchera la violation
du typage.

\begin{SaveVerbatim}[]{LeBonProgramme}
f = Fun () {
    Decl x = 0 in
    return (&x);
}

g = Fun (p) {
    Decl x = 0.0 in
    *p = 1;
    x <- x + 2.0; // (*)
}

h = Fun () {
    Decl p = f() in
    g(p);
}
\end{SaveVerbatim}

\begin{figure}
\hspace{1cm}\BUseVerbatim{LeBonProgramme}
\caption{Cassage du typage par un pointeur fou}
\label{fig:prog-cleanup}
\end{figure}

Pour définir l'opération de nettoyage, on commence par définir une opération de
nettoyage selon un prédicat sur les chemins:

\begin{align*}
\cVP{p}{\widehat{c}} &= \widehat{c} \\
\cVP{p}{\widehat{f}} &= \widehat{f} \\
\cVP{p}{\wphi{φ}} &= \begin{cases}
                       \eNull   & \mbox{ si }    p(φ) \\
                       \wphi{φ} & \mbox{ sinon}
                     \end{cases} \\
\cVP{p}{\widehat{\eStruct{l_1:v_1, …, l_n:v_n}}} &=
    \widehat{\eStruct{l_1:\cVP{p}{v_1}, …, l_n:\cVP{p}{v_n}}} \\
\cVP{p}{\widehat{\eArray{v_1, …, v_n}}} &=
    \widehat{\eArray{\cVP{p}{v_1}, …, \cVP{p}{v_n}}}
\end{align*}

On l'étend ensuite aux cadres de pile, puis aux états mémoire:

\begin{align*}
\cLP{p}{(x_1 ↦ v_1, …, x_n ↦ v_n)} &= (x_1 ↦ \cVP{p}{v_1}, …, x_n ↦ \cVP{p}{v_n}) \\
\\
                \cP{p}{(s, g)} &= (s', g') \\
                 \mbox{où } s' &= (\cLP{p}{s_1}, …, \cLP{p}{s_{|s|}}) \\
                            g' &=  \cLP{p}{g}
\end{align*}

À l'aide de ces fonctions, on définit quatre opérations, permettant de nettoyer
des états mémoire ou des valeurs en enlevant tout un niveau de pile ou seulement
une variable:

\begin{mathpar}
\cCleanup{m} =  \cP{λφ. \cHphi{φ} > |m|}{m}

\cCleanVar{m}{a}  = \cP{λφ. φ = a}{m}

\cV{p}{v}    = \cVP{λφ. \cHphi{φ} > |m|}{v}

\cVV{v}{a}        = \cVP{λφ. φ = a}{v}
\end{mathpar}

Ces 4 fonctions seront utilisées dans plusieurs règles dans la suite de ce
chapitre.

\paragraph{Remarques}
\added{Les 3 paragraphes suivants}

Ces opérations ne sont pas toujours bien définies. Par exemple, $\phxxx\cExtend$
ne peut pas s'appliquer à une pile vide, et $m - x$ n'est défini que si une
variable $x$ existe au sommet de la pile de $m$. Ce caractère partiel ne pose
pas de problème de par la structure des règles qui vont utiliser ces
constructions. Par exemple, à chaque empilement correspond exactement un
dépilement. De plus, les phrases d'un programme ne peuvent pas faire intervenir
de déclaration de variable (une instruction est forcément dans une fonction),
donc $\phxxx\cExtend$ réussit toujours.

Un autre problème se pose si deux variables ont le même nom dans un cadre. Elles
ne peuvent pas être distinguées. On interdit donc ce cas en demandant aux
programmes d'être bien formés: au sein d'une fonction, les paramètres ainsi que
l'ensemble des locales déclarées doivent être de noms différents. En pratique,
une phase préalable d'$α$-conversion peut renommer les variables problématiques.

De plus, le fait d'ajouter cette étape de nettoyage à chaque retour de fonction
peut être assez coûteux. C'est un compromis: si on considère que les programmes
se comportent bien et ne créent pas de pointeurs fous (pointant au-dessus de la
pile), alors cette phase est inutile et peut être remplacée par l'identité.
Autrement dit, il s'agit seulement d'une technique pour s'assurer de ne pas
avoir d'erreurs dans la sémantique. L'ajout d'un ramasse-miette, ou une
vérification préalable par un système de régions~\cite{jfp92} peut garantir
qu'il n'y a pas de telles constructions dangereuses.

\section{Accesseurs}
\label{sec:mem-access}

Le but de cette section est de définir rigoureusement les accès à la mémoire. À
partir d'un état mémoire $m$ et d'une valeur gauche $φ$, on veut pouvoir obtenir:

\begin{itemize}
    \item la valeur accessible au chemin $φ$: $m[φ]_Φ$
    \item l'état mémoire obtenu en remplaçant celle-ci par une nouvelle valeur $v'$ :
            $m[φ ← v']_Φ$
\end{itemize}

Pour définir cette lentille indexée $Φ$, on commence par définir des lentilles
élémentaires, et on les compose pour pouvoir définir des lentilles entre
valeurs.

On commence par définir deux lentilles $I$ et $L$ pour accéder aux structures de
listes. $I$ accède par indice et $L$ par clef (dans une liste d'association,
donc).

Cela permet ensuite de définir $A$, qui extrait une valeur à partir d'un nom de
variable et d'une éventuelle hauteur de pile. Pour cela, on compose les
lentilles $I$ et $L$.

Les autres travaillent sur des valeurs composées, c'est-à-dire sur les
structures et tableaux. La lentille $F$ extrait une sous-valeur
correspondant au champ d'une structure. Le fonctionnement est similaire à la
lentille $L$ puisqu'on accède par nom à une sous-structure.
La lentille $T$, quant à elle, permet d'accéder au $n$\ieme élement d'un
tableau. De ce point de vue, elle est similaire à $I$ mais en travaillant sur
les valeurs.

Enfin, on définit $Φ$ pour accéder à n'importe quelle sous-valeur d'une variable
dans la mémoire. Cela utilise $A$, $F$ et $T$ précédemment définis.

La figure~\ref{fig:dep-lens} résume ces dépendances. Les lignes pleines
indiquent quelles sont les définitions utilisées, et les pointillés relient les
lentilles similaires. À droite, on donne un exemple des lentilles de base. La
valeur entourée correspond au \enquote{curseur} de la lentille, c'est-à-dire la
valeur qui peut être renvoyée ou mise à jour.

\begin{figure}[h]
\centering
\vspace{5mm}
\begin{minipage}[b]{0.4\textwidth}
\begin{tikzpicture}
\node[draw] (I) {I};
\node[draw, below of=I] (L) {L};
\node[draw, below of=L] (F) {F};
\node[draw, below of=F] (T) {T};
\node[draw, right of=L, node distance=2cm, yshift=5mm] (A) {A};
\node[draw, right of=F, node distance=4cm, yshift=-5mm] (P) {$Φ$};

\draw[->] (I) -- (A.north west);
\draw[->] (L) -- (A.south west);
\draw[->] (A) -- (P.north west);
\draw[->] (F) -- (P.west);
\draw[->] (T) -- (P.south west);

\draw[dotted] (L) -- ++(-5mm,  0) |- (F);
\draw[dotted] (I) -- ++(-10mm, 0) |- (T);
\end{tikzpicture}
\end{minipage}
\begin{minipage}[b]{0.5\textwidth}
\begin{align*}
I(3)     &: (3, 14, \circled{15}, 92, 65) \\
L(toto)  &: ((toto, \circled{3}), (tata, 6), (titi, 2)) \\
F(y)     &: \{ x: 0; y: \circled{-3} \} \\
T(0)     &: \eArray{\circled{1}, 2, 3, 5} \\
\end{align*}
\end{minipage}
\vspace{5mm}

\caption{Dépendances entre les lentilles}
\label{fig:dep-lens}

\end{figure}

\subsection*{Accès à une liste par indice : $I$}

\changed{$\textrm{List}(\alpha) \rightarrow \sList{\alpha}$}

On définit une lentille indexée $I : ℕ → \setLens{\sList{α}}{α}$
permettant d'accéder aux éléments d'une liste par leur indice. On rappelle que
les listes sont des suites finies, définies page~\pageref{page:def-listes}. En
outre, I n'est définie que pour $n ∈ [1 ;|l|]$.

\begin{align*}
    l{[n]}_I     &= l_n \mbox{ si } n ∈ [1; |l|] \\
    l{[n ← x]}_I &= l' \\
           \mbox{ où } & l'_n = x \\
                       & ∀ i ≠ n, l'_i = l_i
\end{align*}

\subsection*{Accès à une liste d'associations : $L$}

  Une liste d'association est une liste de paires (clef, valeur) avec
  l'invariant supplémentaire que les clefs sont uniques. Il est donc possible de
  trouver au plus une valeur associée à une clef donnée. L'écriture est
  également possible, en remplaçant un couple par un couple avec une valeur
  différente.

\begin{align*}
l{[x]}_L  &=
    \begin{cases}
        v          & \mbox{ si } ∃! n ∈ [1;|l|], ∃v, l_n = (x ↦ v) \\
        \serr{var} & \mbox{ sinon }
    \end{cases} \\
l{[x ← v]}_L  &=
    \begin{cases}
        l{[n ← (x ↦ v')]}_I & \mbox{ si } ∃! n ∈ [1;|l|], ∃v, l_n = (x ↦ v) \\
        \serr{var} & \mbox{ sinon }
    \end{cases} \\
\end{align*}

\subsection*{Accès par adresse : $A$}

Les états mémoire sont constitués des listes d'association (nom, valeur).

L'accesseur par adresse ${[\cdot]}_A$ permet de généraliser l'accès à ces
valeurs en utilisant comme clef non pas un nom mais une adresse.

Selon cette adresse, on accède soit à la liste des globales, soit à une
des listes de la pile des locales.

On pose $m = (s, g)$.

Les accès aux globales se font de la manière suivante. Si la variable n'existe
pas, notons que $L$ retourne $\serr{var}$.

\[
A((x)) = \mathrm{Snd} \ggg L(x)
\]

$\mathrm{Snd}$ désigne la lentille entre un couple et sa deuxième composante.
Ainsi, par exemple $m{[(x) ← v]}_A = (s, g[x←v]_L)$.

%\begin{align*}
%m {[x]}_A     &= g{[x]}_L \\
%m {[x ← v]}_A &= (s, g{[x←v]}_L) \\
%\end{align*}

Les accès aux locales reviennent à accéder à la bonne variable du bon cadre de
pile. Cela revient naturellement à composer les lentilles $L$ et $I$. On définit
donc une lentille $ℒ_{n,x} = I(|s|-n+1) \ggg L(x)$ qui accède à la variable $x$
du $n$\ieme cadre de pile.

\begin{align*}
m {[(n, x)]}_A     & = \begin{cases}
                           \lensGet{ℒ_{n,x}}{s} & \mbox{ si $n ∈ [1 ; |s|]$} \\
                           \serr{var}          & \mbox{ sinon}
                       \end{cases} \\
m {[(n, x) ← v]}_A & = \begin{cases}
                            (\lensPut{ℒ_{n,x}}{v}{s}, g) & \mbox{ si $n ∈ [1;|s|]$} \\
                            \serr{var}          & \mbox{ sinon}
                       \end{cases} \\
\end{align*}

Les numéros de cadre qui permettent d'identifier les locales (le $n$ dans $(n,
x)$) croissent avec la pile. D'autre part, l'empilement se fait en tête de liste
(près de l'indice 1). Donc pour accéder aux plus vieilles locales (numérotées
1), il faut accéder au dernier élément de la liste. Ceci explique pourquoi un
indice $|s|-n+1$ apparaît dans la définition précédente.

\subsection*{Accès par champ : $F$}

  Les valeurs qui sont des structures possèdent des sous-valeurs, associées à
  des noms de champ.

  L'accesseur ${[ \cdot ]}_F$ permet de lire et de modifier un champ de ces
  valeurs.

  L'erreur $\serr{field}$ est levée si on accède à un champ non existant.

  \begin{align*}
    \eStruct{ l_1 : v_1; … ; l_n : v_n }{[l]}_F &= v_i \mbox{ si } ∃i∈[1;n], l=l_i\\
    \eStruct{ l_1 : v_1; … ; l_n : v_n }{[l]}_F     &= \serr{field} \mbox{ sinon} \\
    \eStruct{ l_1 : v_1; … ; l_n : v_n }{[l ← v]}_F &=
        \eStruct{ l_1 : v_1 \\
           &\quad ; … \\
           &\quad ; l_{p-1} : v_{p-1} \\
           &\quad ; l_p : v \\
           &\quad ; l_{p+1} : v_{p+1} \\
           &\quad ; … \\
           &\quad ; l_n : v_n } \mbox{ si } ∃p ∈ [1;n], l = l_p \\
    \eStruct{ l_1 : v_1; … ; l_n : v_n }{[l ← v]}_F &= \serr{field} \mbox{ sinon}
  \end{align*}

\subsection*{Accès par indice de tableau : $T$}
\label{subsec:lens-array}

  On définit de même un accesseur ${[\cdot]}_T$ pour les accès par indice à des
  valeurs tableaux. Néanmoins le paramètre indice est toujours un entier et pas
  une expression arbitraire. Notons que les accès sont vérifiés dynamiquement:
  il ne peut pas y avoir de débordement de tableau.

  \begin{align*}
    \eArray{ v_1 ; … ; v_n } {[i]}_T   &= v_{i+1} \textrm{ si } i ∈ [0;n-1] \\
    \eArray{ v_1 ; … ; v_n } {[i]}_T   &= \serr{array} \textrm{ sinon} \\
    \eArray{ v_1 ; … ; v_n } {[i←v]}_T &= \eArray{ v'_1 ; … ; v'_n } \textrm{ si } i ∈ [0;n-1] \\
                      \mbox{ où } & \begin{cases}
                                      v'_i = v \\
                                      ∀j≠i, v'_j = v_j \\
                                    \end{cases} \\
    \eArray{ v_1 ; … ; v_n } {[i←v]}_T &= \serr{array} \textrm{ sinon} \\
  \end{align*}

\subsection*{Accès par chemin : $Φ$}
\label{subsec:acces-phi}

  L'accès par chemin $Φ$ permet de lire et de modifier la mémoire en profondeur.

  On peut accéder directement à une variable, et les accès à des sous-valeurs se
  font en composant les accesseurs (définition~\ref{def:lens-comp},
  page~\pageref{def:lens-comp}):

  \begin{align*}
    Φ(a)    &= A(a) \\
    Φ(φ.l)  &= Φ(φ) \ggg F(l) \\
    Φ(φ[i]) &= Φ(φ) \ggg T(i)
  \end{align*}

\paragraph{Remarque}

\added{Ce paragraphe}

Dans toute la suite, lorsque ce n'est pas ambigü, on emploiera la notation
$m[φ]$ pour désigner $m[φ]_Φ$. Il est important de remarquer que $m$ désigne un
état particulier et $φ$ un chemin particulier, mais que $Φ$ est la lentille
indexée globale définie page~\pageref{subsec:acces-phi}.

\section{Contextes d'évaluation}

L'évaluation des expressions repose sur la notion de contextes d'évaluation.
L'idée est que, si on peut évaluer une expression, alors on peut évaluer une
expression qui contient celle-ci.

Par exemple, supposons que $\mm{m}{f(3)}{m}{2}$. Alors on peut ajouter la
constante $1$ à gauche de chaque expression sans changer le résultat:
$\mm{m}{1+f(3)}{m}{1+2}$. On a utilisé le même contexte
$C = 1+\ctxEmpty$.

Pour pouvoir raisonner en termes de contextes, 3 points sont nécessaires:

\begin{itemize}
\item comment découper une expression selon un contexte;
\item comment appliquer une règle d'évaluation sous un contexte;
\item comment regrouper une expression et un contexte.
\end{itemize}

\begin{figure}
\figctx{}

\caption{Contextes d'évaluation}
\label{fig:eval-ctx}
\end{figure}

Le premier point consiste à définir les contextes eux-mêmes
(figure~\ref{fig:eval-ctx}).

Dans cette définition, chaque cas hormis le cas de base fait apparaître
exactement un \enquote{$C$}. Chaque contexte est donc constitué d'exactement une
occurrence de $\ctxEmpty$ (une dérivation de $C$ est toujours linéaire).
L'opération de substitution consiste à remplacer ce trou: $\ctxSub{C}{X}$ est
l'objet syntaxique (instruction, expression ou valeur gauche) obtenu en
remplaçant l'unique $\ctxEmpty$ dans $C$ par $X$. Par exemple,
$\ctxSub{\iDecl{x}{2+\ctxEmpty}{\iPass}}{5}$ est $\iDecl{x}{2+5}{\iPass}$

À titre d'illustration, décomposons l'évaluation de $e_1~\opbin~e_2$ en $v =
v_1~\widehat{\opbin}~v_2$ depuis un état mémoire $m$:

\begin{enumerate}
\item
  on commence par évaluer l'expression $e_1$ en une valeur $v_1$. Le nouvel
  état mémoire est noté $m'$. Soit donc $\mm{m}{e_1}{m'}{v_1}$.
\item
  En appliquant la règle \textsc{Ctx} (définie page suivante) avec $C =
  \ctxOp{\ctxEmpty}{e_2}$ (qui est une des formes possibles pour un contexte
  d'évaluation), on déduit de 1.\ que
  $\mm{m}{e_1~\opbin~e_2}{m'}{v_1~\opbin~e_2}$
\item
  D'autre part, on évalue $e_2$ depuis $m'$. En supposant encore que
  l'évaluation converge, notons $v_2$ la valeur calculée et $m''$ l'état mémoire
  résultant: $\mm{m'}{e_2}{m''}{v_2}$.
\item
  Appliquons la règle~\textsc{Ctx} à 3.\ avec $C = \ctxOp{v_1}{\ctxEmpty}$. On
  obtient $\mm{m}{v_1~\opbin~e_2}{m'}{v_1~\opbin~v_2}$.
\item
  En combinant les résultats de 2.\ et 4.\ on en déduit que
  $\mm{m}{e_1~\opbin~e_2}{m''}{v_1\opbin~v_2}$.
\item D'après la règle~\textsc{Exp-Binop}
    (page~\pageref{rule:exp-binop}),
  $ \mm{m''}{v_1\opbin~v_2}{m''}{v_1~\widehat{\opbin}~v_2}$
\item D'après 5.\ et 6., on a par combinaison
  $\mm{m}{e_1~\opbin~e_2}{m''}{v}$
  en posant
  $v = v_1~\widehat{\opbin}~v_2$.
\end{enumerate}

\changed{Ce cas était défini plus haut}

Le deuxième point sera résolu par la règle d'inférence suivante.

\begin{mathpar}
  \semrule{Ctx}
\end{mathpar}

Enfin, le troisième revient à définir l'opérateur de substitution
$\phxx{\ctxSub}$ présent dans la règle précédente. Notons que puisque $i
\gramisa e$ et $e \gramisa lv$, on peut aussi l'appliquer aux expressions et aux
valeurs gauches : l'opération $\phxx\ctxSub$ est purement syntaxique.

\section{Valeurs gauches}
\label{sec:eval-lv}

Obtenir un chemin à partir d'un nom de variable revient à résoudre le nom de
cette variable: est-elle accessible? Le nom désigne-t-il une variable locale
ou une variable globale?

\begin{mathpar}
  \semrule{Phi-Var}
\end{mathpar}

Les règles portant sur le déréférencement et l'accès à un champ de structure
sont similaires: on commence par évaluer la valeur gauche sur laquelle porte ce
modificateur, et on place le même modificateur sur le chemin résultant. Dans le
cas des champs de structure, l'annotation de structure $S$ n'est pas prise
en compte pour l'évaluation: elle servira uniquement au typage.

\begin{mathpar}
  \semrule{Phi-Struct}
\end{mathpar}

Enfin, pour évaluer un chemin dans un tableau, on commence par procéder comme
précédemment, c'est-à-dire en évaluant la valeur gauche sur laquelle porte
l'opération d'indexation. Puis on évalue l'expression d'indice en une valeur qui
permet de construire le chemin résultant.

\begin{mathpar}
  \semrule{Phi-Array}
\end{mathpar}

Notons qu'en procédant ainsi, on évalue les valeurs gauches en allant de gauche
à droite: dans l'expression $x[e_1][e_2][e_3]$, $e_1$ est évalué en premier,
puis $e_2$, puis $e_3$.

La règle portant sur le déréférencement est particulière. On peut penser que la
bonne définition de $φ$ consiste à se calquer sur la définition de $lv$, en
remplaçant les noms de variable par leur adresse résolue et en évaluant les
indices de tableau, et à ajouter une règle qui transforme $*φ$ en
$\widehat{*}φ$.

Cela ne fonctionne pas, car alors les déréférencements sont évalués trop tard:
au moment de l'affectation dans la valeur gauche plutôt qu'à sa définition. La
figure~\ref{fig:lazy-deref} illustre ce problème.

%{{{
\begin{SaveVerbatim}[]{LeLazy}
Decl s0 = { .f : 0 } in
Decl s1 = { .f : 1 } in
Decl x  = & s0 in
Decl p = & ((*x).f) in
/* (a) */
x <- & s1
/* (b) */
\end{SaveVerbatim}

\begin{figure}[h]

\hspace{1cm}\BUseVerbatim{LeLazy}

\caption{Évaluation stricte ou paresseuse des valeurs gauches}
\label{fig:lazy-deref}
\end{figure}%}}}

On s'intéresse à l'évaluation
de l'expression \texttt{*p} aux points \texttt{(a)} et \texttt{(b)}. Avec une
sémantique paresseuse (en ajoutant un $\widehat{*}φ$), la valeur de
\texttt{p}
est $\widehat{\&}~((*(1,x)).f)$, donc
\texttt{*p} est évalué à $0$ en
\texttt{(a)}
et $1$ en
\texttt{(b)}.
Au contraire, avec une sémantique stricte (correcte),
\texttt{p} vaut
$\widehat{\&}~(((1,s0).f)$ et donc
\texttt{*p} est évalué à $0$ en
\texttt{(a)}
et en
\texttt{(b)}.

Dans le cas où la valeur référencée n'a pas la forme $\widehat{\&} φ$ ou
$\widehat{\eNull}$, aucune règle ne peut s'appliquer (comme lorsqu'on cherche à
réduire l'addition d'une fonction et d'un entier, par exemple). Cela est
préférable à renvoyer $\serr{ptr}$ car on montrera que ce cas est toujours évité
dans les programmes typés (théorème~\ref{thm:progres}).

\begin{mathpar}
  \semrule{Phi-Deref}

  \semrule{Phi-Deref-Null}
\end{mathpar}

Par exemple, $lv = x.l_S[2*n].g_T$ pourra s'évaluer en $φ = (2, x).l[4].g$.

\section{Expressions}
\label{sec:eval-exp}

Évaluer une constante est le cas le plus simple, puisqu'en quelque sorte
celle-ci est déjà évaluée. À chaque constante syntaxique $c$, on peut associer
une valeur sémantique $\widehat{c}$. Par exemple, au chiffre (symbole) $3$, on
associe le nombre (entier) $\widehat{3}$.

\begin{mathpar}
  \semrule{Exp-Cst}
\end{mathpar}

De même, une fonction est déjà évaluée:

\begin{mathpar}
  \semrule{Exp-Fun}
\end{mathpar}

Pour lire le contenu d'un emplacement mémoire (valeur gauche), il faut tout d'abord
l'évaluer en un chemin.

\begin{mathpar}
  \semrule{Exp-Lv}
\end{mathpar}

Pour évaluer une expression constituée d'un opérateur, on évalue une
sous-expression, puis l'autre (l'ordre d'évaluation, est encore imposé: de
gauche à droite). À chaque opérateur $\opbin$, correspond un opérateur
sémantique $\widehat{\opbin}$ qui agit sur les valeurs. Par exemple, l'opérateur
$\widehat{+}$ est l'addition entre entiers machine
(page~\pageref{page:entiers-bits}). Comme précisé dans la
section~\ref{sec:sem-ops}, la division par zéro via $/$, $\%$ ou $\floatop{/}$
provoque l'erreur $\serr{div}$.

\label{rule:exp-binop}
\begin{mathpar}
  \semrule{Exp-UnOp}

  \semrule{Exp-BinOp}
\end{mathpar}

Il est nécessaire de dire un mot sur les opérations $\widehat{+_p}$ et
$\widehat{-_p}$ définissant l'arithmétique des pointeurs. Celles-ci sont
uniquement définies pour les références mémoire à un tableau, c'est-à-dire
celles qui ont la forme $\widehat{\&}~φ[n]$. On a alors:

\begin{align*}
    \widehat{\&}~φ[n]~\widehat{+}_p~i & = \widehat{\&}~φ[n~\widehat{+}~i] \\
  \widehat{\&}~φ[n]~]~\widehat{-}_p~i & = \widehat{\&}~φ[n~\widehat{-}~i] \\
\end{align*}

Cela implique qu'on ne peut pas faire faire d'arithmétique de pointeurs au sein
d'une même structure. Autrement c'est une erreur de manipulation de
pointeurs\footnote{ Cela est cohérent avec la norme C99:
\enquote{If the pointer operand points to an element of an array object, and the
    array is large enough, […] ; otherwise, the behavior is undefined.
    }~\cite[6.5.6~§8]{AnsiC}}
et l'opérateur $\widehat{\opbin}$ renvoie $\serr{ptr}$.\label{page:def-arith-ptr-error}

\added{Cas où on sort de l'espace alloué}

Si l'indice calculé ($n~\widehat{+}~i$ ou $n~\widehat{-}~i$) sort de l'espace
alloué, alors l'erreur sera faite au moment de l'accès: la lentille $T$ renverra
$\serr{array}$ (page~\pageref{subsec:lens-array}).

Une left-value s'évalue en le chemin correspondant.

\begin{mathpar}
  \semrule{Exp-Addr}
\end{mathpar}

L'affectation se déroule en 3 étapes. D'abord, l'expression est évaluée en une
valeur $v$. Ensuite, la valeur gauche est évaluée en un chemin $φ$. Enfin, un
nouvel état mémoire est construit, où la valeur accessible par $φ$ est remplacée
par $v$. Comme dans le langage C, l'expression d'affectation produit une valeur,
qui est celle qui a été affectée.

\begin{mathpar}
  \semrule{Exp-Set}
\end{mathpar}

\subsection*{Expressions composées}

Les littéraux de structures sont évalués en leurs constructions syntaxiques
respectives. Puisque les contextes d'évaluation sont de la forme $ [ v_1 ; … ; C
; … ; e_n ] $, l'évaluation se fait toujours de gauche à droite.

\begin{mathpar}
  \semrule{Exp-Struct}

  \semrule{Exp-Array}
\end{mathpar}

%\shorthandoff{!}
%\begin{figure}%{{{

  %\centering

  %\begin{tikzpicture}
    %[node distance=2cm]

    %\node (m0) {$m_0$};
    %\node[right of=m0, node distance=3cm] (m1) {$m_1$};
    %\node[below of=m1]                    (m2) {$m_2$};
    %\node at ($(m0)!(m2)!(0,1)$)          (m3) {$m_3$};
    %\node[below of=m3]                    (m4) {$m_4$};

    %\draw [->] (m0) -- node[auto] {$\mathrm{Push}(\vec{a}↦\vec{v})$}    (m1);
    %\draw [->] (m1) -- node[auto] {$i → \iReturn{v_r}$}                 (m2);
    %\draw [->] (m2) -- node[auto] {$\mathrm{Pop}$}                      (m3);
    %\draw [->] (m3) -- node[auto] {$\mathrm{Cleanup}$}                  (m4);

    %\draw [dashed, ->] ($ (m0) + (-5mm, -3mm) $)
                    %-- node[auto, swap] {$f(\vec{e}) → v$}
                       %($ (m4) + (-5mm, 3mm) $);

  %\end{tikzpicture}

  %\caption[Appel d'une fonction]
  %{ Appel d'une fonction.
    %La taille de la pile croît de gauche à droite,
    %et les réductions se font de haut en bas.
  %}
%\label{fig:fcall-details}

%\end{figure}%}}}
%\shorthandon{!}

L'appel de fonction est traité de la manière suivante. On ne peut pas facilement
relier un pas d'évaluation de $i$ à un pas d'évaluation de $\eFun{a}{i}
(v_1,…,v_n)$, et donc un contexte $C \gramisa \eFun{a}{\ctxEmpty} (v_1,…,v_n)$
n'est pas à considérer. En effet, l'empilement suivi du dépilement modifie la
mémoire.

On emploie donc une règle \textsc{Exp-Call-Ctx} qui relie un pas interne
$\mm{m_1}{i}{m_2}{i'}$ à un pas externe. Une fois l'instruction interne réduite
d'un pas, on évalue les arguments en des valeurs $v'_i$. Ils correspondent aux
nouvelles valeurs à passer à la fonction.

Les autres fonctions permettent de transférer le flot de contrôle: en retournant
la même instruction pour une instruction terminale, ou en propageant une erreur.
Dans le cas où on retourne de la fonction par $i = \iReturn{v}$, il faut alors
supprimer les références aux variables qui ont disparu grâce aux opérateurs
$\phx\cCleanup$ et $\phxx\cV$.

\label{page:return-fonction}
On suppose deux choses sur chaque fonction: d'une part, les noms de ses
arguments sont deux à deux différents et, d'autre part, son corps se termine par
une instruction $\iReturn{\cdot}$. Cela veut dire que la dernière instruction
doit être soit de cette forme, soit par exemple une alternative dans laquelle
les deux branches se terminent par un $\iReturn{\cdot}$. C'est une propriété qui
peut être détectée statiquement avant l'exécution. Néanmoins, dans la syntaxe
concrète, on peut supposer qu'un $\iReturn{\eUnit}$ est inséré automatiquement
en fin de fonction lorsqu'aucun $\iReturn{\cdot}$ n'est présent dans son corps.

\begin{mathpar}
    \semrule{Exp-Call-Ctx}

    \semrule{Exp-Call-Err}

    \semrule{Exp-Call-Return}
\end{mathpar}

\section{Instructions}
\label{sec:eval-instr}

Les cas de la séquence et de l'évaluation d'une expression sont sans surprise.

\begin{mathpar}
  \semrule{Seq}

  \semrule{Exp}
\end{mathpar}

L'évaluation de $\iDecl{x}{v}{i}$ sous $m$ se fait de la manière suivante,
similaire à l'appel de fonction.
La règle principale est \textsc{Decl-Ctx} qui relie un pas d'évaluation sous une
déclaration à un pas d'évaluation externe : pour ce faire, on étend l'état
mémoire en ajoutant $x$, on effectue le pas, puis on enlève $x$. L'instruction
résultante est la déclaration de $x$ avec la nouvelle valeur $v'$ de $x$ après
le pas d'exécution\footnote{
    On peut remarquer qu'il est impossible de définir un contexte d'évaluation
    $C \gramisa \iDecl{x}{v}{C}$. En effet, puisque celui-ci nécessiterait
    d'ajouter une variable, il ne préserve pas la mémoire.
}.

\label{page:decl-masquage}
On suppose qu'il n'y a pas de masquage au sein d'une fonction, c'est-à-dire que
le nom d'une variable déclarée n'est jamais dans l'environnement avant cette
déclaration.

Si $i$ est terminale ($\iPass$ ou $\iReturn{v}$), alors on peut s'évaluer en $i$
en nettoyant l'espace mémoire des références à $x$ qui peuvent subsister.

Enfin, si une erreur se produit elle est propagée.

\begin{mathpar}
 \semrule{Decl-Pass}

 \semrule{Decl-Return}

 \semrule{Decl-Ctx}

 \semrule{Decl-Err}
\end{mathpar}

Pour traiter l'alternative, on a besoin de 2 règles. Elles commencent de la même
manière, en évaluant la condition. Si le résultat est 0 (et seulement dans ce
cas), c'est la règle \textsc{If-False} qui est appliquée et l'instruction
revient à évaluer la branche \enquote{\emph{else}}. Dans les autres cas, c'est
la règle \textsc{If-True} qui s'applique et la branche \enquote{\emph{then}} qui
est prise.

\begin{mathpar}
  \semrule{If-False}

  \semrule{If-True}
\end{mathpar}

On exprime la sémantique de la boucle comme une simple règle de réécriture:

\begin{mathpar}
  \semrule{While}
\end{mathpar}

Enfin, si un $\phx\iReturn$ apparaît dans une séquence, on peut supprimer la suite:

\begin{mathpar}
  \semrule{Return}
\end{mathpar}

\section{Erreurs}
\label{sec:eval-errors}

Les erreurs se propagent des données vers l'interprète; c'est-à-dire que si
une expression ou instruction est réduite en une valeur d'erreur $Ω$, alors une
transition est faite vers cet état d'erreur.

Cela est aussi vrai d'une sous-expression ou sous-instruction: si l'évaluation
de $e_1$ provoque une erreur, l'évaluation de $e_1 + e_2$ également. La notion
de sous-expression ou sous instruction est définie en fonction des contextes
$C$. Notons que, dans \textsc{Eval-Err}, $\ctxSub{C}{e}$ peut être une expression
ou une instruction.

\begin{mathpar}
  \semrule{Exp-Err}

  \semrule{Eval-Err}
\end{mathpar}

\section{Phrases et exécution d'un programme}

Un programme est constitué d'une suite de phrases: déclarations de fonctions,
de variables et de types, et évaluation d'expressions.

Donc l'évaluation d'une phrase $p$ fait passer d'un état mémoire $m$
à un autre $m'$, ce que l'on note $\ph{m}{p}{m'}$.

\changed{Initialiseur $\rightarrow$ valeur initiale}

L'évaluation d'une expression est uniquement faite pour ses effets de bord. Par
exemple, après avoir défini les fonctions du programme, on pourra appeler
\texttt{main()}. La déclaration d'une variable globale, quant à elle, consiste à
évaluer sa valeur initiale et à étendre l'état mémoire avec ce couple (variable,
valeur). On suppose que les variables globales ont toutes des noms différents.
Notons que ces évaluations se font à grands pas.

Enfin, l'exécution d'un programme est sans surprise l'exécution de ses phrases,
les unes à la suite des autres.

\begin{mathpar}
  \semrule{ET-Exp}

  \semrule{ET-Var}

  \semrule{Prog}

\end{mathpar}

\section{Exemple}

Considérons le programme suivant:

\begin{Verbatim}
x = 0
f = fun(q) {
    *q <- 1
}
f(&x)
\end{Verbatim}

Ce programme est constitué des phrases $p_1, p_2$ et $p_3$. On rappelle que par
rapport à la syntaxe concrète, un \iReturn{\eUnit} est inséré automatiquement,
donc $p_2 = f=\eFun{q}{*q←1;\iReturn{\eUnit}}$

D'après \textsc{Prog}, l'évaluer va revenir à évaluer à la suite ces 4 phrases.

\paragraph{Déclaration de \texttt{x}}

D'après \textsc{Prog}, on part d'un état mémoire $m_0 = ([~], [~])$. Pour
trouver $m_1$ tel que $\ph{m_0}{x = 0}{m_1}$, il faut appliquer la règle
\textsc{ET-Var}.

Celle-ci va étendre l'ensemble (vide) des globales mais demande d'évaluer
l'expression $0$. D'après \textsc{Exp-Cst}, $\mm{m_0}{0}{m_0}{\widehat{0}}$.

Donc $\ph{m_0}{x = 0}{m_1}$ en posant $m_1 = ([~], [x↦\widehat{0}])$.

\paragraph{Déclaration de \texttt{f}}

On se trouve encore dans le cas de la déclaration d'une variable globale. Il
faut comme auparavant évaluer l'expression. C'est la règle \textsc{Exp-Fun} qui
s'applique: $\mm{m_1}{\eFun{q}{*q←1;\iReturn{\eUnit}}}
{m_1}{\widehat{\eFun{q}{*q←1;\iReturn{\eUnit}}}}$ (ce qui revient à dire que la
fonction est directement placée en mémoire).

Ainsi $\ph{m_1}{f = \eFun{q}{*q←1;\iReturn{\eUnit}}}{m_2}$ où:

\[ m_2 = ([~], [f↦\widehat{\eFun{q}{*q←1;\iReturn{\eUnit}}};x↦\widehat{0}])\]

\paragraph{Appel de \texttt{f}}

Ici, on évalue une expression pour ses effets de bords. La règle à appliquer est
\textsc{ET-Exp}, qui a comme prémisse $\mm{m_2}{f(\&x)}{m_3}{v}$.

D'après la forme de l'expression, la règle à appliquer va être
\textsc{Exp-Call-Return}. Mais il va falloir d'abord réécrire l'expression à
l'aide de \textsc{Ctx} (pour que l'expression appellée ait la forme
$\widehat{f}$ et l'argument soit évalué) et \textsc{Exp-Call-Ctx} (pour que le
corps de la fonction ait pour forme $\iReturn{\eUnit}$).

Tout d'abord on applique donc \textsc{Ctx} avec $C = \ctxEmpty (\&x)$. Comme on
a via \textsc{Phi-Var} puis \textsc{Exp-Lv}:

\[
    \mm{m_2}{f}{m_2}{\widehat{\eFun{q}{*q←1;\iReturn{\eUnit}}}}
\]

On en déduit que:

\[
    \mm{m_2}{f(\&x)}{m_2}{\widehat{\eFun{q}{*q←1;\iReturn{\eUnit}}}(\&x)}
\]

Une application supplémentaire de \textsc{Ctx} permet d'en arriver à la ligne
suivante (en évaluant $\&x$ en $\widehat{\&}(x)$ via \textsc{Exp-Addr}):

\[
    \mm{m_2}{f(\&x)}{m_2}{\widehat{\eFun{q}{*q←1;\iReturn{\eUnit}}}(\widehat{\&}(x))}
\]

La fonction et son argument sont évalués, donc on peut appliquer
\textsc{Exp-Call-Ctx}. En posant $m_2' = \cPush{m_2}{(q ↦ \widehat{\&}(x))}$, le
but est de trouver $m_2''$ et $v$ tels que:

\[
    \mm{m_2'}{*q←1;\iReturn{\eUnit}}{m_2''}{\iReturn{v}}
\]

Puisque l'instruction est une séquence, on va appliquer \textsc{Seq}. La
première partie n'étant pas $\iPass$, il faut l'évaluer grâce à la règle
\textsc{Ctx} avec $C = \ctxEmpty;\iReturn{\eUnit}$.

Le nouveau but est de trouver un $m_2''$ tel que
$\mm{m_2'}{*q←1}{m_2'}{\iPass}$.

En appliquant \textsc{Phi-Deref} sous $C = *\ctxEmpty ← 1$, on obtient
$\mm{m_2'}{*q←1}{m_2'}{(x)←1}$.

Puis on applique \textsc{Exp-Cst} sous $C = (x)←\ctxEmpty$ et
$\mm{m_2'}{*q←1}{m_2'}{(x)←\widehat{1}}$.

Maintenant que les deux côtés de $←$ sont évalués, on peut appliquer
\textsc{Exp-Set}, et $\mm{m_2'}{*q←1}{m_2''}{\iPass}$ où:

\[m_2'' = ([[q↦(x)]],
[f↦\widehat{\eFun{q}{*q←1;\iReturn{\eUnit}}}
;x↦\widehat{1}])
\]

Alors, d'après \textsc{Seq}, $\mm{m_2'}{*q←1}{m_2''}{\iReturn{()}}$.
Avec \textsc{Exp-Cst} sous $C = \iReturn{\ctxEmpty}$, on a donc
$\mm{m_2'}{*q←1}{m_2''}{\iReturn{\widehat{()}}}$.

On peut enfin appliquer \textsc{Exp-Call-Ctx} pour en déduire que:

\[
    \mm{m_2}{f(\&x)}{m_2'''}{ \eFun{q}{\iReturn{\widehat{\eUnit}}}(\widehat{\&}(x)) }
\]

Donc d'après \textsc{Exp-Call-Return} (car on a $m_2''' = \cCleanup{m_2'''}$ et
$ \widehat{\eUnit} = \cV{|m_2'''|}{\widehat{\eUnit}} $):

\[
    \mm{m_2}{f(\&x)}{m_2'''}{ \widehat{\eUnit} }
\]

En posant $m_3 = m_2'''$, on a $\ph{m_2}{f(\&x)}{m_3}$.

Donc pour conclure (grâce à \textsc{Prog}), on a
$\phprog{[p_1,p_2,p_3]}{m_3}$.

% TODO intégrer une structure
% TODO intégrer if

\section*{Conclusion}

On vient de définir un langage impératif, \langname. Le but est que celui-ci
serve de support à des analyses statiques, afin notamment de montrer une
propriété de sécurité sur les pointeurs. Pour le moment, on a seulement défini
ce que sont les programmes (leur syntaxe) et comment ils s'exécutent (leur
sémantique). Sur ces deux points, on note que nous sommes restés suffisamment
proches de C, tout en utilisant pour la mémoire un modèle plus structuré qu'une
simple suite d'octets. Les définitions de la syntaxe ainsi que de la sémantique
sont rappelées dans l'annexe~\ref{anx:eval} (sections~\ref{sec:anx-sem-start}
à~\ref{sec:anx-sem-end}). \added{sections entre parenthèses}

Afin de manipuler les états mémoire dans la sémantique d'évaluation, nous avons
utilisé le concept des lentilles, qui permettent de chaîner des accesseurs entre
eux et d'accéder simplement à des valeurs profondes de la mémoire, en utilisant
le même outil pour la lecture et l'écriture.

Pour le moment, on ne peut rien présager de l'exécution d'un programme bien
formé syntaxiquement. Pour la grande majorité des programmes bien formés (à la
syntaxe correcte), l'évaluation s'arrêtera soit par une erreur, soit parce
qu'aucune règle d'évaluation ne peut s'appliquer. Dans les
chapitres~\ref{cha:typbase} et~\ref{cha:qualifs}, nous allons donc définir un
système de types qui permet de rejeter ces programmes se comportant mal à
l'exécution.

% vim: spelllang=fr


\chapter{Typage}

\label{cha:typbase}
Dans ce chapitre, nous enrichissons le langage défini dans le
chapitre~\ref{cha:lang} d'un système de types. Celui-ci permet de séparer les
programmes bien formés, comme celui de la figure~\ref{fig:progwf:good} des
programmes mal formés comme celui de la figure~\ref{fig:progwf:bad}.

\begin{SaveVerbatim}[]{VerbProgGood}
f()
(x=0)
{
  x = 1
  return x
}
\end{SaveVerbatim}

\begin{SaveVerbatim}[]{VerbProgBad}
f()
(x=0)
{
  x = 1
  return (*x)
}
\end{SaveVerbatim}

\begin{figure}

  \centering

  \subfloat[][Programme bien formé]{
    \label{fig:progwf:good}
    \BUseVerbatim{VerbProgGood}
  }
  \hspace{2cm}
  \subfloat[][Programme mal formé]{
    \label{fig:progwf:bad}
    \BUseVerbatim{VerbProgBad}
  }

  \caption{Programmes bien et mal formés}
  \label{fig:progwf}

\end{figure}

Le but d'un tel système de types est de rejeter les programmes qui sont
"évidemment faux", c'est à dire dont on peut prouver qu'il provoqueraient des
erreurs à l'exécution dues à une incompatibilité entre valeurs. En ajoutant
cette étape, on restreint la classe d'erreurs qui pourraient bloquer la
sémantique.

\section{Principe}

Le principe est d'associer à chaque construction syntaxique une étiquette
représentant le genre de valeurs qu'elle produira. Dans le programme de la
figure~\ref{fig:progwf:good}, la variable $x$ est initialisée avec la valeur
$0$, c'est donc un entier. Cela signifie que dans tout le programme, toutes les
instances de cette variable
\footnote{Deux variables peuvent avoir le même nom dans deux fonctions
  différentes, par exemple. Dans ce cas il n'y a aucune contrainte particulière
  entre ces deux variables. L'analyse de typage se fait toujours dans un
  contexte précis.
}
porteront ce type. La première instruction est l'affectation de la constante $1$
(entière) à $x$ dont on sait qu'elle porte des valeurs entières, ce qui est donc
correct. Le fait de rencontrer $\iReturn{x}$ permet de conclure que le type de
la fonction est $() → \tInt$.

Dans la seconde fonction, au contraire, l'opérateur $*$ est appliqué à $x$ (le
début de l'analyse est identique et permet de conclure que $x$ porte des valeurs
entières). Or cet opérateur prend un argument d'un type pointeur de la forme
$t*$ et renvoie alors une valeur de type $t$. Ceci est valable pour tout $t$
(\tInt, \tFloat où même $t'*$ : le déréférencement d'un pointeur sur pointeur
donne un pointeur), mais le type de $x$, \tInt, n'est pas de cette forme. Ce
programme est donc mal typé.

\section{Définitions}

\begin{figure}

  \gramlr{Type}{
    \begin{align*}
      t \gramisa & \tInt                         & \textrm{Entier}
      \\ \gramor & \tFloat                       & \textrm{Flottant}
      \\ \gramor & t[]                           & \textrm{Tableau}
      \\ \gramor & t*                            & \textrm{Pointeur}
      \\ \gramor & \tstruct{s}                   & \textrm{Structure}
      \\ \gramor & (t_1, …, t_n) \rightarrow t_r & \textrm{Fonction}
    \end{align*}
  }

  \gramlr{Environnement de typage}{
    \begin{align*}
      Γ \gramisa & []         & \textrm{Environnement vide}
      \\ \gramor & (a, t)::Γ' & \textrm{Extension}
    \end{align*}
  }

  \caption{Types et environnements de typage}

  \label{fig:les-types}

\end{figure}

Les types associés aux expressions sont décrits dans la
figure~\ref{fig:les-types}.

Pour maintenir les contextes de typage, un environnement $Γ$ associe un type à
un ensemble de variables.

Plus précisément, un environnement $Γ$ est une liste de couples (variable,
type).

Par exemple, $(p, \tInt*) ∈ Γ$ permet de typer (sous $Γ$) l'expression $p$ en
$\tInt*$, $*p$ en $\tInt$ et $p +_p 4$ en $\tInt*$.

Le type des fonctions semble faire apparaître un n-uplet $(t_1, …, t_n)$ mais ce
n'est qu'une notation : il n'y a pas de n-uplets de première classe, ils sont
toujours présents dans un type fonctionnel.

\begin{definition}[Typage d'une expression]

  On note de la manière suivante le fait qu'une expression $e$ (telle que
  définie dans la figure~\ref{fig:stx-data}) ait pour type $t$ dans le contexte
  $Γ$.

  \[
    \ty{Γ}{e}{t}
  \]

\end{definition}

\begin{definition}[Typage d'une instruction]

  Les instructions n'ont en revanche pas de type. Mais il est tout de même
  nécessaire de vérifier que troutes les sous-expressions apparaissant dans une
  instruction sont cohérentes ensemble.

  On note de la manière suivante le fait que sous l'environnement $Γ$
  l'instruction $i$ est bien typée :

  \[
    \tyi{Γ}{i}
  \]

\end{definition}

\begin{definition}[Typage d'une phrase]

  De par leur nature séquentielle, les phrases qui composent un programme
  altèrent l'environnement de typage. Par exemple, la déclaration d'une variable
  globale ajoute une valeur dans l'environnement.

  On note

  \[
    \typh{Γ}{p}{Γ'}
  \]

  si le typage de la phrase $p$ transforme l'environnement $Γ$ en $Γ'$.

  On étend cette notation aux suites de phrases :

  \[
    \begin{cases}
      \typhstar{Γ}{[]}{Γ}  \\
      \typhstar{Γ}{p::ps}{Γ'} \mbox{ si }
        ∃ Γ'',
            \begin{cases}
              \typh{Γ}{p}{Γ''}  \\
              \typhstar{Γ''}{ps}{Γ'}
            \end{cases}
    \end{cases}
  \]

\end{definition}

\begin{definition}[Typage d'un programme]

  Un programme est bien typé si on peut typer sa suite de phrases en partant
  d'un environnement vide. C'est à dire s'il existe un environnement final $Γ_f$
  tel que

  \[
    \typhstar{[]}{P}{Γ_f}
  \]

  Cela est indépendant de tout environnement ; on note alors :

  \[
    ⊢ P
  \]

\end{definition}

\section{Expressions}

\subsection*{Littéraux}

Le typage des littéraux numériques ne dépend pas de l'environnement de typage :
ce sont toujours des entiers ou des flottants.

\begin{mathpar}

  \disprule{Cst-Int}

  \disprule{Cst-Float}

\end{mathpar}

Le pointeur nul, quant à lui, est compatible avec tous les types pointeur.

\begin{mathpar}
  \disprule{Cst-Null}
\end{mathpar}

Enfin, le littéral unité a le type $\tUnit$.

\begin{mathpar}
  \disprule{Cst-Unit}
\end{mathpar}

\subsection*{Left-values}

Rappelons que l'environnement de typage $Γ$ contient le type des variables
accessibles du programme. Le cas où la left-value à typer est une variable est
donc direct : il suffit de retrouver son type dans l'environnement.

\begin{mathpar}
  \disprule{Lv-Var}
\end{mathpar}

Dans le cas d'un déréférencement, on commence par typer la left-value
déréférencée. Si elle a un type pointeur, la valeur déréférencée est du type
pointé.

\begin{mathpar}
  \disprule{Lv-Deref}
\end{mathpar}

Pour une left-value indexée (l'accès à tableau), on s'assure que l'indice soit
entier, et que la left-value a un type tableau : le type de l'élement est encore
une fois le type de base du type tableau ($t$ pour $t[]$).

\begin{mathpar}
  \disprule{Lv-Index}
\end{mathpar}

\begin{mathpar}
  \disprule{Lv-Field}
\end{mathpar}

% TODO

\subsection*{Opérateurs}

Un certain nombre d'opérations est possible sur le type \tInt.

\begin{mathpar}
  \disprule{Op-Int}
\end{mathpar}

De même sur \tFloat.

\begin{mathpar}
  \disprule{Op-Float}
\end{mathpar}

Les opérateurs de comparaison peuvent s'appliquer à deux opérandes qui sont d'un
type qui supporte l'égalité. Ceci est représenté par un jugement
$\textsc{Eq}(t)$ qui est vrai pour les types \tInt, \tFloat et pointeurs. Les
opérateurs $=$ et $≠$ renvoient alors un \tInt.

\begin{mathpar}
\disprule{Eq-Num}

\disprule{Eq-Ptr}

\disprule{Eq-Array}

\disprule{Op-Eq}
\end{mathpar}

%TODO Eq-Struct

Les comparaisons sont plus restrictives, et ne s'appliquent qu'aux types
primitifs (on ne peut pas comparer deux pointeurs, ou deux tableaux).

\begin{mathpar}
  \disprule{Op-Comparable}
\end{mathpar}

% TODO extensions aux structs et tableaux

Les opérateurs unaires "$+$" et "$-$" appliquent aux \tInt, et leurs équivalents
"$+.$" et "$-.$" aux \tFloat.

\begin{mathpar}
\disprule{Unop-Plus-Int}
\and
\disprule{Unop-Plus-Float}

\disprule{Unop-Minus-Int}
\and
\disprule{Unop-Minus-Float}
\end{mathpar}

Les opérateurs de négation unaires, en revanche, ne s'appliquent qu'aux
entiers.

\begin{mathpar}
  \disprule{Unop-Not}
\end{mathpar}

L'arithmétique de pointeurs préserve le type des pointeurs.

\begin{mathpar}
  \disprule{Ptr-Arith}
\end{mathpar}

\subsection*{Autres expressions}

Prendre l'adresse d'une left-value rend un type pointeur sur le type de
celle-ci.

\begin{mathpar}
  \disprule{Addr}
\end{mathpar}

Pour typer une affectation, on vérifie que la left-value (à gauche) et
l'expression (à droite) ont le même type. C'est alors le type résultat de
l'expression d'affectation.

\begin{mathpar}
  \disprule{Set}
\end{mathpar}

Un littéral tableau a pour type $t[]$ où $t$ est le type de chacun de ses
éléments.

\begin{mathpar}
  \disprule{Array}
\end{mathpar}

Pour qu'un littéral de structure soit bien typé, il faut que chacun de ses noms
de champs corresponde à un même nom de type structure, avec le bon type pour
chaque champ.

\begin{mathpar}
  \disprule{Struct}
\end{mathpar}

Pour typer un appel de fonction, on s'assure que la fonction a bien un type
fonctionnel. On type alors chacun des arguments avec le type attendu. Le
résultat est du type de retour de la fonction.

\begin{mathpar}
  \disprule{Call}
\end{mathpar}

\section{Instructions}

La séquence est simple à traiter : l'instruction vide est toujours bien typée,
et la suite de deux instructions est bien typée si celles-ci le sont également.

\begin{mathpar}
  \disprule{Pass}

  \disprule{Seq}
\end{mathpar}

Une instruction constituée d'une expression est bien typée si celle-ci peut être
typée dans ce même contexte.

\begin{mathpar}
  \disprule{Exp}
\end{mathpar}

Les constructions de contrôle sont bien typées si leurs sous-instructions sont
bien typées, et si la condition est d'un type entier.

\begin{mathpar}
  \disprule{If}

  \disprule{While}
\end{mathpar}

\section{Fonctions}

Le typage des fonctions fait intervenir une variable virtuelle $\vRet$. Cela
revient à typer l'instruction $\iReturn{e}$ comme $\vRet ← e$.

\begin{mathpar}
  \disprule{Return}
\end{mathpar}

Pour typer une définition de fonction, on commence par créer un nouvel
environnement de typage $Γ'$ obtenu par la suite d'opérations suivantes :

\begin{itemize}
\item
  on enlève (s'il existe) le couple $\vRet : t_f$ correspondant à la
  valeur de retour de la fonction appelante
\item
  on ajoute les types des arguments $a_i : t_i$
\item
  on ajoute les types des variables locales $l_i : t'_i$
\item
  on ajoute le type de la valeur de retour de la fonction appelée,
  $\vRet : t$
\end{itemize}

Il reste alors à vérifier que les initialiseurs $e_i$ ont le bon type $t'_i$ et
que le corps de la fonction est bien typé sous $Γ'$. Le type de la fonction est
alors $(t_1, …, t_n) → t$.

\begin{mathpar}
  \disprule{Fun}
\end{mathpar}

\section{Phrases}

L'évaluation d'une expression est le cas le plus simple. En effet, il y a juste
à vérifier que celle-ci est bien typable (avec ce type) dans l'environnement de
départ. L'environnement n'est pas modifié.

\begin{mathpar}
  \disprule{Ph-Exp}
\end{mathpar}

La déclaration d'une variable globale commence de la même manière, mais on
enrichit l'environnement de cette nouvelle valeur.

\begin{mathpar}
  \disprule{Ph-Var}
\end{mathpar}

\begin{mathpar}
  \disprule{Ph-Struct}
\end{mathpar}

% TODO

\section{Sûreté du typage}

Comme dit l'adage :

\begin{quote}
"Well-typed programs don't go wrong." (Robin Milner)
\end{quote}

C'est à dire qu'un programme bien typé possède des propriétés de sûreté.

\paragraph{Progrès :} l'évaluation d'un terme bien typé ne reste pas bloquée;
il y a toujours une règle qui s'applique.

\paragraph{Préservation :} l'évaluation d'un terme bien typé produit un terme
bien typé.

Puisqu'il s'agit de propriétés reliant la syntaxe à la sémantique, deux types
d'environnement sont utilisés en même temps. D'une part, les environnement de
typage $Γ$ pour la syntaxe, et d'autre part les environnements mémoire $m$ pour
la sémantique. Il est nécessaire de définir une relation de compatibilité entre
ces deux monde, pour exprimer par exemple qu'un état mémoire contient dans la
variable $x$ un entier.

On commence par énoncer quelques lemmes utiles dans la démonstration de ces
théorèmes.

Les règles précédentes ont la particularité suivante : pour chaque forme
syntaxique, il n'y a souvent qu'une règle qui peut s'appliquer. Cela permet de
déduire quelle règle il faut appliquer pour vérifier (ou inférer) le type d'une
expression.

\begin{lemma}[Inversion] \label{lemma:inversion}

  À partir d'un jugement de typage, on peut en déduire des inforamtions sur les
  types de ses sous-expressions.

\begin{itemize}
\item
  Constantes
  \begin{itemize}
    \item si $Γ ⊢ n : t$, alors $t = \tInt$
    \item si $Γ ⊢ d : t$, alors $t = \tFloat$
    \item si $Γ ⊢ \eNull : t$, alors $∃ t', t = t'*$
    \item si $Γ ⊢ \eUnit : t$, alors $t = \tUnit$
  \end{itemize}

\item Références mémoire :
  \begin{itemize}
    \item
      si $Γ ⊢ x : t$, $x : t ∈ Γ$
    \item
      si $Γ ⊢ *lv : t$, alors $Γ ⊢ lv : t*$
    \item
      si $Γ ⊢ lv[e] : t$, alors $Γ ⊢ lv : t[]$ et $Γ ⊢ e : \tInt$
    \item
      si $Γ ⊢ lv.l : t$, alors $Γ ⊢ lv : \{ l: t ; … \}$
    % TODO structure
  \end{itemize}
\item
  Appel de fonction : si $Γ ⊢ e (e_1, …, e_n) : t$, il existe $(t_1, …, t_n)$
  tels que

  \[
    \begin{cases}
      Γ ⊢ e : (t_1, …, t_n) → t \\
      ∀ i ∈ [1;n], Γ ⊢ e_i : t_i
    \end{cases}
  \]

\item Fonction : si $Γ ⊢ \mathrm{fun} (a_1, …, a_n) ((l_1, e_1), …, (l_p, e_p))
  \{i\} : t$, alors il existe $(t_1, …, t_n)$ et $t'$ tels que $t' = (t_1, …,
  t_n) → t$.

  % TODO ...
\end{itemize}

\end{lemma}

\begin{proof}

  Pour chaque jugement, on considère les règles qui peuvent amener à cette
  conclusion.

\begin{itemize}
\item
  Références mémoire :

  \begin{itemize}

    \item $Γ ⊢ x : t$

      La seule règle de cette forme est \textsc{Lv-Var}. Puisque sa prémisse est
      vraie, on en conclut que $x : t ∈ Γ$.

    \item $Γ ⊢ *φ : t$

      De même, seule la règle \textsc{Lv-Deref} convient. On en conclut que $Γ ⊢ φ : t*$.

    \item $Γ ⊢ φ[] : t$

      Idem avec \textsc{Lv-Index}.

    \item $Γ ⊢ φ.l : t$

      $Γ ⊢ φ : \{ l: t ; … \}$ %TODO sûrement différent sans S

  \end{itemize}

\item
  Appel de fonction : pour en arriver à $Γ ⊢ e (e_1, …, e_n) : t$, seule la
  règle \textsc{Call} s'applique, ce qui permet de conclure.

\item Fonction : la seule règle possible pour conclure une dérivation de

  \[
    Γ ⊢ \mathrm{fun} (a_1, …, a_n) ((l_1, e_1), …, (l_p, e_p)) \{i\} : t
  \]
  est \textsc{Fun}.

\end{itemize}
\end{proof}

Il est aussi possible de réaliser l'opération inverse : à partir du type d'une
valeur, on peut déterminer sa forme syntaxique. C'est bien sûr uniquement
possible pour les valeurs, pas pour n'importe quelle expression (par exemple
l'expression $x$ (variable) peut avoir n'importe quel type $t$ dans le contexte
$Γ = x:t$).

\begin{lemma}[Formes canoniques] \label{lemma:canon}

  Il est possible de déterminer la forme syntaxique d'une valeur étant donné son
  type.

  \begin{itemize}
  \item si $Γ ⊢ v : \tInt$, $v = n$.
  \item si $Γ ⊢ v : \tFloat$, $v = d$.
  \item si $Γ ⊢ v : \tUnit$, $v = \eUnit$.
  \item si $Γ ⊢ v : t*$, $v = φ$ ou $v = \eNull$.
  \item si $Γ ⊢ v : (t_1, …, t_n) → t$, $v = f$.
  \item si $Γ ⊢ v : t*$, $v = φ$.
  \item si $Γ ⊢ v : t[]$, $v = \{v_1; …; v_n\}$.
    %TODO structures
  \item si $Γ ⊢ v : (t_1, …, t_n) → t$, $v = \mathrm{fun}
    (a_1, …, a_n) ((l_1, e_1), …, (l_p, e_p))$.
  \end{itemize}

\end{lemma}

Ces lemmes étant établis, on énonce maintenant le théorème de progrès accompagné
d'une notion de compatibilité entre envionnements de typage $Γ$ et états mémoire
$m$.

\begin{definition}[Compatibilité mémoire]

  Soient $Γ$ un environnement de typage et $m = x_1 ↦ v_1, …, x_n ↦ v_n$ un état
  mémoire. On dit que $m$ est compatible avec $Γ$ si

  \[
    ∃ (t_1, …, t_n),
    ∀ i ∈ [1;n],
    \begin{cases} Γ ⊢ x_i : t_i
               \\ Γ ⊢ v_i : t_i
    \end{cases}
  \]

  On note alors $\mcomp{Γ}{m}$.

  % TODO les états mémoires ne sont pas linéaires
  % TODO quid de ce qui n'est pas visible (ie, le tas)

\end{definition}

\begin{theorem}[Progrès]
  \label{thm:progres}

  Supposons que $Γ ⊢ e : t$. Soit $m$ un état mémoire tel que $\mcomp{Γ}{m}$.
  Alors l'un des cas suivant est vrai :

\begin{itemize}
  \item $∃ v, e = v$
  \item $∃ (e', m'), \mcomp{Γ}{m'} ∧ \mm{m}{e}{m'}{e'}$
  \item $∃ Ω ∈ \{Ω_{div},Ω_{array},Ω_{ptr}\}, \msi{m}{e} → Ω$
\end{itemize}

  C'est à dire, soit :

\begin{itemize}
  \item $e$ est complètement évaluée
  \item un pas d'évaluation préservant la compatibilité mémoire est possible
  \item une erreur de division, tableau ou pointeur se produit
\end{itemize}

\end{theorem}

\begin{proof}

  On procède par induction sur la dérivation de $Γ ⊢ e : t$, et plus précisément
  sur la dernière règle utilisée.

  \paragraph{\textsc{Cst-Int} :} % {{{
$e$ est alors de la forme $n$, qui est une valeur.
%}}}
  \paragraph{\textsc{Cst-Float} :} % {{{
$e$ est alors de la forme $d$, qui est une valeur.
%}}}
  \paragraph{\textsc{Cst-Null} :} % {{{
$e$ est alors égale à $\eNull$, qui est une valeur.
%}}}
  \paragraph{\textsc{Cst-Unit} :}%{{{
$e$ est alors égale à \eUnit, qui est une valeur.
%}}}
\paragraph{\textsc{Lv-Var} :}%{{{

Puisque $(x, t) ∈ Γ$ $\mcomp{Γ}{m}$, il existe $(x ↦ v) ∈ m$. La règle
d'évaluation \textsc{Phi-Var} s'applique donc.

% TODO prouver le coup du lookup
%}}}
\paragraph{\textsc{Lv-Deref} :}%{{{

  Appliquons l'hypothèse de récurrence à $lv$ (vu en tant qu'expression).

\begin{itemize}
\item
  $lv = v$. Puisque $Γ ⊢ v : t*$, on déduit du
  lemme~\ref{lemma:canon} que $v = φ$ ou $v = \eNull$.

  Dans le premier cas, la règle \textsc{Phi-Deref} s'applique :
  $\mms{e}{\widehat{*}φ}$.
  Dans le second, puisque $\msi{m}{*\eNull} → Ω_{ptr}$, on a
  $\msi{m}{e} → Ω_{ptr}$.

\item
  $\mm{m}{lv}{m'}{e'}$.
  De \textsc{Ctx} avec $C = *\ctxEmpty$, on obtient
  $\mm{m}{e}{m'}{*e'}$.

\item
  $\msi{m}{lv} → Ω$.
  En appliquant \textsc{Eval-Err} avec $C = *\ctxEmpty$, on obtient
  $\msi{m}{e} → Ω$.

\end{itemize}

% }}}
\paragraph{\textsc{Lv-Index} :} % %{{{

De même, on applique l'hypothèse de récurrence à $lv$.

\begin{itemize}
\item $lv = v$.

Comme $Γ ⊢ v : t[]$, on déduit du lemme~\ref{lemma:canon} que
$v = \{v_1; …; v_p\}$.
Appliquons l'hypothèse de récurrence à $e$.

\begin{itemize}
\item $e = v'$. Puisque $Γ ⊢ e : \tInt$, on réapplique le
lemme~\ref{lemma:canon} et $v' = n$. Deux cas sont à distinguer :
si $n ∈ [0;p-1]$, la règle \textsc{Phi-Array} s'applique et
$\mm{m}{lv[e]}{m}{v_{n+1}}$.
% TODO réécrire bien Phi-Array
% TODO attention à l'off by one
Sinon, $\msi{m}{lv[e]} → Ω_{array}$.
% TODO pourquoi ?

\item $\mm{m}{e}{m'}{e'}$.
En appliquant \textsc{Ctx} avec $C = v[\ctxEmpty]$, on en déduit
\item $\mm{m}{lv[e]}{m'}{lv[e']}$.

\item $\msi{m}{e} → Ω$.
Avec \textsc{Eval-Err} sous ce même contexte,
$\msi{m}{lv[e]} → Ω$
\end{itemize}

\item $\mm{m}{lv}{m'}{e'}$.
On applique alors \textsc{Ctx} avec $C = \ctxEmpty[e]$, et
$\mm{m}{lv[e]}{m'}{e'[e]}$.

\item $\msi{m}{lv} → Ω$.
Toujours avec $C = \ctxEmpty[e]$, de \textsc{Eval-Err} il vient
$\msi{m}{lv[e]} → Ω$.

\end{itemize}
%}}}
\paragraph{\textsc{Lv-Field} :} % TODO
  \paragraph{\textsc{Op-Int} :} % {{{

  Cela implique que $e = e_1~\opbin~e_2$. Par le lemme~\ref{lemma:inversion}, on
  en déduit que $Γ ⊢ e_1 : \tInt$ et $Γ ⊢ e_2 : \tInt$.

  Appliquons l'hypothèse de récurrence sur $e_1$. Trois cas peuvent se produire.

\begin{itemize}

  \item $e_1 = v_1$. On a alors $\mm{m}{e_1}{m'}{v_1}$ avec $m' = m$.

    On applique l'hypothèse de récurrence à $e_2$.

      \begin{itemize}

        \item $e_2 = v_2$ : alors $\mm{m'}{e_2}{m''}{v_2}$ avec $m'' = m$. On
          peut alors appliquer \textsc{Exp-BinOp}.

        \item $∃ (e'_2, m''), \mm{m'}{e_2}{m''}{e'_2}$.

          En appliquant \textsc{Ctx} avec $C = \ctxOp{v_1}{\ctxEmpty}$, on
          en déduit $\mm{m'}{v_1~\opbin~e_2}{m''}{v_1~\opbin~e'_2}$ soit
          $\mm{m}{e}{m''}{v_1~\opbin~e'_2}$.

        \item $\msi{m'}{e_2} → Ω$.
          De \textsc{Eval-Err} avec $C = \ctxOp{v_1}{\ctxEmpty}$
          vient alors $\msi{m}{e} → Ω$.

      \end{itemize}

  \item $∃(e', m'), \mm{m}{e_1}{m'}{e'_1}$.
    En appliquant \textsc{Ctx} avec $C = \ctxOp{\ctxEmpty}{e_2}$, on obtient
    $\mm{m}{e_1~\opbin~e_2}{m'}{e'_1~\opbin~e_2}$, ou
    $\mm{m}{e}{m'}{e'_1~\opbin~e_2}$.

  \item $\msi{m}{e_1} → Ω$.
    D'après \textsc{Eval-Err} avec $C = \ctxOp{\ctxEmpty}{e_2}$, on a
    $\msi{m}{e} → Ω$.

\end{itemize}

% }}}
\paragraph{\textsc{Op-Float} :} % {{{
Ce cas est similaire à \textsc{Op-Int}.
% TODO quid du lemme d'inversion (premiere ligne de Op-Int)?
%}}}
\paragraph{\textsc{Op-Eq} :} %{{{
Ce cas est similaire à \textsc{Op-Int}.
% TODO expand un peu
% TODO quid du lemme d'inversion (premiere ligne de Op-Int)?
%}}}
\paragraph{\textsc{Op-Comparable} :} %{{{
Ce cas est similaire à \textsc{Op-Int}.
% TODO expand un peu
% TODO quid du lemme d'inversion (premiere ligne de Op-Int)?
%}}}
\paragraph{\textsc{Unop-Plus-Int} :} % {{{

Alors $e = +~e_1$. En appliquant l'hypothèse d'induction sur $e_1$ :

\begin{itemize}
\item
  soit $e_1 = v_1$. Alors en appliquant \textsc{Exp-UnOp},
  $\mm{m}{+~v_1}{m}{\widehat{+}~v_1}$, c'est à dire en posant $v =
  \widehat{+}~v_1$, $\mm{m}{e}{m}{v}$.
% TODO écrire la règle
\item
  soit $∃ e'_1, m', \mm{m}{e_1}{m'}{e'_1}$. Alors en appliquant \textsc{Ctx}
avec $C = +~\ctxEmpty$, on obtient $\mm{m}{e}{m'}{e'_1}$.
\item
  soit $\msi{m}{e_1} → Ω$.
  De \textsc{Eval-Err} avec $C = +~\ctxEmpty$ il vient$\msi{m}{e} → Ω$.
\end{itemize}

% }}}
\paragraph{\textsc{Unop-Plus-Float} :} % {{{
Ce cas est similaire à \textsc{Unop-Plus-Int}.
% }}}
\paragraph{\textsc{Unop-Minus-Int} :} % {{{
Ce cas est similaire à \textsc{Unop-Plus-Int}.
% }}}
\paragraph{\textsc{Unop-Minus-Float} :} % {{{
Ce cas est similaire à \textsc{Unop-Plus-Int}.
% }}}
\paragraph{\textsc{Unop-Not} :}%{{{
Ce cas est similaire à \textsc{Unop-Plus-Int}.
%}}}
\paragraph{\textsc{Ptr-Arith} :} % %{{{ TODO

\[ \left( \disprule{Ptr-Arith} \right) \]

%}}}
\paragraph{\textsc{Addr} :} % {{{ TODO

\[ \left( \disprule{Addr} \right) \]

%}}}
\paragraph{\textsc{Set} :} % {{{ TODO

\[ \left( \disprule{Set} \right) \]

%}}}
\paragraph{\textsc{Array} :} % {{{ TODO

\[ \left( \disprule{Array} \right) \]

%}}}
\paragraph{\textsc{Struct} :} % {{{ TODO

\[ \left( \disprule{Struct} \right) \]

%}}}
\paragraph{\textsc{Call} :} % {{{ TODO

\[ \left( \disprule{Call} \right) \]

%}}}
\paragraph{\textsc{Fun} :} % {{{ TODO

\[ \left( \disprule{Fun} \right) \]

%}}}
% TODO vérifier qu'il n'y a pas d'autre règle
% TODO général : que désigne e. Ex avec lv[e]
\end{proof}

\begin{lemma}[Permutation]
  L'ordre dans lequel les variables apparaissent dans un environnement
  n'influe pas sur la relation de typage.

  Pour toute permutation $σ$ de $[1;n]$, on note $σ(x_1 : t_1, …, x_n : t_n) =
  x_{σ(1)} : t_{σ(1)}, … x_{σ(n)} : t_{σ(n)}$.

  Alors : si $Γ ⊢ e : t$ et $Γ' = σ(Γ)$, alors $Γ' ⊢ e : t$.
\end{lemma}

\begin{lemma}[Affaiblissement]
  De même que l'ordre n'influe pas le typage, on peut aussi ajouter des
  associations supplémentaires dans l'environnement sans modifier les typages
  dans cet environnement.

  Si $Γ ⊢ e : t$ et $x ∉ \mathrm{dom}(Γ)$, alors $Γ, x : t' ⊢ e : t$.
\end{lemma}

\begin{lemma}[Substitution]
  Si dans une expression $e$ il apparait une variable $x$ de type $t'$, le
  typage est préservé lorsqu'on remplace ses occurrences par une expression $e'$
  de même type.

  Si $Γ, x : t' ⊢ e : t$ et $Γ ⊢ e' : t'$, alors $Γ ⊢ e [x/e'] : t$.
\end{lemma}

Ces lemmes permettent de prouver le théorème suivant :

\begin{theorem}[Préservation]
  \label{thm:preservation}

  Si une expression est typable et que son évaluation produit une valeur, alors
  cette valeur est du même type que l'expression.

  Si $Γ ⊢ e : t$ et $e → v$ % TODO

  alors $Γ ⊢ v : t$.

\end{theorem}

\section*{TODO}

\begin{itemize}
\item ordre des sections
\item versions $Γ ⊢ i$ des propriétés
\item preuve de progres : état mémoire : doublet/triplet
\item définir les opérations d'ajout/remplacement sur les contextes de typage
\end{itemize}


%\chapter{Sémantique statique}

%\label{cha:typbase}
%Ici nous enrichissons le langage défini dans le chapitre~\ref{cha:lang} d'un
système de types. Celui-ci permet d'obtenir plus de garanties que celui de C tel
que décrit dans \cite{AnsiC}.

Il permet le polymorphisme sur les types pointeurs, permettant par exemple de
typer :

\begin{mathpar}
⊢ \text{memcpy} : ∀ a . (a^*, a^*, \tInt) \rightarrow \tVoid
\end{mathpar}

\section{Règles de typage}

\subsection{Types}

Dans cette section, on définit la notion de programme bien typé. L'analyse par
typage permet de vérifier qu'à chaque expression on peut associer un type, et ce
de manière cohérente entre plusieurs utilisations d'une variable.

Les types des valeurs sont :

\gramlr{Types}{
\begin{align*}
τ  \gramisa & \tInt, \tFloat, \tVoid        & \textrm{Constante}
\\ \gramor  &  a                            & \textrm{Variable}
\\ \gramor  & (τ_1, …, τ_n) \rightarrow τ_r & \textrm{Fonction}
\\ \gramor  & [τ]                           & \textrm{Tableau}
\\ \gramor  & τ*                            & \textrm{Pointeur}
\\ \gramor  & \{ f_1:τ_1
              ,    …
              , f_n:τ_n \}                  & \textrm{Structure}
\end{align*}
}

L'ensemble des types possibles (défini inductivement ci-dessus) sera noté
\textsc{Typ}, et l'ensemble des variables de type par \textsc{VarTyp}.

\subsection{Schémas de type}

On va associer à chaque variable globale un type. Mais faire de même pourrait
être trop restrictif. En effet, une fonction comme memcpy peut être utilisée
pour copier des tableaux d'entiers, mais aussi de flottants. On va donc associer
un schéma de types à chaque fonction.

\gramlr{Schémas}{
\begin{align*}
σ \gramisa & ∀ \vec{a} . τ
\end{align*}
}

Un schéma de types correspond à un ensemble de types. Prenons l'exemple de la
fonction identité : elle a pour schéma de types $∀ a . a \rightarrow a$, ce qui
signifie que pour chaque type $τ$, on peut l'utiliser avec le type $τ
\rightarrow τ$. Plus précisément, cela veut dire que puisque $a$ est quantifiée,
on peut le substituer par n'importe quel autre type.

\begin{definition}[Substitution]
Une substitution est une fonction partielle de \textsc{VarTyp} dans
\textsc{Typ}. Elle sera notée par exemple $s = \{ a ↦ \tInt, b ↦ (\tFloat
\rightarrow \tInt)\}$.

Le support d'une substition, noté $\textrm{Supp}(s)$, est son domaine de
définition.

On définit aussi l'application d'une substitution sur un type quelconque : si
$s$ est une substitution, $\overline{s}$ est son extension définie par :

\begin{align*}
\overline{s}(a)   =~& s(a) & \textrm{si $a$ est une variable}  \\
\overline{s}(c)   =~& c    & \textrm{si $c$ est une constante} \\
\overline{s}([τ]) =~& [\overline{s}(τ)] & \\
\overline{s}(τ*)  =~& \overline{s}(τ)*   & \\
\overline{s}(τ_1, …, τ_n) \rightarrow τ_r =~& (\overline{s}(τ_1), …, \overline{s}(τ_n)) \rightarrow \overline{s}(τ_r) \\
\overline{s}(\{ f_1:τ_1, … ,f_n:τ_n \}) =~& \{ f_1:\overline{s}(τ_1), …, f_n:\overline{s}(τ_n) \}
\end{align*}

Par souci de simplicité, on notera $s$ pour $\overline{s}$.
\end{definition}

\begin{definition}[Instanciation]
Un schéma de types $σ = ∀ \vec{a} . τ$ peut être instancié en un type $μ$ s'il
existe une substition $s$ telle que :

\begin{itemize}
\item $\textrm{Supp}(s) ⊆ \vec{a}$
\item $s(τ) = μ$
\end{itemize}

On note alors $μ ≤ σ$.
\end{definition}

\begin{definition}[Variables libres]
Les variables libres d'un type sont l'ensemble des variables de types qui
apparaissent dans celui-ci :

\begin{align*}
\tFV{a}   =~& ∅    & \textrm{si $a$ est une variable}  \\
\tFV{c}   =~& ∅    & \textrm{si $c$ est une constante} \\
\tFV{[τ]} =~& \tFV{τ} & \\
\tFV{τ*}  =~& \tFV{τ} & \\
\tFV{(τ_1, …, τ_n) \rightarrow τ_r} =~& \tFV{τ_1} ∪ … ∪ \tFV{τ_n} ∪ \tFV{τ_r} \\
\tFV{\{ f_1:τ_1, … ,f_n:τ_n \}} =~& \tFV{τ_1} ∪ … ∪ \tFV{τ_n}
\end{align*}
\end{definition}

\begin{definition}[Généralisation]
La généralisation consiste à construire un schéma de type à partir d'un type, en
quantifiant sur les variables libres :

\[
\textrm{Gen}(τ) = ∀ \vec{a} . τ \quad \textrm{où} \quad \vec{a} = \tFV{τ}
\]
\end{definition}

En associant un schéma de type $σ$ à une fonction $f$, on indique que la
fonction pourra être utilisée avec tout type $τ$ qui est une instanciation de
$σ$.

\subsection{Environnements de typage}

Chaque jugement de typage est effectué dans un environnement de typage $Γ$
particulier, qui contient le contexte nécessaire : ici, le type des fonctions et
variables du programme.

\gramlr{Environnements}{
\begin{align*}
Γ  \gramisa & (Γ_{fun}, Γ_{var}) & \textrm{Fonctions, variables}
\end{align*}

\begin{align*}
Γ_{fun} \gramisa & ε           & \textrm{Environnement vide}
\\      \gramor  & Γ_{fun},f:σ & \textrm{Ajout d'une fonction}
\end{align*}

\begin{align*}
Γ_{var} \gramisa & ε           & \textrm{Environnement vide}
\\      \gramor  & Γ_{var},v:τ & \textrm{Ajout d'une variable}
\end{align*}
}

Lorsque ce n'est pas ambigü, si $Γ=(Γ_{fun},Γ_{var})$ on notera les extensions
d'environnement $Γ,f:σ$ pour $((Γ_{fun},f:σ), Γ_{var})$ et $Γ,x:τ$ pour
$(Γ_{fun},(Γ_{var},x:τ))$. De même, on notera $(f, σ) ∈ Γ$ si $(f, σ) ∈
Γ_{fun}$, ou $(x, τ) ∈ Γ$ si $(x, τ) ∈ Γ_{var}$.

\subsection{Jugements de typage}

Un des principes du typage est d'associer à chaque expression un type, qui
décrit le genre des valeurs produites par l'évaluation de cette fonction.

\begin{definition}[Jugement de typage]
Un jugement de typage est de la forme $Γ ⊢ e : τ$ et se lit "sous $Γ$, $e$ est
typable en $τ$".
\end{definition}

Les instructions et blocs, au contraire, n'ont pas de type.

\begin{definition}[Bloc bien typé]
On note $Γ ⊢ i$ pour "sous $Γ$, $i$ est bien typé", c'est à dire que ces sous
expressions sont typables en accord avec le flot de données (par exemple, pour
que l'instruction $lv \leftarrow e$ soit bien typée sous $Γ$, il faut que les
types de $lv$ et de $e$ puissent avoir le même type sous $Γ$).
\end{definition}

Le cas des fonctions est particulier puisque celles-ci ont un schéma de types
qui leur est associé.

\begin{definition}[Fonction bien typée]
On note $Γ ⊢ f : σ$ le fait qu'une fonction $f$ est typable en un schéma $σ$ dans $Γ$.
\end{definition}

Enfin, la notion de programme bien typé est intrinsèque : elle se fait
indépendemment d'un environnement externe.

\begin{definition}[Programme bien typé]
Un programme $P$ est bien typé s'il existe un environnement $Γ$ permettant de bien
typer toutes les composantes (fonctions, globales et bloc d'initialisation) d'un
programme. On notera alors $⊢ P$.
\end{definition}

\subsection{Programme}

Au niveau global, un programme P est bien typé (noté $⊢ P$) s'il existe un
environnement $Γ = (\vec{σ}, \vec{τ})$ permettant de typer ses composantes (les fonctions, les
globales et le bloc d'initialisation).

\[
\irule{Prog}{
  (\vec{σ}, \vec{τ}) ⊢ \vec{f} : \vec{σ} \\
  (\vec{σ}, \vec{τ}) ⊢ \vec{x} : \vec{t} \\
  (\vec{σ}, \vec{τ}) ⊢ b
}{
  ⊢ (\vec{f}, \vec{x}, b)
}
\]

\subsection{Flot de contrôle}

Les règles suivantes permettent de définir les jugements $Γ ⊢ i$. De manière
générale, les instructions sont bien typées si leurs sous-instructions sont bien typées.

\begin{mathpar}
\irule{Pass}{
}{
  Γ ⊢ ε
}
\and
\irule{Seq}{
  Γ ⊢ s \\
  Γ ⊢ b
}{
  Γ ⊢ s;b
}
\and
\irule{While}{
  Γ ⊢ b
}{
  Γ ⊢ \npkWhile{b}
}
\and
\irule{Goto}{
}{
  Γ ⊢ \npkGoto{l}
}
\and
\irule{DoWith}{
  Γ ⊢ b
}{
  Γ ⊢ \npkDoWith{b}{l}
}
\end{mathpar}

Dans le cas de la conditionnelle, il est en plus nécessaire de vérifier que la
condition est un entier.

\begin{mathpar}
\irule{If}{
  Γ ⊢ e : \tInt \\
  Γ ⊢ i_t \\
  Γ ⊢ i_f
}{
  Γ ⊢ \npkIf{e}{i_t}{i_f}
}
\end{mathpar}

\subsection{Left values}

On associe à chaque left-value un type, qui est aussi le type des valeurs que
celle-ci peut contenir. Le cas des variables est direct : si un couple
(variable, type) est dans l'environnement de typage, la variable possède ce
type.

\begin{mathpar}
\irule{Lv-Var}{
  (v, τ) ∈ Γ
}{
  Γ ⊢ v : τ
}
\end{mathpar}

Si une expression a un type pointeur, en la déréférençant on obtient une valeur
du type pointé.

\begin{mathpar}
\irule{Lv-Deref}{
  Γ ⊢ e : τ*
}{
  Γ ⊢ *e : τ
}
\end{mathpar}

\todo{accès à un élément d'un type composite}

\begin{mathpar}
\irule{Lv-Field}{
  Γ ⊢ lv : τ_s \\
  (f, τ_f) ∈ τ_s
}{
  Γ ⊢ lv.f : τ_f
}
\and
\irule{Lv-Array}{
  Γ ⊢ lv : [τ] \\
  Γ ⊢ e : \tInt
}{
  Γ ⊢ lv[e] : τ
}
\end{mathpar}

\subsection{Expressions}

Les constantes ont leurs types habituels. Notons que le pointeur nul (\cNil) a
un type polymorphe.

\begin{mathpar}
\irule{Const-Int}{
}{
  Γ ⊢ n : \tInt
}
\and
\irule{Const-Float}{
}{
  Γ ⊢ f : \tFloat
}
\and
\irule{Const-Nil}{
}{
  Γ ⊢ \cNil : τ*
}
\end{mathpar}

Un certain nombre d'opérations est possible sur le type \tInt.

\begin{mathpar}
\irule{Op-Int}{
  \textrm{op} ∈ \{+,-,\times,/,\&,|,\opxor,\&\&,||,⋘,⋙\} \\
  Γ ⊢ e_1 : \tInt \\
  Γ ⊢ e_2 : \tInt
}{
  Γ ⊢ e_1~\textrm{op}~e_2 : \tInt
}
\end{mathpar}

De même sur \tFloat.

\begin{mathpar}
\irule{Op-Float}{
  \textrm{op} ∈ \{+.,-.,\times.,/.\} \\
  Γ ⊢ e_1 : \tFloat \\
  Γ ⊢ e_2 : \tFloat
}{
  Γ ⊢ e_1~\textrm{op}~e_2 : \tFloat
}
\end{mathpar}

Les opérateurs de comparaison peuvent s'appliquer à deux opérandes de types
"comparables". On introduit donc un jugement $\textsc{Comparable}(τ)$ qui est
vrai pour les types \tInt, \tFloat et pointeurs. Les comparaisons renvoient
alors un \tInt.

\begin{mathpar}
\irule{Cmp-Num}
  {τ ∈ \{ \tInt, \tFloat \}}
  {\textsc{Comparable}(τ)}
\and
\irule{Cmp-Ptr}
  { }
  {\textsc{Comparable}(τ*)}
\and
\irule{Op-Cmp}{
  \textrm{op} ∈ \{=,≠,≤,≥,<,>\} \\
  Γ ⊢ e_1 : τ \\
  Γ ⊢ e_2 : τ \\
  \textsc{Comparable}(τ)
}{
  Γ ⊢ e_1~\textrm{op}~e_2 : \tInt
}
\end{mathpar}

L'opérateur unaire "$-$" peut être appliqué à un \tInt ou à un \tFloat.

\begin{mathpar}
\irule{Unop-Minus}{
  τ ∈ \{ \tInt, \tFloat \} \\
  Γ ⊢ e : τ
}{
  Γ ⊢ -e : τ
}
\end{mathpar}

Les opérateurs de négation unaires, en revanche, ne s'appliquent qu'aux
entiers.

\begin{mathpar}
\irule{Unop-Not}{
  \textrm{op} ∈ \{ \textasciitilde, ! \} \\
  Γ ⊢ e : \tInt
}{
  Γ ⊢ \textrm{op}~e : \tInt
}
\end{mathpar}

On peut prendre l'adresse d'une left-value, et obtenir un pointeur vers
celle-ci. Le type du résultat est un type pointeur vers le type de base.

\begin{mathpar}
\irule{AddrOf}{
  Γ ⊢ lv : τ
}{
  Γ ⊢ \&lv : τ*
}
\end{mathpar}

La règle concernant les pointeurs sur fonction est similaire, à ceci près qu'une
fonction a un schéma de type, qu'il faut instancier afin d'obtenir un type
pointeur.

\begin{mathpar}
\irule{AddrOfFun}{
  Γ ⊢ f : σ \\
  τ ≤ σ
}{
  Γ ⊢ \&f : τ
}
\end{mathpar}

\subsection{Fonctions}

Pour typer une fonction, on commence par ajouter ses paramètres dans
l'environnement de typage, et on type la définition de la fonction. Le type
résultant est généralisé.

\todo{Vérifier qu'il n'y a pas de problème entre polymorphisme et mutabilité}

\todo{Notation vecteur}

\[
\irule{Fun}{
  Γ' = Γ ⊕ \{ args(f) = \vec{τ} \} \\
  Γ' ⊢ body(f) \\
  Γ' ⊢ !ret : τ_r
}{
  Γ ⊢ f : Gen(\vec{τ} \rightarrow τ_r, Γ)
}
\]

\subsection{Instructions}

Pour typer une déclaration, il suffit de rajouter la variable avec n'importe
quel type dans l'environnement de typage.

\begin{mathpar}
\irule{Decl}{
  Γ , x : τ ⊢ b
}{
  Γ ⊢ \npkDecl{x}{b}
}
\end{mathpar}

Une affectation est bien typée si elle faite entre une left-value et une
expression de même type.

\begin{mathpar}
\irule{Assign}{
  Γ ⊢ lv : τ \\
  Γ ⊢ e  : τ
}{
  Γ ⊢ lv \leftarrow e
}
\end{mathpar}

\wip{}

\begin{mathpar}
\irule{Fcall}{
  Γ ⊢ lv : τ_{ret} \\
  Γ ⊢ fe : σ \\
  Γ ⊢ \vec{e} : \vec{τ} \\
  (\vec{τ} \rightarrow τ_r) ≤ σ
}{
  Γ ⊢ lv \leftarrow fe(\vec{e})
}
\end{mathpar}

\section{Limitations}
\subsection{Programmes non typables}
\subsection{Incohérences}


\chapter{Analyse de provenance de pointeurs}

\label{cha:qualifs}
\epigraph{
  A friend of mine in a compiler writing class produced a compiler with
  one error message ``you lied to me when you told me this was a program''.
  }
  {--- \textup{Pete Fenelon}}

Jusqu'ici \langname est un langage impératif généraliste, ne prenant pas en
compte les spécificités de l'adressage utilisé dans les systèmes d'exploitation.

Dans ce chapitre, on commence par l'étendre en ajoutant des constructions
modélisant les variables présentes dans l'espace utilisateur (cf.
chapitre~\ref{cha:os}). Celles-ci restent accessibles, mais nécessitent
d'utiliser un opérateur de déréférencement sûr qui vérifie à l'exécution que
l'invariant suivant est respecté :

\begin{quote}
Les pointeurs dont la valeur est contrôlée par l'utilisateur, pointent vers
l'espace utilisateur.
\end{quote}

Pour prouver que cette approche est bien fondée, on procède en plusieurs étapes.

Tout d'abord, on définit un nouveau type d'erreur $\serr{taint}$, déclenché
lorsqu'un pointeur contrôlé par l'utilisateur et pointant vers le noyau (le cas
que l'on cherche à éviter) est déréférencer\footnote{Il est important de noter
  que ce cas d'erreur est ``virtuel'', la distinction entre pointeur contrôlé
  par l'utilisateur et par le noyau n'étant normalement pas visible à
l'exécution}.

Ensuite, on montre qu'avec cet ajout, si on étend naïvement le système de types
en donnant le même type aux pointeurs contrôlés par l'utilisateur et le noyau,
le théorème de progrès (\ref{thm:progres}) n'est plus valable.

L'étape suivante est d'étendre, à son tour, le système de types de \langname en
distinguant les types des pointeurs contrôlés par l'utilisateur des pointeurs
contrôlés par le noyau. Puisqu'on veut interdire le déréférencement des premiers
par l'opérateur \texttt{*}, on modifie également le typage de celui-ci.

Enfin, une fois ces modifications faites, on prouve enfin que les propriétés de
progrès et de préservation sont rétablies.

\section{Extensions noyau pour \langname}

On ajoute à \langname la notion de valeur provenant de l'espace utilisateur.
Pour marquer la séparation entre les deux espaces d'adressage, on ajoute une
construction $φ \gramisa \vTainted{φ'}$. Le chemin interne $φ'$ correspond à une
variable classique (un pointeur noyau) qui, vu sous la forme $\vTainted{φ'}$,
est un pointeur utilisateur.

En plus du déréférencement par \texttt{*} qui devra donc renvoyer
$\serr{taint}$ pour les valeurs de la forme $\vTainted{φ'}$), il faut aussi
ajouter des constructions de lecture et d'écriture à travers les pointeurs
utilisateurs. On note $\uGet{lv}{e}$ la lecture sûre de $e$ (pointeur
utilisateur) dans $lv$ et $\uPut{e}{e'}$ l'écriture de $e$ dans l'espace pointé
par le pointeur utilisateur $e'$.

% TODO pourquoi pas des expressions ?

Illustrons ceci par un exemple. Imaginons un anière, pour écrire un appel
système fictif qui renvoie la version du noyau, en remplissant par pointeur une
structure contenant les champs entiers \texttt{major}, \texttt{minor} et
\texttt{patch}.

% TODO depend de eFun
% TODO dépend de tStruct

\begin{align*}
  & \mathrm{sys\_getver} = \mathrm{fun}(p) \{                      &             & & \\
  & \indent \uPut{ 3}{p.\texttt{major}_S} ; & \mathrm{où} & & S = \{ & \texttt{major} : \tInt ; \\
  & \indent \uPut{14}{p.\texttt{minor}_S} ; &             & &        & \texttt{minor} : \tInt ; \\
  & \indent \uPut{15}{p.\texttt{patch}_S}   &             & &        & \texttt{patch} : \tInt   \\
  & \}                                      &             & &     \}
\end{align*}

% TODO le champ utilisateur ça passe bien ?

Le but est en

% TODO
\begin{center}\rule{3in}{0.4pt}\end{center}

On procède en trois étapes (figure~\ref{fig:qualif-changes}) :

\begin{figure}%{{{

\begin{align*}
\gramdef{Expressions}{e}
  { … }{}
  { \eTaint{e} }{Expression souillée}
  { \uGet{lv}{e}}{ Lecture depuis l'espace utilisateur}
  { \uPut{e}{e'}}{ Écriture dans l'espace utilisateur}
  {END}
\\
\gramdef{Contextes}{C}
  { … }{}
  { \eTaint{C}}{}
  {END}
\\
\gramdef{Chemins}{φ}
  { … }{}
  { \vTainted{φ} }{Valeur souillée}
  {END}
\\
\gramdef{Erreurs}{Ω}
  { … }{}
  { \serr{taint} }{Erreur de souillure}
  {END}
\end{align*}

\caption{Ajouts liés aux pointeurs utilisateurs}
\label{fig:qualif-changes}
% TODO définir les états et ms

\end{figure}%}}}

\begin{itemize}
\item
  tout d'abord, on ajoute une expression d'annotation sur les variables
  indiquant que celles-ci sont contrôlés par un utilisateur non privilégié,
  ainsi que des opérateurs de copie sûre.
\item
  ensuite, on étend l'ensemble des valeurs possibles pour les pointeurs
  à une valeur $\vTainted{φ}$ signifiant que l'objet pointé se situe en
  % TODO [E] ?
  espace utilisateur
\item
  enfin, on définit une nouvelle erreur $\serr{taint}$ produite par le
  déréférencement d'un pointeur ayant une telle valeur.
\end{itemize}

Pour adapter l'évaluation, plusieurs cas sont à rajouter. D'une part, la
présence de $\phx{\eTaint}$ dans une instruction consiste à ajouter un
$\phx{\vTainted}$ dans la valeur construite. Ceci ne peut être fait que dans le
cas où la valeur est un chemin $φ$, c'est-à-dire que la construction
$\phx{\eTaint}$ ne peut se faire que sur une expression de type pointeur.

% TODO passe globale c'est-a-dire & co (pas de ')

\begin{mathpar}
  \irule{Expr-Tainted}
    { }
    { \mm{m}{\eTaint{φ}}{m}{\vTainted{φ}} }
\end{mathpar}

D'autre part, une règle accède à la mémoire : \textsc{Exp-Lv}; pour
rappel :

\begin{mathpar}
  \semrule{Exp-Lv}
\end{mathpar}

Puisque la définition des chemins $φ$ a été changée, il est aussi nécessaire de
redéfinir la lentille $Φ$ utilisée ci-dessus (définition~\ref{def:acces-phi}).

% TODO [E] à détailler

On rajoute donc le cas :

\begin{align*}
Φ(\vTainted{φ}) =& \serr{taint}
\end{align*}

Pour accéder à ces valeurs, il faut utiliser les opérateurs $\phxx\uGet$ et
$\phxx\uPut$.

% TODO leur sémantique

\section{Insuffisance des types simples}

% TODO[E] c'est ok

Étant donné \langname augmenté de cette extension sémantique, on peut étendre
trivialement le système de types avec la règle suivante :

\begin{mathpar}
  \irule{Taint-Ignore}
    { Γ ⊢ e : t* }
    { Γ ⊢ \eTaint{e} : t* }
\end{mathpar}

Cette règle est compatible avec l'extension, sauf qu'elle introduit des termes
qui sont bien typables mais dont l'évaluation provoque une erreur autre que
$\serr{div}$, $\serr{array}$ ou $\serr{ptr}$, violant ainsi le
théorème~\ref{thm:progres}.

Par exemple, supposons que $x$ soit une variable globale entière, et posons $e =
* \eTaint {\& x}$.

$e$ est alors bien typée sous $Γ = x : \tInt$ :

\begin{mathpar}
  \irule{Lv-Deref}
    {
      \irule{Taint-Ignore}
        {
          \irule{Lv-Deref}
            {
              \irule{Lv-Var}
                { x:\tInt ∈ Γ }
                { Γ ⊢ x : \tInt }
            }
            { Γ ⊢ \& x : \tInt*}
        }
        { Γ ⊢ \eTaint {\& x} : \tInt*}
    }
    { Γ ⊢ * \eTaint {\& x} : \tInt}
\end{mathpar}

Posons alors $m = ([[x↦0]], [])$ (on a bien $\mcomp{Γ}{m}$). L'évaluation de $e$
sous $m$ provoque une erreur, comme le montre la dérivation suivante.

\begin{mathpar}
\inferrule*
  {
    \irule{Exp-Lv}
      {
        \inferrule*
          { }
          {m[*\vTainted{x}] = \serr{taint}}
      }
      { \mm{m}{* \eTaint {\& x}}{m}{\serr{taint}} }
    \\
    \irule{Eval-Err}
      { }
      {
        \msi{m}{\serr{taint}} → \serr{taint}
      }
  }
  {\msi{m}{* \eTaint {\& x}} → \serr{taint}}
\end{mathpar}

\section{Extensions du système de types}

On présente ici un système de types plus expressif permettant de capturer les
extensions de sémantique. \emph{In fine}, cela permettra de prouver le
théorème~\ref{thm:progres-qual} qui est l'équivalent du
théorème~\ref{thm:progres} mais pour le nouveau jugement de typage.

% TODO [E] ??

Définir un nouveau système de types revient à définir un nouveau jugement de
typage $\cdot ⊢_q \cdot : \cdot$, à partir d'un ensemble de règles. Pour la
plupart, les règles seront identiques, donc sauf mention contraire, les règles
portant sur $\cdot ⊢ \cdot : \cdot$ s'appliqueront aussi sur $\cdot ⊢_q \cdot :
\cdot$. Naturellement, la plupart des différences porteront sur le traitement des
pointeurs.

La différence principale est qu'à chaque pointeur, on ajoute un
\emph{qualificateur} qui représente \emph{qui} contrôle sa valeur
(section~\ref{sec:secu-syscalls}). Les deux qualificateurs sont :

\begin{itemize}
\item
  \qKernel : il s'applique aux pointeurs contrôlés par le noyau. Par exemple,
  prendre l'adresse d'un objet de la pile noyau donne un pointeur noyau.
\item
  \qUser : il s'applique aux pointeurs qui proviennent de l'espace utilisateur.
  Ces pointeurs proviennent toujours d'interfaces particulières, comme les
  appels système ou les paramètres de la fonction \texttt{ioctl}.
% TODO [E] tu comptes malloc ? &x ?
\end{itemize}

Cet ajout est précisé dans la figure~\ref{fig:qualif-changes-typ}.

\begin{figure}%{{{

\begin{align*}
\gramdef{Qualificateurs}{q}
  { \qKernel }{Donnée noyau (sûre)}
  { \qUser   }{Donnée utilisateur (non sûre)}
  {END}
\\
\gramdef{Types}{t}
  { … }{}
  { \msout{t~*} }{\textrm{Pointeur}}
  { t~q~* }{Pointeur qualifié}
  {END}
\end{align*}

\caption{Changements liés aux qualificateurs de types}
\label{fig:qualif-changes-typ}
\end{figure}%}}}

Au niveau du système de types, la principale restriction est que seuls les
pointeurs \qKernel peuvent être déréférencés de manière sûre :

\begin{mathpar}
\irule{Lv-Deref-Kernel}{
  Γ ⊢_q e : τ~\textsc{Kernel}*
}{
  Γ ⊢_q *e : τ
}
\end{mathpar}

% TODO [E] et que se passe-t'il quand on a une expression e de type τ K* que
% l'on cherche à déréférencer ?

L'opérateur $\eTaint{\cdot}$ transforme un pointeur selon la règle de souillure
suivante :

% TODO définir les stubs

\begin{mathpar}
\irule{Taint}
  { Γ ⊢_q e : t~q~* }
  { Γ ⊢_q \eTaint{x} : t~\qUser~* }
\end{mathpar}

Les opérateurs $\phxx\uGet$ et $\phxx\uPut$ sont typés de la manière suivante :

\begin{mathpar}
    \disprule{GetU}

    \disprule{PutU}
\end{mathpar}

% TODO adresse = kernel

% TODO parler de sous typage

La définition du typage sémantique doit aussi être étendue au cas $φ =
\vTainted{φ'}$. Les références mémoires sont ``nettoyées'' pour accéder à la
left-value encapsulée.

\begin{mathpar}
  \irule{S-Tainted}
    { m ⊧           φ'  : t~\qKernel~* }
    { m ⊧ \vTainted{φ'} : t~\qUser~* }
\end{mathpar}

\wip

\section{Sûreté du typage}

Le déréférencement d'un pointeur dont la valeur est contrôlée par l'utilisateur
ne peut se faire qu'à travers une fonction qui vérifie la sûreté de celui-ci.

\begin{theorem}[Progrès pour les types qualifiés]
  \label{thm:progres-qual}

  Supposons que $Γ ⊢_q e : t$. Soit $m$ un état mémoire tel que $\mcomp{Γ}{m}$.
  Alors l'un des cas suivant est vrai :

\begin{itemize}
  \item $∃ v ≠ Ω, e = v$
  \item $∃ (e', m'), \mcomp{Γ}{m'} ∧ \mm{m}{e}{m'}{e'}$
  \item $∃ Ω ∈ \{\serr{div},\serr{array},\serr{ptr}\}, \msi{m}{e} → Ω$
\end{itemize}
\end{theorem}

%\begin{proof}

  %% TODO

%\end{proof}

Et nous donnons un équivalent du théorème~\ref{thm:preservation}.

\begin{theorem}[Préservation pour les types qualifiés]

  Si une expression est typable et que son évaluation produit une valeur, alors
  cette valeur est du même type que l'expression.

  Si $Γ ⊢_q e : t$ et $e → v$ % TODO

  alors $Γ ⊢_q v : t$.

\end{theorem}

% vim: spelllang=fr


%\chapter{Analyse de terminaison des chaînes C}

%\label{cha:qualsz}
%\input{chap/06b-qualsz.tex}

\partccl{part:lang}
\label{cha:ccl2}
On vient de décrire en détail un langage impératif, \langname: tout d'abord sa
syntaxe et sa sémantique d'évaluation dans le chapitre~\ref{cha:lang}. Une des
particularités de cette sémantique est l'utilisation de lentilles pour modifier
les valeurs composées en profondeur.

\label{page:ccl-II-monad}
Il y a plusieurs alternatives à cette présentation. La première est la solution
classique qui consiste à décrire les modifications de la mémoire en extension.
C'est en général long et laborieux puisqu'il faut définir les accès en lecture
et écriture à chaque étape (avec des lentilles on décrit ces deux opérations
uniquement sur les briques du calcul, et la composition fait le reste). La
seconde solution est d'employer une sémantique monadique. Les transisitions sont
alors encodées comme des actions monadiques qui représentent les modifications
de la mémoire. Un des avantages de cette solution est qu'elle est très
extensible. Par exemple, la propagation des erreurs, ou l'ajout de continuations
(c'est-à-dire le support des fonctions \texttt{setjmp} et \texttt{longjmp})
peuvent facilement être exprimées dans un formalisme monadique. Nous avons
préféré une présentation classique qui reste plus accessible, et suffisante
compte tenu de la simplicité des constructions à interpréter dans le langage.

Ensuite, dans le chapitre~\ref{cha:typbase} nous avons ajouté un système de
types à \langname. Le but est de restreindre le genre d'erreurs qui peuvent
arriver lors de l'évaluation d'un programme. Par le théorème de progrès
(théorème~\ref{thm:progres} , page~\pageref{thm:progres}), on interdit les
erreurs qui signalent une manipulation de valeurs incompatibles, l'accès à un
champ de structure inconnu, et l'accès à une variable inexistante. Et le
théorème de préservation (théorème~\ref{thm:preservation},
page~\pageref{thm:preservation}) formalise le résultat classique qu'une étape
d'évaluation ne modifier pas le typage.
Une particularité de \langname est que son état mémoire est structuré, avec une
pile de variables locales explicite. On retrouve donc cette distinction dans le
typage: les variables globales et les variable locales sont séparées dans les
environnement de typage $Γ$ (page~\pageref{page:gamma-split}).

Enfin, on étend dans le chapitre~\ref{cha:qualifs} le langage pour exprimer la
notion de pointeurs utilisateur. Cela permet d'écrire des fonctions qui se
comportent comme des appels système. On commence par montrer qu'une extension
naïve du système de types ne suffit pas, car le théorème de progrès est alors
invalidé. On ajoute donc un type dédié aux pointeurs utilisateur. Les valeurs de
ce type sont créées explicitement et passés aux appels système. La règle de
typage du déréférencement est restreinte aux pointeurs noyau, ce qui permet de
ré-établir les théorèmes de progrès et préservation.

Notre technique de typage permet donc d'exprimer correctement les problèmes liés
à la manipulation mémoire lors des appels système, ainsi que décrits dans le
chapitre~\ref{cha:os}: c'est une méthode simple pour détecter et empêcher les
problèmes de sécurité qui proviennent des pointeurs utilisateur.

Comme nous l'avons fait remarquer dans le chapitre~\ref{cha:etatdelart},
utiliser une technique de typage pour étudier des propriétés sur les données a
déjà été explorée dans l'outil CQual~\cite{pldi99}, en particulier sur les
problèmes de pointeurs utilisateurs~\cite{cquk-usenix04}.

En effet, si on remplace «$\ptrK{t}$» par «$\qKernel~t~*$» et «$\ptrU{t}$»
par «$\qUser~t~*$», on obtient un début de système de types qualifiés.

En revanche, il y a une différence importante : CQual modifie fondamentalement
l'ensemble du système de types, pas \langname. Le jugement de typage de CQual a
pour forme générale $Γ ⊢ e : q~t$ (où $Γ$ est un environnement de typage, $e$
une expression, $q$ un qualificateur et $t$ un type), alors que le nôtre a la
forme plus classique $Γ ⊢ e : t$.

En intégrant $q$ à la relation de typage, on ajoute un qualificateur à chaque
type, même les expressions pour lesquels il n'est pas directement pertinent de
déterminer qui les contrôle (comme par exemple, un entier). Dans CQual, ceci
permet de traiter de manière correcte le transtypage. Par exemple, si $e$ a pour
type qualifié $\qUser~\tInt$, alors $(\tFloat~*)~e$ aura pour type qualifié
$\qUser~\tFloat~*$, et déréférencer cette expression produira une erreur de
typage. \langname, dans son état actuel, ne permet pas de traiter les
conversions de type et ne permet donc pas de traiter ce cas.

Nous prenons au contraire, l'approche de ne modifier le système de types que là
où cela est nécessaire, c'est-à-dire sur les types pointeurs. Cela permet de ne
pas avoir à modifier en profondeur un système de types existant.

Le modèle d'exécution est aussi très différent. CQual s'appuie sur un
langage proche de ML: un noyau de lambda-calcul avec des références. Le système
de types sous-jacent est proche de celui d'OCaml: du polymorphisme de premier
ordre (avec la restriction habituelle de généralisation des références) et du
sous-typage structurel. En outre, leur approche repose sur une gestion
automatique de la mémoire. De notre côté, nous nous appuyons sur un modèle
mémoire plus proche de C, reposant sur une pile de variables, et une gestion
manuelle de la mémoire.

Une autre différence fondamentale est que le système de types de CQual fait
intervenir une relation de sous-typage. Le cas particulier du problème de
déréférencement des pointeurs utilisateurs peut être traité dans ce cadre en
posant $\qKernel \preceq \qUser$ pour restreindre certaines opérations aux
pointeurs $\qKernel$.

Notre approche, au contraire, n'utilise pas de sous-typage, mais consiste à
définir un type abstrait $\ptrU{t}$ partageant certaines propriétés avec
$\ptrK{t}$ (comme la taille et la représentation) mais incompatible avec
certaines opérations. C'est à rapprocher des types abstraits dans les langages
comme OCaml et Haskell.

Les perspectives de travaux futurs sont également très différentes. Dans le cas
des pointeurs, même si le noyau Linux (et la plupart des systèmes
d'exploitation) ne comportent que deux espaces d'adressage, il est commun dans
les systèmes embarqués de manipuler des pointeurs provenant d'espaces mémoire
indépendants : par exemple, de la de la mémoire flash, de la RAM, ou une EEPROM
de configuration. Ces différentes mémoires possèdent des adresses, et un
pointeur est interprété comme faisant référence à une ou l'autre selon le code
dont il est tiré.

% vim: spelllang=fr


\long\def\partintro{
% TODO contexte!!
On décrit ici la démarche expérimentale liée à l'implémentation des analyses
décrites dans la partie~\ref{part:lang}.

Le chapitre~\ref{cha:implem} décrit l'implémentation en elle-même : comment le
code source C est compilé vers \langname, et comment les types du programme sont
vérifiés.

Ensuite, dans le chapitre~\ref{cha:etudedecas}, le cas d'un \emph{bug} de pilote
graphique dans le noyau Linux est étudié. On montre que les analyses précédentes
permettent de distinguer statiquement entre le cas incorrect et le cas corrigé.

Enfin, le chapitre~\ref{cha:conclusion} conclut par un résumé des contributions
présentées ici, les limitations de cette approche, ainsi que les perspectives de
travaux futurs.

% TODO un peu court

}

\part{Expérimentation}
\label{part:xp}


\chapter{Implantation}

\label{cha:implem}
Dans ce chapitre, nous décrivons la mise en œuvre des analyses statiques
précédentes. Nous commençons par décrire la représentation intermédiaire
utilisée: Newspeak. La chaîne de compilation est explicitée, partant de C pour
aller jusqu'à Tyspeak, sur lequel nous décrivons d'un algorithme d'inférence de
types à la Hindley-Milner, reposant sur l'unification et le partage de
références. Toutes ces étapes sont implantées dans le langage
OCaml~\link{ocaml}.

En réalité, Newspeak (et donc Tyspeak) utilise une représentation de plus bas
niveau que \langname donc les règles sont légèrement différentes.

% TODO mieux décrire les 3

\section{Newspeak}
\label{sec:npk}

Newspeak est un langage intermédiaire conçu pour être un bon support d'analyses
statiques, contrairement à des langages conçus pour les programmeurs comme C.
Plusieurs points en font un langage pratique à utiliser. Tout d'abord, il est
suffisamment expressif pour compiler tout C, y compris de nombreuses extensions
GNU utilisées dans le noyau Linux (section~\label{sec:gnuc}), et des traducteurs
automatiques depuis C et Ada existent (section~\label{sec:compil}).

Ensuite, ses instructions sont orthogonales et minimales: il existe en général
une seule manière de faire les choses. Par exemple, le flot de contrôle est
restreint à la boucle infinie et au saut en avant («\emph{break}» généralisé).

Enfin, lorsque certaines constructions sont ambigües, un choix est fait. Par
exemple, l'évaluation des arguments d'une fonction est faite de manière
explicite, les tailles des types sont indiqués à chaque déclaration de variable,
etc. Sa sémantique d'exécution (ainsi qu'une partie des étapes de compilation
depuis C) est décrites dans \cite{newspeak}. Sa syntaxe est donnée dans la
figure~\ref{fig:stx-npk}.

\begin{figure}

\def\npkstyle#1{\mathrm{#1}}

\begin{align*}
\gramdef{Instruction}{s}
    {\npkstyle{Set} (lv, e, st}{Affectation}
    {\npkstyle{Copy} (lv, lv, n)}{Copie}
    {\npkstyle{Guard} (e)}{Garde}
    {\npkstyle{Decl} (var, t, blk)}{Déclaration}
    {\npkstyle{Select} (blk, blk)}{Branchement}
    {\npkstyle{InfLoop} (blk)}{Boucle infinie}
    {\npkstyle{DoWith} (blk, x)}{Nommage de bloc}
    {\npkstyle{Goto} (x)}{Saut}
    {\npkstyle{Call} ([(e_i, t_i)], f, [(lv_i, t_i)])}
            {Appel de fonction}
    {END}
\\ \\
\gramdef{Bloc}{blk}
    {[i_i]}{Liste d'instructions}
    {END}
\\ \\
\gramdef{Left-value}{lv}
    {\npkstyle{Local} (x)}{Locale}
    {\npkstyle{Global} (x)}{Globale}
    {\npkstyle{Deref} (e, n)}{Déréférencement}
    {\npkstyle{Shift} (lv, e)}{Décalage}
    {END}
\\ \\
\gramdef{Expression}{e}
    { \npkstyle{CInt} (n) }{Entier}
    { \npkstyle{CFloat} (d)}{Flottant}
    { \npkstyle{Nil} }{ Pointeur nul }
    { \npkstyle{Lval} (lv, t)}{Accès mémoire}
    { \npkstyle{AddrOf} (lv) }{Adresse de variable}
    { \npkstyle{AddrOfFun} (x, [t_i], [t_i])}{Adresse de fonction}
    { \npkstyle{UnOp} (unop, e) }{Opérateur unaire}
    { \npkstyle{BinOp} (binop, e_1, e_2) }{Opérateur binaire}
    {END}
\\ \\
\gramdef{Fonction}{f}
    { \npkstyle{FunId} (x)}{Appel par nom}
    { \npkstyle{FunDeref} (e)}{Appel par pointeur}
    {END}
\\ \\
\gramdef{Type}{t}
    { \npkstyle{Scalar} (st) }{Type scalaire}
    { \npkstyle{Array} (t, n)}{Tableau}
    { \npkstyle{Region} ([(n_i, t_i)], n')}{Structure/union}
    {END}
\\ \\
\gramdef{Type scalaire}{st}
    { \npkstyle{Int} (n) }{Entier}
    { \npkstyle{Float} (n) }{Flottant}
    { \npkstyle{Ptr} }{ Pointeur sur données}
    { \npkstyle{FunPtr}}{ Pointeur sur fonction}
    {END}
\end{align*}
\caption{Syntaxe simplifiée de Newspeak}
\label{fig:stx-npk}
\end{figure}

\section{Chaîne de compilation}
\label{sec:compil}

\begin{figure}
  \centering
  \begin{tikzpicture}\shorthandoff{!}
\tikzstyle{file}=[draw, shape=rectangle, node distance=2.2cm, minimum
height=1cm, shade, top color=white,
    bottom color=blue!50!black!20, draw=blue!40!black!60, very thick];

\node [file] (c1) {\textcolor{black}{.c}};
\node [file, below of=c1] (c2) {\textcolor{black}{.c}};
\node [node distance=2.2cm, below of=c2] (c3) {};

\node [file, minimum height=0, node distance=5mm, above of=c3,draw] (c3b) {.adb};
\node [file, minimum height=0, node distance=5mm, below of=c3,draw] (c3s) {.ads};

\path (c3b.north west) ++(-3mm,3mm) [draw,dotted] rectangle ($(c3s.south east)+(3mm,-3mm)$);

\node [below of=c3, node distance=2.2cm] (c4){};

\node [file, right of=c1] (cc1) {\textcolor{black}{.c}};
\node [file, right of=c2] (cc2) {\textcolor{black}{.c}};
\node [node distance=2.2cm, right of=c3] (cc3) {};

\node [below of=cc3, node distance=2.2cm](cc4){};

\draw[->] (c1) -- node[above] {{\tiny \ttfamily gcc -E}} (cc1);
\draw[->] (c2) -- (cc2);

\node [file, right of=cc1] (no1) {\textcolor{black}{.no}};
\node [file, right of=cc2] (no2) {\textcolor{black}{.no}};
\node [file, right of=cc3] (no3) {\textcolor{black}{.no}};

\node [below of=no3, node distance=2.2cm](no4){};

\draw[->] (cc1) -- node[above] {{\tiny \ttfamily c2newspeak -c}} (no1);
\draw[->] (cc2) --  (no2);

\draw[->] ($ (c3b.north east)!0.5!(c3s.south east) + (3mm,0) $) -- node[above] {\tiny \ttfamily ada2newspeak -c} (no3);


\node [file, right of=no2] (npk) {\textcolor{black}{.npk}};

\node [right of=no3, node distance=2.2cm](npk2){};
\node [right of=no4, node distance=2.2cm](npk3){};

\node[draw, ellipse, right of=no2, minimum height=3cm]{};

\draw[->] (no1) -- node {} (npk);
\draw[->] (no2) -- node[above] {{\tiny \ttfamily c2newspeak}} (npk);
\draw[->] (no3) -- node {} (npk);


\node [file, right of=npk](warn) {\textcolor{black}{.warn}};

\node [right of=npk3, node distance=3cm](analyser2){};

\draw[->] (npk) -- node[above] {{\tiny \ttfamily penjili}} (warn);

\draw[->] (c4)   -- node[above, text depth=3pt] {\footnotesize Pr\'etraitement} (cc4);
\draw[->] (cc4)  -- node[above, text depth=3pt] {\footnotesize Compilation} (no4);
\draw[->] (no4)  -- node[above, text depth=3pt] {\footnotesize \'Edition de liens} (npk3);
\draw[->] (npk3) -- node[above, text depth=3pt] {\footnotesize Analyse} (analyser2);
\end{tikzpicture}

  \caption{Compilation depuis Newspeak}
\label{fig:compil-npk}
\end{figure}

%TODO Mettre à jour la figure

La compilation vers C est faite en trois étapes (figure~\ref{fig:compil-npk}):
prétraitement du code source, compilation de C prétraité vers \newspeak{}, puis
compilation de \newspeak{} vers ce langage.

La première étape consiste à prétraiter les fichiers C source avec le logiciel
\texttt{cpp}, comme pour une compilation normale. Cette étape interprète les
directives comme \texttt{\#include}, \texttt{\#ifdef}, À cet étape, les
commentaires sont aussi supprimés.

Une fois cette passe effectuée, le résultats est un ensemble de fichiers C
prétraités, c'est-à-dire des unités de compilation.

Puisque la directive \texttt{\#include} est textuelle, ces fichiers sont très
grands et donnent lieu à beaucoup de duplication dans les passes suivantes.

Sur cette représentation (du C prétraité), il est possible d'ajouter des
annotations de la forme \texttt{/*!npk [...] */} qui pourront être accessible
dans l'arbre de syntaxe abstrait des passes suivantes.

À ce niveau, les fichiers sont passés à l'outil \texttt{c2newspeak} qui les
traduit vers Newspeak. Comme il sera décrit dans la section~\ref{sec:gnuc}, la
plupart des extensions GNU C sont acceptés en plus du C ANSI. Dans cette étape,
les types et les noms sont résolus, et le programme est annoté de manière à
rendre les prochaines étapes indépendante du contexte. Par exemple, chaque
déclaration de variable est adjointe d'une description complète du type.

Lors de cette étape, le flôt de contrôle est également simplifié. C en effet
propose de nombreuses constructions ambigües ou redondantes.

Au contraire, Newspeak propose un nombre réduit de constructions. Rappelons que
le but de ce langage est de faciliter l'analyse statique: des constructions
orthogonales permettent donc d'éviter la duplication de règles sémantique, ou de
code lors de l'implémentation d'un analyseur.

Par exemple, plutôt que de fournir une boucle \emph{while}, une boucle
\emph{do/while} et une boucle for, Newspeak fournit une unique boucle
\npkWhile{}. La sortie de boucle est compilée vers un \npkGoto{}\cite{goto}, qui
est toujours un saut vers l'avant (similaire à un «\emph{break}» généralisé).

Newspeak est conçu pour l'analyse statique par interprétation abstraite. Il a
donc une vue de bas niveau sur les programmes. Par exemple, aucune distinction
n'est faite entre l'accès à un champ et l'accès à un l'élément d'un tableau
(tous deux sont traduits par un décalage numérique depuis le début de la zone
mémoire). Pour supprimer cette ambiguïté, il faut s'interfacer dans les
structures internes de \texttt{c2newspeak}, où les informations nécessaires sont
encore présentes.

Ensuite, les différents fichiers sont liés ensemble. Cette étape consiste
principalement à s'assurer que les hypothèses faites par les différentes unités
de compilation sont cohérentes entre elles. Les objets marqués \texttt{static},
invisibles à l'extérieur de leur unité de compilation, sont renommés afin qu'ils
aient un nom unique.

Enfin, l'implantation d'un algorithme d'inférence pour les systèmes de types
décrits dans les chapitres~\ref{cha:typbase} et~\ref{cha:qualifs} est assez
simple. Puisqu'ils sont suffisament proches du lambda calcul simplement typé, on
peut utiliser une variante de l'algorithme W de Damas et
Milner~\cite{DamasMilner}. On utilise l'optimisation classique qui consiste à se
reposer sur le partage de références pour réaliser l'unification, plutôt que de
faire des substitutions explicites. Puisque ces systèmes de types sont
monomorphes, on présente une erreur si des variable de type libres sont
présentes.

À la fin de cette étape, on obtient soit un programme complètement annoté, soit
une erreur de type.

\begin{figure} % fig:exunif:c {{{

  \insertcode{ex-unif-c.c}

  \caption{Exemple de programme C à compiler}
\label{fig:exunif:c}
\end{figure} % }}}

Prenons l'exemple de la figure~\ref{fig:exunif:c} et typons-le «à la main». On
commence par oublier toutes les étiquettes de type présentes dans le programme.
Celui-ci devient alors:

\insertcode{exunif-code.c}

La premiere ligne introduit deux variables. On peut noter leurs types respectifs
(inconnus pour le moment) $t_1$ et $t_2$. La première affectation \texttt{p =
\&x} implique que les deux côtés du signe «$=$» ont le même type. À gauche, le
type est $t_2$, et à droite $Ptr(t_1)$. On applique le même raisonnement à la
seconde affectation: à gauche, le type est $t_1$ et à droite \texttt{Int}. On en
déduit que le type de x est \texttt{Int} et celui de p est \verb!Ptr(Int)!.

\insertcode{lambda-types.ml}

% TODO c'est le code d'un typeur naif

Pour implanter cet algorithme, on représente les types de données du programmes
à typer par une valeur de type \texttt{ml\_type}. En plus des constantes de
types comme int ou float, et des constructeurs de type comme pair et fun, le
constructeur Var permet d'exprimer les variables de types (inconnues ou non).

Celles-ci sont numérotées par un int, on suppose avoir à disposition deux
fonctions manipulant un compteur global d'inconnues.

\insertcode{exunif-counter.ml}

De plus, un module gère les environnements de typage. Il peut être implanté avec
des listes d'association ou des tables de hachage, par exemple. Sa signature
est:

\insertcode{exunif-env.ml}

Reprenons l'exemple précédent. Partant d'un environnement vide (Env.empty), on
commence par l'étendre de deux variables. Comme on n'a aucune information, il
fait allouer des nouveaux noms d'inconnues (qui correspondent à $t_1$ et $t_2$):

\insertcode{exunif-1.ml}

La première instruction indique que les deux côtés de l'affectation doivent
avoir le même type.

\insertcode{exunif-2.ml}

\begin{figure}
  \centering
  \begin{tikzpicture}
    [scale=2.3]

    \node at (0,  0) (n0a) {t1:};
    \node at (0, -1) (n0b) {t2:};

    \node at (1,  0) (n1a) {Var\_type};
    \node at (2,  0) (n2a) {ref};
    \node at (3,  0) (n3a) {Unknown 0};

    \node at (1, -1) (n1b) {Var\_type};
    \node at (2, -1) (n2b) {ref};
    \node at (3, -1) (n3b) {Instanciated};
    \node at (4, -1) (n4b) {Ptr};

    \draw [->] (n1a) -- (n2a);
    \draw [->] (n2a) -- (n3a);

    \draw [->] (n1b) -- (n2b);
    \draw [->] (n2b) -- (n3b);
    \draw [->] (n3b) -- (n4b);

    \draw [->] (n4b) -- ++ (0, 0.5) -| (n1a);

  \end{tikzpicture}
  \caption{Partage lors de l'unification}
\label{fig:unifsharing}
\end{figure}

Ici il se passe plusieurs choses intéréssantes. D'une part nous faisont appel à
une fonction externe typeof qui retourne le type d'une expression sous un
environnement (dans une implantation complète il s'agirait d'un appel récursif).
Dans ce cas, \texttt{typeof lhs1 env} est identique à \texttt{Env.get lhs1 env}
et \texttt{typeof rhs1 env} à \texttt{Ptr\_type t1}. L'autre aspect intéressant
est la dernière ligne: la fonction \texttt{unify} va modifier en place les
représentations des types afin de les rendre égales. L'implantation de
\texttt{unify} sera décrite plus tard. Dans ce cas précis, elle va faire pointer
la référence dans t2 vers t1 (figure~\ref{fig:unifsharing}).

Enfin, la seconde affectation se déroule à peu près de la même manière.

\insertcode{exunif-3.ml}

Ici \texttt{typeof lhs2 env} est identique à
\verb!Ptr_type (Env.get "p" env)! et \texttt{typeof lhs2 env} à \texttt{Const\_type Int\_type}. Et dans cas,
l'unification doit se faire entre t1 et \texttt{Const\_type Int\_type}: cela
mute la référence derrière t1 (figure~\ref{fig:typeunifref}).

\begin{figure}
  \begin{tikzpicture}
    [scale=2.3]

    \node at (0,  0) (n0a) {t1:};
    \node at (0, -1) (n0b) {t2:};

    \node at (1,  0) (n1a) {Var\_type};
    \node at (2,  0) (n2a) {ref};
    \node at (3,  0) (n3a) {Instanciated};
    \node at (4,  0) (n4a) {Const\_type};
    \node at (5,  0) (n5a) {Int\_type};

    \node at (1, -1) (n1b) {Var\_type};
    \node at (2, -1) (n2b) {ref};
    \node at (3, -1) (n3b) {Instanciated};
    \node at (4, -1) (n4b) {Ptr};

    \draw [->] (n1a) -- (n2a);
    \draw [->] (n2a) -- (n3a);
    \draw [->] (n3a) -- (n4a);
    \draw [->] (n4a) -- (n5a);

    \draw [->] (n1b) -- (n2b);
    \draw [->] (n2b) -- (n3b);
    \draw [->] (n3b) -- (n4b);

    \draw [->] (n4b) -- ++ (0, 0.5) -| (n1a);

  \end{tikzpicture}
  \caption{Unification par mutation de références}
\label{fig:typeunifref}
\end{figure}

L'essence de l'algorithme d'inférence se situe donc dans 2 fonctions. D'une
part, \texttt{unify} qui réalise l'unification des types grâce à au partage des
références. D'autre part, la \texttt{typeof} qui encode les règles de typage
elles-mêmes et les applique à l'aide de \texttt{unify}.

\subsection*{Algorithme d'unification}

Voici une implantation de la fonction \texttt{unify}.

Celle-ci prend en entrée deux types $t_1$ et $t_2$. À l'issue de l'exécution de
\texttt{unify}, ces deux types doivent pouvoir être considérés comme égaux. Si
ce n'est pas possible, une erreur sera levée.

La première étape est de réduire ces deux types, c'est-à-dire à transformer les
constructions \texttt{Var (ref (Instanciated t))} en \texttt{t}.

Ensuite, cela dépend des formes qu'ont les types réduits:

\begin{figure}
  \centering
  \begin{tikzpicture}
    [every node/.style={node distance=2cm}
    ,>=latex
    ]

    \node (n1) {Ptr};
    \node[right of=n1] (n2) {Var};
    \node[right of=n2] (n3) {ref};
    \node[right of=n3] (n4) {Instanciated};

    \draw [->] (n1) -- (n2);
    \draw [->] (n2) -- (n3);
    \draw [->] (n3) -- (n4);

    \draw [->] (n4) -- ++ (0, -1cm) -| (n1);

  \end{tikzpicture}
  \caption{Cycle dans le graphe de types}
\label{fig:typecycle}
\end{figure}

\begin{itemize}

\item si les deux types sont inconnus (de la forme \texttt{Var (ref
(Instanciated t))}), on fait pointer l'une des deux références vers le premier
type. Notons que cela crée un type de la forme \texttt{Var (ref (Instanciated
(Var (ref (Unknown n)))))} qui sera réduit lors d'une prochaine étape
d'unification.

\item si un type est inconnu et pas l'autre, il faut de la même manière affecter la
référence. Mais en faisant ça inconditionnellement, cela peut poser problème:
par exemple en tentant d'unifier \texttt{a} avec \verb!Ptr(a)! on pourrait
créer un cycle dans le graphe (figure~\ref{fig:typecycle}).
Pour éviter cette situation, il suffit de s'assurer que le type inconnu n'est
pas présent dans le type à affecter.

\item si les deux types sont des types de base (comme $\tInt$ ou $\tFloat$)
égaux, on ne fait rien.

\item si les deux types sont des constructeurs de type, il faut que les
constructeurs soient égaux. On unifie en outre leurs arguments deux à deux.

\item si les deux types sont des types structure, on étudie les types de leurs
champs. On commence par unifier les types des champs de même nom, et on
on ajoute à l'un les champs exclusifs à l'autre.

\item dans les autres cas, l'algorithme échoue.

\end{itemize}

\section{Architecture de \ptrtype}
\label{sec:ptrtype-archi}

L'outil \ptrtype{} lit un programme Newspeak (ou un fichier C), et réalise
l'inférence de qualificateurs. En sortie, il affiche soit le programme
complètement annoté, soit une erreur. Le cœur de l'outil est dans la

Si l'argument à \ptrtype{} est un fichier C, il est tout d'abord compilé vers
Newspeak grâce à l'utilitaire \texttt{c2newspeak}. Ensuite, les autres passes
travaillent sur une représentation intermédiaire proche de Newspeak, mais où des
étiquettes de type supplémentaires sont ajoutées. Ce type de représentation
intermédiaire (polymorphe en le type des étiquettes) est \texttt{'a Tyspeak.t}.

\begin{figure}
\insertcode{implem-process.ml}
\caption{Fonction principale de \ptrtype{}}
\label{fig:implem-process}
\end{figure}

Le reste de l'outil est résumé dans la fonction
\texttt{process\_npk} (figure~\ref{fig:implem-process}):

\begin{itemize}

\item Grâce à la fonction \verb!convert_unit : Newspeak.t -> unit Tyspeak.t!,
  on ajoute des étiquettes «vides» (toutes égales à \verb!() : unit!).

\item L'ensemble des fonctions du programme est trié topologiquement selon la
  relation~$\preceq$ définie par $f \preceq g \eqdef \textrm{«} g
  \textrm{ apparaît dans la définition de } f \textrm{»}$. Cela est fait en
  construisant une représentation de $\preceq$ sous forme de graphe, puis en
  faisant un parcours en largeur de celui-ci.

% TODO et pour les fonctions mutuellement récursives?

\item Les annotations extérieures sont alors lues (variable \texttt{exttbl}), ce
  qui permet de créer un environnement initial.

\item Les types de chaque fonction sont inférés, par le biais de la fonction
  suivante:

\begin{Verbatim}
val infer : Newspeak.fid list
         -> Types.simple Env.t
         -> 'a Tyspeak.t
         -> Types.simple Tyspeak.t
\end{Verbatim}

\item Le programme obtenu, de type \texttt{Types.simple Tyspeak.t}, est affiché
  sur le terminal.

\end{itemize}

\section{Inférence de types}

L'inférence de types consiste à remplacer les étiquettes de type \texttt{unit}
par des étiquettes de type \texttt{simple} (autrement dit de vraies
représentations de types).

Cette étape se fait de manière impérative: on peut créer de nouveaux types avec
\texttt{new\_unknown} et unifier deux types
avec \texttt{unify}. Leurs types sont:

\begin{verbatim}
val new_unknown : unit -> Types.simple
val unify : Types.simple -> Types.simple -> unit
\end{verbatim}

La fonction \texttt{infer} s'appuie sur un ensemble de fonctions récursivement
définies portant sur chaque type de fragment: \texttt{infer\_fdec} pour les
déclarations de fonction, \texttt{infer\_exp} pour les expressions,
\texttt{infer\_stmtkind} pour les instructions, etc.

Les règles de typage sont implantées par \texttt{new\_unknown} et
\texttt{unify}. Par exemple, pour typer une déclaration
(figure~\ref{fig:implem-unif-stmt}), on crée un nouveau
type \texttt{t0}. On étend l'environnement courant avec cette nouvelle
association et sous ce nouvel environnement, on type le bloc de portée de la
déclaration.

\begin{figure}[h] %{{{

\insertcode{infer-unif-stmt.ml}

\caption{Inférence des déclarations de variable et appels de
         fonction}

\label{fig:implem-unif-stmt}
\end{figure}%}}}

De même, pour typer un appel de fonction, on infère le type de ses arguments et
left-values de retour. On obtient également le type de la fonction (à partir du
type de la fonction présent dans l'environnement, ou du type du pointeur de
fonction qui est déréférencé), et on unifie ces deux informations.

On peut donner quelques exemples, comme:

\begin{itemize}
  \item addition de deux flottants (dans \texttt{infer\_binop}):

\begin{verbatim}
let infer_binop op (_, a) (_, b) =
  match op with
    (* [...] *)
    | N.PlusF _ ->
        unify a Float;
        unify b Float;
        Float
\end{verbatim}

  \item adresse d'une left-value (dans \texttt{lval\_type}):

\begin{verbatim}
| T.AddrOf lv ->
    let lv' = infer_lv env lv in
    let ty = lval_type env lv in
    (T.AddrOf lv', Ptr (Kernel, ty))
\end{verbatim}

  \item déréférencement d'une left-value (dans \texttt{lval\_type}):

\begin{verbatim}
| T.Deref(e, _sz) ->
    let (_, te) = infer_exp env e in
    let t = new_unknown () in
    unify (Ptr (Kernel, t)) te;
    t
\end{verbatim}

\end{itemize}

\section{Unification}

\begin{figure}
\insertcode{implem-types.ml}
\caption{Représentation des types}
\label{fig:implem-typ}
\end{figure}

Les types de données utilisés sont donnés dans la figure~\ref{fig:implem-typ}.
Les types \langname{} sont représentés soit par un constructeur de type
«résolu» immutable comme \texttt{Int}, soit par une référence dans le cas
d'une variable inconnue (placée alors derrière le constructeur \texttt{Var}).

Ces références contiennent une valeur du type
\texttt{simple var\_type}\footnote{
  Ce type est polymorphe pour pouvoir être repris dans l'unification des
  qualificateurs (figure~\ref{fig:implem-qualifs}).
}, c'est-à-dire:

\begin{itemize}
\item soit un numéro d'inconnue (constructeur \texttt{Unknown})\footnote{On
  place le numéro dans un enregistrement pour abstraire l'entier sous-jacent, en
  empêchant par exemple de faire de l'arithmétique sur celui-ci.}.

\item soit un type résolu (constructeur \texttt{Instanciated}) si cette
  inconnue a été unifiée avec un type concret.
\end{itemize}

Ce système peut créer des représentations de types arbitrairement longues, comme
par exemple:

\texttt{Var (ref (Instanciated (Var (ref (Instanciated Int)))))}

Cela est dû au fait que \texttt{fun x -> Var (ref (Instanciated x))} est typée
\texttt{simple -> simple} et peut donc être appliquée à loisir. Pour éviter cet
effet d'allongement, on définit une fonction \texttt{shorten} qui supprime ces
chaînes (figure~\ref{fig:implem-shorten}).

\begin{figure}

  \insertcode{implem-shorten.ml}

\caption{Fonction de racourcissement des représentations de types}
\label{fig:implem-shorten}
\end{figure}

La fonction d'unification, quant à elle, commence par raccourcir ses deux
arguments puis faire une analyse de cas par filtrage.

Les cas principaux pour unifier deux types réduits \texttt{sta} et \texttt{stb}
sont:

\begin{itemize}

\item \texttt{sta} et \texttt{stb} sont deux inconnues. Alors on modifie l'un
  pour pointer sur l'autre. Les références étant uniques et partagées, cela
  revient à substituer l'un par l'autre dans toutes les représentations de
  types.

\item Un type est concret, l'autre une inconnue. Dans ce cas on modifie le
  second comme étant un \texttt{Instanciated} du premier. Il faut vérifier que
  le type inconnu remplacé n'apparaît pas dans le type concret, sinon on créé un
  cycle. Ce cas a lieu par exemple quand on cherche à unifier $t$ avec
  $\ptrK{t}$. Il faut alors signaler une erreur, ce que fait la fonction
  \texttt{occurs\_check\_failed}.

\item Les deux types sont des types concrets. Alors ils sont de la forme
  respective
  $C (t_1, …, t_n)$
  et
  $D (u_1, …, u_m)$ où $C$ et $D$ sont des constructeurs de type avec
  respectivement $n$ et $m$ arguments\footnote{
    Le cas $n = 0$ ou $m = 0$ correspond aux types concrets comme $\tInt$ ou
    $\tFloat$.
  }. Si $C = D$ et
  $n = m$, alors on unifie récursivement chaque $t_i$ avec $u_i$. Sinon on lève
  une erreur (fonction \texttt{type\_clash}).

\end{itemize}

Cette analyse de cas est implantée dans la fonction \texttt{unify_now}
(figure~\ref{fig:implem-unify-now}).

\begin{figure}
  \insertcode{implem-unify.ml}
\caption{Fonction d'unification}
\label{fig:implem-unify-now}
\end{figure}

Pour les pointeurs, il est nécessaire de définir des fonctions similaires sur
les qualificateurs (figure~\ref{fig:implem-qualifs}).

\begin{figure}

  \insertcode{implem-qualifs.ml}

  \caption{Qualificateurs}
\label{fig:implem-qualifs}

\end{figure}

\begin{figure}

\insertcode{unify-main.ml}

\caption{Unification directe ou retardée}
\label{fig:implem-lazy}
\end{figure}

La fonction appellée directement par le reste du code, appellée \texttt{unify},
peut retarder l'unification (figure~\ref{fig:implem-lazy}). Dans ce cas, la
paire de types à unifier est mise dans une liste d'attente qui sera unifiée
après le parcours du programme. Le but est d'instrumenter l'inférence de types
afin de pouvoir en faire une exécution «pas à pas».

\section{Exemple}

Lançons l'analyse sur un petit exemple:

\begin{verbatim}
int f(int *x) { return (*x + 1); }
\end{verbatim}

L'exécution de notre analyseur affiche un programme complètement annoté:

\begin{verbatim}
 % ptrtype example.c
f : (Ptr (Int)) -> (Int)
Int (example.c:1#4)^f(Ptr (Int) x) {
  (.c:3#4)^!return =(int32)
    (coerce[-2147483648,2147483647]
      ( ( ([(x_Ptr (Int) : Ptr (Int))]32_Int
            : Int
          )
          + (1 : Int)
        ) : Int
      ) : Int
    );
}
\end{verbatim}

L'opérateur \texttt{coerce[a,b]} est un artefact de Newspeak, destiné à détecter
les débordements d'entiers lors d'une analyse de valeurs par interprétation
abstraite. Dans le cas de notre analyse, les valeurs ne sont pas pertinentes et
cet opérateur peut être vu que comme un «plus» unaire typé $(\tInt) → \tInt$.

Un exemple de cas d'erreur sera décrit dans la section~\ref{sec:demo-unif}.

% TODO pas de conclusion

% TODO trop de figures

% vim: spelllang=fr


\chapter{Étude de cas : un pilote de carte graphique}

\label{cha:etudedecas}
Dans ce chapitre, un exemple de mise en œuvre du système de type décrit dans le
chapitre~\ref{cha:qualifs} et implanté dans le chapitre~\ref{cha:implem}.

\section{Linux}

Le noyau Linux, abordé dans le chapitre~\ref{cha:os}, est un noyau de système
d'exploitation développé depuis le début des années 90 et ``figure de proue'' du
mouvement open-source. Au départ écrit par Linus Torvalds sur son ordinateur
personnel, il a au fil des années était porté sur de nombreuses architectures et
s'est enrichi de nombreux pilotes de périphériques. En 20XX, son code source
comporte 8.8 millions de lignes de code (en grande majorité du C) dont 55\% de
pilotes.

% TODO fix le XX
% TODO mettre à jour

Même si le noyau est monolithique (la majeure partie des traitements s'effectue
au sein d'un même fichier objet), les sous-systèmes sont indépendants. C'est ce
qui permet d'écrire des pilotes de périphériques et des modules.

\section{GNU C}

Linux est écrit dans le langage C, mais pas sa version normalisée. Il utilise le
dialecte GNU C qui est celui que supporte GCC.

Pour traduire ce dialecte, il a été nécessaire d'adapter \texttt{c2newspeak}. La
principale particularité est la notation \texttt{\_\_attribute\_\_((...))} qui
peut décorer les déclarations de fonctions, de variables ou de types. De
nouvelles fonctionnalités sont aussi présentes.

Par exemple, il est possible de manipuler des étiquettes de première classe : si
\texttt{lbl:} est présent avant une instruction, on peut capturer l'adresse de
celle-ci avec \texttt{void *p = \&\&lbl} et y sauter indirectement avec
\texttt{goto *p}.

Une autre fonctionnalité est le concept d'instruction-expression :
\texttt{(\{bloc\})}) est une expression, dont la valeur est celle de la dernière
expression évaluée lors de \texttt{i}.

Les attributs, quant à eux, rentrent dans trois catégories :

\begin{itemize}
  \item les annotations de compilation ; par exemple, \texttt{used} désactive
    l'avertissement ``cette variable n'est pas utilisée ''.

  \item les optimisations ; par exemple, les objets marqués \texttt{hot} sont
    groupés de telle manière qu'ils se retrouvent en cache ensemble.

  \item les annotations de bas niveau ; par exemple, \texttt{aligned(n)}
    spécifie qu'un objet doit être aligné sur au moins \texttt{n} bits.
\end{itemize}

Dans notre cas, toutes ces annotations peuvent être ignorées.

\section{Configuration}

Pour que le code noyau soit compilable, il est nécessaire de définir certaines
macros. En particulier, le système de configuration de Linux utilise des macros
nommées \texttt{CONFIG\_*} pour inclure ou non certaines fonctionnalités. Il a
donc fallu faire un choix ; nous avons choisi la configuration par défaut. Pour
analyser des morceaux plus importants du noyau, il faudrait définir un fichier
de configuration plus important.

\section{Bug}

On décrit le cas d'un pilote video qui contenait un bug de pointeur utilisateur.
Il est répertorié sur \url{http://freedesktop.org} en tant que bug \#29340. 

Pour changer de mode graphique, les pilotes de GPU peuvent supporter le
\emph{Kernel Mode Setting} (KMS).

Pour configurer un périphérique, l'utilisateur communique avec le pilote noyau
avec le mécanisme d'\emph{ioctls}\footnote{
  Ce nom vient de la fonction \texttt{ioctl()} pour
  \emph{Input/Output Control}.
}.
Ils sont similaires à des appels système, mais spécifique à un périphérique
particulier. Les fonctions implantant un \emph{ioctl} sont donc vulnérables à la
même classe d'attaques que les appels système, et donc doivent être écrits avec
une attention particulière.

Le code suivant est présent dans le pilote KMS pour les GPU AMD Radeon :

\begin{verbatim}
/* from drivers/gpu/drm/radeon/radeon_kms.c */
int radeon_info_ioctl(struct drm_device *dev,
                      void *data,
                      struct drm_file *filp) {
	struct radeon_device *rdev =
            dev->dev_private;
	struct drm_radeon_info *info;
	struct radeon_mode_info *minfo =
            &rdev->mode_info;
	uint32_t *value_ptr;
	uint32_t value;
	struct drm_crtc *crtc;
	int i, found;

	info = data;
	value_ptr = (uint32_t *)
            ((unsigned long)info->value);
	value = *value_ptr;
        [...]
}
\end{verbatim}

On peut voit que l'argument \texttt{data} est converti en un \texttt{struct
drm\_radeon\_info *}. Un pointeur \texttt{value\_ptr} est extrait de son champ
\texttt{value}, et finalement ce pointeur est déréferencé.

Cependant, l'argument \texttt{data} est un pointeur vers une structure (allouée
en espace noyau) du type suivant, dont les champs proviennent d'un appel
utilisateur de \texttt{ioctl()}.

\begin{verbatim}
/* from include/drm/radeon_drm.h */
struct drm_radeon_info {
  uint32_t    request;
  uint32_t    pad;
  uint64_t    value;
};
\end{verbatim}

Pour mettre ce problème en évidence, nous avons annoté la fonction
\texttt{radeon\_info\_ioctl} de telle manière que son second paramètre soit un
pointeur noyau vers une structure contenant un champ \qUser, \texttt{value}.
Ceci est possible puisqu'avant la traduction, on efface les types présents dans
le programme C. Ainsi, un pointeur ne peut pas être distingué d'un entier
transtypé en un pointeur. Avec cette configuration, nous obtenons une erreur de
type à la ligne 16.

L'intégralité de ce code peut être trouvée en annexe~\ref{cha:code-noyau}.

\begin{figure}
  \insertcode{radeon-git-mini.diff}

  \caption{Patch résolvant le problème de pointeur utilisateur.}
  \label{fig:linux-patch}
\end{figure}

% TODO décrire syntaxe unified diff

La bonne manière de faire a été publiée avec le numéro de \emph{commit}
\texttt{d8ab3557} (figure~\ref{fig:linux-patch}) (\texttt{DRM\_COPY\_FROM\_USER}
étant une simple macro pout \texttt{copy\_from\_user}). Dans ce cas, on
n'obtient pas d'erreur de typage.

% TODO et l'assembleur inline?

\section{Détails}

Pour utiliser notre système de types, plusieurs étapes sont nécessaires en plus
de traduire le noyau linux en \langname.

Afin de réaliser l'analyse, il faut annoter les sources pour créer un
environnement initial (via la variable \texttt{exttbl} décrite en
section~\ref{sec:ptrtype-archi}).

% TODO dire ce qu'on annote : les sources de pointeurs souillés

Ensuite, il faut réécrire les fonctions de manipulation de pointeurs fournies
par le noyau : \texttt{get\_user}, \texttt{put\_user},
\texttt{copy\_from\_user}, \texttt{copy\_to\_user}, etc. Leur implantation
revient à réaliser un test puis à faire la copie. À leur place on utilise les
primitives $\phxx{\uGet}$ et $\phxx{\uPut}$.

% TODO ... ou on ne met pas ces primitives là en tyspeak et on ajoute juste un
% bon type dans l'environnement initial

Enfin, on peut lancer l'inférence de type.


\partccl{part:xp}
Après avoir décrit notre solution théorique dans la partie~\ref{part:lang}, nous
avons présenté ici notre démarche expérimentale. Dans le
chapitre~\ref{cha:implem}, nous avons détaillé l'implantation de notre
prototype. Pour ce faire, nous avons ajouté des étiquettes de type au langage
\newspeak et implanté un algorithme d'inférence de types. Ce prototype est
distribué sur~\link{penjili} sous le nom de \ptrtype.

Ensuite, le but du chapitre~\ref{cha:etudedecas} est d'appliquer notre analyse
(à l'aide de ce prototype) sur le noyau Linux. Après avoir décrit le
fonctionnement des appels système sur ce noyau, on présente deux bugs qui ont
touché respectivement un pilote de carte graphique et l'implantation d'un appel
système. Ils sont la manifestation d'un problème de pointeur utilisateur mal
déréférencé dans le noyau, ainsi que décrit dans le chapitre~\ref{cha:os}. En
lançant notre analyse sur le code présentant un problème, l'erreur est détectée.
Au contraire, en la lançant sur le code après application du correctif, aucune
erreur n'est trouvée.

En s'appuyant sur le langage \newspeak, on gagne beaucoup par rapport à d'autres
représentations intermédiaires. Le fait d'avoir un langage avec peu de
constructions permet de ne pas avoir à exprimer plusieurs fois la même règle
(par exemple, une fois sur la boucle \emph{for} et une autre sur la boucle
\emph{while}).

\label{page:ccl-npk-spk}

Un des inconvénients est que le modèle mémoire utilisé par \newspeak est assez
différent de celui de \langname (ainsi que décrit dans le
chapitre~\ref{cha:lang}). \newspeak est en effet prévu pour implanter des
analyses précises de valeur reposant sur l'interprétation abstraite, et
nécessite donc un modèle mémoire de plus bas niveau (où on peut créer des
valeurs à partir d'une suite d'octets, par exemple).

Le prototype d'implantation peut évoluer dans deux directions: d'une part, en
continuant à s'appuyer sur \newspeak, on peut réaliser des pré-analyses de
typage qui permettent de guider une analyse de valeurs plus précise, par exemple
en choisissant un domaine abstrait différent en fonction des types de données
rencontrés. D'autre part, il est possible de faire une implantation plus fidèle
à \langname, qui permette d'ajouter de nouvelles fonctionnalités plus éloignées
de C. Par exemple, un système de régions comme~\cite{jfp92} permettrait de
simplifier l'environnement d'exécution en enlevant l'opération de nettoyage
mémoire $\phx\cCleanup$. Le système de types peut également être enrichi, pour
ajouter par exemple du polymorphisme. Cela rapprocherait le langage source de
Rust. Le chapitre~\ref{cha:conclusion} présente quelques unes de ces extensions
possibles.

L'expérimentation, quant à elle, est pour le moment limitée, mais on peut
l'étendre à des domaines de plus en plus importants dans le noyau Linux. Tout
d'abord, le module graphique définit d'autres fonctions implantant des
\emph{ioctls}. Celles-ci reçoivent donc également des pointeurs utilisateur et
sont susceptibles d'être vulnérables à ce genre d'erreurs de programmation.
Ensuite, d'autres modules exposent une interface similaire, à commencer par les
autres pilotes de cartes graphiques. Ceux-ci sont également un terrain sur
lequel appliquer cette analyse.

De manière générale, toutes les interfaces du noyau manipulant des pointeurs
utilisateur gagnent à être analysées. Outre les implantations des \emph{ioctls}
dans chaque pilote et les appels système, les systèmes de fichiers manipulent
aussi de tels pointeurs via leurs opérations de lecture et d'écriture.

% vim: spelllang=fr


\part{Conclusion}

\chapter{Conclusion}

\label{cha:conclusion}
\section{Limitations}

Un seul type d'entier est supporté.

Notre système ne supporte pas certains cas de transtypage (\emph{casts}). Les
casts entre types entiers ne posent pas de problèmes. Ceux entre types pointeurs
ne posent pas de problème non plus puisqu'au moment de la traduction en
\langname on efface les étiquettes de types sur les pointeurs.

Cependant, les casts entre entiers et pointeurs posent problème : si on autorise
les casts, la seule manière sensible de les typer est la suivante.

\begin{mathpar}
\irule{PtrInt-Bad}
  { Γ ⊢ e : t~q~*}
  { Γ ⊢ (\tInt)~e : \tInt }

\irule{IntPtr-Bad}
  { Γ ⊢ e : \tInt }
  { Γ ⊢ (\textsc{Ptr})~e : t~q~*}
\end{mathpar}

Cela pose problème car il est alors possible de créer une fonction qui convertit
n'importe quel type pointeur en n'importe quel autre type pointeur :

\[
  ⊢ \eFun{p}{ \iReturn{(\textsc{Ptr})~(\tInt)~p} }
  : (t_a~q_a~*) → t_b~q_b~*
\]

% TODO et à l'exec?
% TODO on peut faire un dynamic check ou un type UInt ou un qualif sur tous les
% types

% TODO le modèle mémoire de Newspeak est bcp plus bas niveau que celui de
% Safespeak

\section{Travaux futurs}


\appendix
\renewcommand{\appendixpagename}{Annexes}
\renewcommand{\appendixtocname}{\appendixpagename}
\appendixpage*

\chapter{Module Radeon KMS}

\label{cha:code-noyau}
On inclut ici le code analysé dans le chapitre~\ref{cha:etudedecas}. On inclut à
la suite le contexte nécessaire pour comprendre ce code.

{ \small
\insertcode{radeon-buggy.c}
}


\chapter{Règles d'évaluation}

\begin{figure}

\begin{mathpar}
    \semrule{Ctx}

    \semrule{Ctx-Lv}

    \semrule{Ctx-Instr}
\end{mathpar}

    \figctx{}

\caption{Règles d'évaluation -- contextes}
\end{figure}

\begin{figure}

\begin{mathpar}
    \ruleheader{$Ξ → Ω$}

    \semrule{Exp-Err}

    \semrule{Eval-Err}
\end{mathpar}

\caption{Règles d'évaluation -- erreurs}
\end{figure}

\begin{figure}

\begin{mathpar}
    \ruleheader{$\mm{m}{lv}{m}{φ}$}

    \semrule{Phi-Var}

    \semrule{Phi-Deref}

    \semrule{Phi-Struct}

    \semrule{Phi-Array}

    \ruleheader{$\mm{m}{e}{m}{v}$}

    \semrule{Exp-Cst}

    \semrule{Exp-Fun}

    \semrule{Exp-Lv}

    \semrule{Exp-UnOp}

    \semrule{Exp-BinOp}

    \semrule{Exp-AddrOf}

    \semrule{Exp-Set}

    \semrule{Exp-Struct}

    \semrule{Exp-Array}

    \semrule{Exp-Call}
\end{mathpar}

\caption{Règles d'évaluation -- left-values et expressions }
\end{figure}

\begin{figure}

\begin{mathpar}
    \ruleheader{$\mm{m}{i}{m}{i'}$}

    \semrule{Seq}

    \semrule{Pass}

    \semrule{Exp}

    \semrule{Decl}

    \semrule{Decl-Return}

    \semrule{If-False}

    \semrule{If-True}

    \semrule{While}

    \semrule{Return}

    \ruleheader{$\ph{m}{p}{m'}$}

    \semrule{T-Exp}

    \semrule{T-Var}

    \ruleheader{$\phstar{m}{p}{m'}$}

    \semrule{T*-Nil}

    \semrule{T*-Cons}

    \ruleheader{$\phprog{P}{m}$}

    \semrule{Prog}
\end{mathpar}

\caption{Règles d'évaluation -- instructions et phrases}
% TODO programmes
\end{figure}

\begin{figure}

\begin{mathpar}
  \semrule{User-Get-OK}

  \semrule{User-Get-Err}

  \semrule{User-Put-OK}

  \semrule{User-Put-Err}
\end{mathpar}

\caption{Règles d'évaluation -- extensions noyau}
\end{figure}


\chapter{Règles de typage}

\begin{figure}[h]

\begin{mathpar}
    \ruleheader{$Γ ⊢ c : t$}

    \disprule{Cst-Int}

    \disprule{Cst-Float}

    \disprule{Cst-Null}

    \disprule{Cst-Unit}

    \ruleheader{$Γ ⊢ lv : t$}

    \disprule{Lv-Var}

    \disprule{Lv-Deref}

    \disprule{Lv-Index}

    \disprule{Lv-Field}
\end{mathpar}
\caption{Règles de typage -- constantes et variables}
\end{figure}

\begin{figure}
\begin{mathpar}
    \ruleheader{$Γ ⊢ \opun~e : t$}

    \disprule{Unop-Plus-Int}

    \disprule{Unop-Plus-Float}

    \disprule{Unop-Minus-Int}

    \disprule{Unop-Minus-Float}

    \disprule{Unop-Not}

    \ruleheader{$Γ ⊢ e_1~\opbin~e_2 : t$}

    \disprule{Op-Int}

    \disprule{Op-Float}

    \disprule{Op-Eq}

    \disprule{Ptr-Arith}

    \ruleheader{$\textsc{Eq}(t)$}

    \disprule{Eq-Num}

    \disprule{Eq-Ptr}

    \disprule{Eq-Array}

    \disprule{Eq-Struct}
\end{mathpar}
\caption{Règles de typage -- opérateurs}
\end{figure}

\begin{figure}
\begin{mathpar}
    \ruleheader{$Γ ⊢ e : t$}

    \disprule{Addr}

    \disprule{Struct}

    \disprule{Call}

    \disprule{Set}

    \disprule{Array}

    \disprule{Fun}

    \ruleheader{$Γ ⊢ i$}

    \disprule{Pass}

    \disprule{Seq}

    \disprule{Exp}

    \disprule{Decl}

    \disprule{If}

    \disprule{While}

    \disprule{Return}

    \ruleheader{$\typh{Γ}{p}{Γ'}$}

    \disprule{T-Exp}

    \disprule{T-Var}
\end{mathpar}

\caption{Règles de typage}
\end{figure}


\chapter{Preuves}

On présente ici les preuves de certains résultats établis dans le manuscrit: le
caractère bien fondé de la composition de deux lentilles, et les théorèmes de
sûreté du typage.

\section{Composition de lentilles}

\label{proof:compo-lens}

\begin{proof}

On cherche à prouver que si 
$ℒ_1 ∈ \setLens{A}{B}$
et
$ℒ_2 ∈ \setLens{B}{C}$, alors
$ℒ = ℒ_1 \ggg ℒ_2 ∈ \setLens{A}{C}$ ($\ggg$ est la composition de lentilles,
définie page~\pageref{def:lens-comp}).

Il suffit pour cela d'établir les trois propriétés caractéristiques qui
définissent les lentilles: \textsc{PutPut}, \textsc{GetPut} et \textsc{PutGet}.
Cela est essentiellement calculatoire: on utilise la définition de $\ggg$ et les
propriétés caractéristiques sur $ℒ_1$ et $ℒ_2$.

\paragraph{PutPut}%{{{

\begin{conteq}[onecolumn]
  \lensPut{ℒ}{a'}{
    \lensPut{ℒ}{a}{r}
  } \\
= \lensPut{ℒ}{a'}{
    \lensPut{ℒ_1}
      {
        \lensPut{ℒ_2}
          {a}
          { \lensGet{ℒ_1}{r} }
      }
      { r }
  }
& définition de \lensPutX{ℒ} \\
= \lensPut{ℒ_1}
    {
      \lensPut{ℒ_2}
        {a'}
        { \lensGet{ℒ_1}{
            \lensPut{ℒ_1}
              {
                \lensPut{ℒ_2}
                  {a}
                  { \lensGet{ℒ_1}{r} }
              }
              { r }
          }
        }
    }
    {
        \lensPut{ℒ_1}
          {
            \lensPut{ℒ_2}
              {a}
              { \lensGet{ℒ_1}{r} }
          }
          { r }
    }
& définition de \lensPutX{ℒ} \\
= \lensPut{ℒ_1}
    {
      \lensPut{ℒ_2}
        {a'}
        {
          \lensPut{ℒ_2}
            {a}
            { \lensGet{ℒ_1}{r} }
        }
    }
    {
        \lensPut{ℒ_1}
          {
            \lensPut{ℒ_2}
              {a}
              { \lensGet{ℒ_1}{r} }
          }
          { r }
    }
& \textsc{GetPut} sur $ℒ_1$ \\
= \lensPut{ℒ_1}
    {
      \lensPut{ℒ_2}
        {a'}
        { \lensGet{ℒ_1}{r} }
    }
    {
        \lensPut{ℒ_1}
          {
            \lensPut{ℒ_2}
              {a}
              { \lensGet{ℒ_1}{r} }
          }
          { r }
    }
& \textsc{PutPut} sur $ℒ_2$ \\
= \lensPut{ℒ_1}
    {
      \lensPut{ℒ_2}
        {a'}
        { \lensGet{ℒ_1}{r} }
    }
    { r
    }
& \textsc{PutPut} sur $ℒ_1$ \\
= \lensPut{ℒ}{a'}{r}
& définition de $\ggg$ \\
\end{conteq}
%}}}

\paragraph{GetPut}%{{{

{

\def\xa{
  \lensPut{ℒ}{
    \lensGet{ℒ}{r}
    }{
      r
    }
  }


\def\xb{
  \lensPut{ℒ}
    {
      \lensGet{ℒ_2}{
        \lensGet{ℒ_1}{r}
      }
    }
    { r }
}

\def\xab{ définition de \lensGetX{ℒ} }

\def\xc{
  \lensPut{ℒ_1}
          {
            \lensPut{ℒ_2}
                    {
                      \lensGet{ℒ_2}
                        {
                          \lensGet{ℒ_1}{r}
                        }
                    }
                    { \lensGet{ℒ_1}
                              {r}
                    }
          }
          {r}
}

\def\xbc{ définition de \lensPutX{ℒ} }

\def\xd{
  \lensPut{ℒ_1}
          {
            \lensGet{ℒ_1}{r}
          }
          {r}
}

\def\xcd{ \textsc{GetPut} sur $ℒ_2$ }

\def\xe{r}

\def\xde{\textsc{GetPut} sur $ℒ_1$}

\begin{conteq}
    \xa{} \\
  = \xb{} & \xab{} \\
  = \xc{} & \xbc{} \\
  = \xd{} & \xcd{} \\
  = \xe{} & \xde{}
\end{conteq}

}

%}}}

\paragraph{PutGet}%{{{

\begin{conteq}
  \lensGet{ℒ}
    {
      \lensPut{ℒ}
        {a}
        {r}
    } \\
= \lensGet{ℒ_2}
    {
      \lensGet{ℒ_1}
        {
          \lensPut{ℒ}
            {a}
            {r}
        }
    }
& définition de \lensGetX{ℒ} \\
= \lensGet{ℒ_2}
    {
      \lensGet{ℒ_1}
        {
          \lensPut{ℒ_1}
            {
              \lensPut{ℒ_2}
                {a}
                { \lensGet{ℒ_1}{r} }
            }
            { r }
        }
    }
& définition de \lensPutX{ℒ} \\
= \lensGet{ℒ_2}
    {
      \lensPut{ℒ_2}
        {a}
        { \lensGet{ℒ_1}{r} }
    }
& \textsc{PutGet} sur $ℒ_1$ \\
= a
& \textsc{PutGet} sur $ℒ_2$ \\
\end{conteq}
%}}}

\end{proof}

\section{Progrès}
\label{proof:progres}

On rappelle l'énoncé du théorème~\ref{thm:progres}.

\begin{theorem}[Progrès]

Supposons que $Γ ⊢ i$. Soit $m$ un état mémoire tel que $\mcomp{Γ}{m}$.

Alors l'un des cas suivants est vrai:
\begin{itemize}
\item $i = \iPass$
\item $∃v, i = \iReturn{v}$
\item $∃ (i', m'), \mm{m}{i}{m'}{i'}$
\item $∃ Ω ∈ \{\serr{div},\serr{array},\serr{ptr}\}, \msi{m}{i} → Ω$
\end{itemize}

\jolibreak

  Supposons que $Γ ⊢ e : t$. Soit $m$ un état mémoire tel que $\mcomp{Γ}{m}$.
  Alors l'un des cas suivant est vrai:

\begin{itemize}
  \item $∃ v ≠ Ω, e = v$
  \item $∃ (e', m'), \mm{m}{e}{m'}{e'}$
  \item $∃ Ω ∈ \{\serr{div},\serr{array},\serr{ptr}\}, \msi{m}{e} → Ω$
\end{itemize}

\jolibreak

Supposons que $Γ ⊢ lv : t$. Soit $m$ un état mémoire tel que $\mcomp{Γ}{m}$.

Alors l'un des cas suivants est vrai:
\begin{itemize}
\item $∃φ, lv = φ$
\item $∃ (lv', m'), \mm{m}{lv}{m'}{lv'}$
\item $∃ Ω ∈ \{\serr{div},\serr{array},\serr{ptr}\}, \msi{m}{lv} → Ω$
\end{itemize}

\end{theorem}

C'est-à-dire, soit:

\begin{itemize}
  \item l'entité (instruction, expression ou valeur gauche) est complètement
évaluée.
  \item un pas d'évaluation est possible.
  \item une erreur de division, tableau ou pointeur se produit.
\end{itemize}

\begin{proof}

On procède par induction sur la dérivation du jugement de typage. Puisque les
jugements $Γ ⊢ i$, $Γ ⊢ e : t$ et $Γ ⊢ lv : t$ sont interdépendants, on traite
tous les cas par récursion mutuelle.

Le squelette de cette preuve est une analyse de cas selon la dernière règle
utilisée. La plupart des cas ont la même forme: on utilise l'hypothèse de
récurrence sur les sous-éléments syntaxiques (en appliquant éventuellement le
lemme~\ref{lemma:inversion} d'inversion pour établir qu'ils sont bien typés).
Dans le cas \enquote{valeur}, on appelle une règle qui permet de transformer une
opération syntaxique en opération sémantique (par exemple, on transforme le $+$
unaire en un $\widehat{+}$ sémantique). Dans le cas \enquote{évaluation}, on
applique la règle \textsc{Ctx} avec un contexte particulier qui permet de passer
d'un jugement $\mm{m}{a}{m'}{a'}$ à un jugement $\mm{m}{b}{m'}{b'}$ (où $a$
apparaît dans $b$). Enfin, dans le cas \enquote{erreur}, on utilise
\textsc{Eval-Err} avec ce même contexte $C$.

Ceci est valable pour la majorité des cas. Il faut faire attention en
particulier aux opérations sémantiques qui peuvent produire des erreurs (comme
la division, ou l'opérateur $\phxx\cLookup$).

%\begin{itemize}
%\item $i = \iPass$.
%\item $i = \iReturn{v}$.
%\item $\mm{m}{i}{m'}{i'}$.
%\item $\msi{m}{i} → Ω$.
%\end{itemize}

%\begin{itemize}
%\item $e = v$.
%\item $\mm{m}{e}{m'}{e'}$.
%\item $\msi{m}{e} → Ω$.
%\end{itemize}

%\begin{itemize}
%\item $lv = φ$.
%\item $\mm{m}{lv}{m'}{lv'}$.
%\item $\msi{m}{lv} → Ω$.
%\end{itemize}

\subsection*{Instructions}

\paragraph{\textsc{Pass}  :} %{{{
Ce cas est immédiat.
% }}}
\paragraph{\textsc{Return}:} %{{{

Partant de $i = \iReturn{e}$, on applique le lemme d'inversion. Il nous donne
l'existence de $t$ tel que $Γ ⊢ e : t$. On applique alors l'hypothèse de
récurrence à $e$.

\begin{itemize}
\item $e = v$. Alors $i = \iReturn{v}$, ce qui nous permet de conclure.

\item $\mm{m}{e}{m'}{e'}$.
Alors en appliquant \textsc{Ctx} avec
$C = \iReturn{\ctxEmpty}$, on conclut que
\linebreak
$\mm{m}{\iReturn{e}}{m'}{\iReturn{e'}}$.

\item $\msi{m}{e} → Ω$. On applique \textsc{Eval-Err} avec ce même $C$.

\end{itemize}
% }}}
\paragraph{\textsc{Seq}   :} %{{{
Avec $i = i_1;i_2$, on applique l'hypothèse de récurrence à $i_1$.

\begin{itemize}
\item $i_1 = \iPass$. On peut donc appliquer la règle \textsc{Seq} et donc
$\mms{i}{i_2}$.

\item $i_1 = \iReturn{v}$. Alors on peut appliquer la règle \textsc{Return}:
    $\mms{i}{\iReturn{v}}$.

\item $\mm{m}{i_1}{m'}{i_1'}$. Soit $C = \ctxEmpty;i_2$. Par \textsc{Ctx} il
    vient $\mm{m}{i}{m'}{i_1';i_2}$.
    \item $\msi{m}{i_1} → Ω$. Avec ce même $C$ dans \textsc{Eval-Err} on trouve
    $\msi{m}{i} → Ω$.

\end{itemize}
% }}}
\paragraph{\textsc{Exp}   :} %{{{

Ici $i = e$. On peut appliquer l'hypothèse de récurrence à $e$ qui est
\enquote{plus petit} que $i$ ($i \gramisa e$ introduit un constructeur
implicite).

\begin{itemize}
\item $e = v$. Alors on peut appliquer \textsc{Exp}: $\mms{e}{\iPass}$.
\item $\mm{m}{e}{m'}{e'}$. Alors $\mm{m}{i}{m'}{e'}$ (cela revient à appliquer
    \textsc{Ctx} au constructeur implicite mentionné ci-dessus).
\item $\msi{m}{e} → Ω$. C'est-à-dire $\msi{m}{i} → Ω$.
\end{itemize}

% }}}
\paragraph{\textsc{Decl}  :} %{{{

Ici $i = \iDecl{x}{e}{i'}$. On commence par appliquer l'hypothèse de récurrence
à $e$.

\begin{itemize}
\item $e = v$. On applique alors l'hypothèse de récurrence à $i'$ sous
    $Γ' = \extendLocal{Γ}{x}{t}$ et avec $m' = \cExtend{m}{x}{v}$.

    \begin{itemize}
    \item $i' = \iPass$. Dans ce cas la règle \textsc{Decl-Pass} s'applique.

    \item $i' = \iReturn{v}$. Idem avec \textsc{Decl-Return}.

    \item $\mm{m'}{i'}{m''}{i''}$. On peut alors appliquer la règle
          \textsc{Decl-Ctx}.

    \item $\msi{m'}{i'} → Ω$. On applique \textsc{Decl-Err}.

    \end{itemize}

\item $\mm{m}{e}{m'}{e'}$. On pose $C = \iDecl{x}{\ctxEmpty}{i'}$ et on conclut avec
    la règle \textsc{Ctx}.
\item $\msi{m}{e} → Ω$. Idem avec \textsc{Eval-Err}.
\end{itemize}



% }}}
\paragraph{\textsc{If}    :} %{{{
Ici $i = \iIf{e}{i_1}{i_2}$. On applique l'hypothèse de récurrence à $e$.

\begin{itemize}
\item $e = v$.

    Si $v ≠ 0$, on applique \textsc{If-True}. Dans le cas contraire, on
    applique \textsc{If-False}.

\item $\mm{m}{e}{m'}{e'}$. On pose $C = \iIf{\ctxEmpty}{i_1}{i_2}$ et on conclut
    avec \textsc{Ctx}.

\item $\msi{m}{e} → Ω$. Avec ce même $C$ et \textsc{Eval-Err}.
\end{itemize}

% }}}
\paragraph{\textsc{While} :} %{{{

Ce cas est direct: on applique la règle d'évaluation \textsc{While}.

% }}}

\subsection*{Expressions}

  \paragraph{\textsc{Cst-Int}:} % {{{
$e$ est alors de la forme $n$, qui est une valeur.
%}}}
  \paragraph{\textsc{Cst-Float}:} % {{{
$e$ est alors de la forme $d$, qui est une valeur.
%}}}
  \paragraph{\textsc{Cst-Null}:} % {{{
$e$ est alors égale à $\eNull$, qui est une valeur.
%}}}
  \paragraph{\textsc{Cst-Unit}:}%{{{
$e$ est alors égale à \eUnit, qui est une valeur.
%}}}
\paragraph{\textsc{Fun}:} % {{{

Ce cas est direct: la règle \textsc{Exp-Fun} s'applique.

%}}}
  \paragraph{\textsc{Op-Int}:} % {{{

  Cela implique que $e = e_1~\opbin~e_2$. Par le lemme~\ref{lemma:inversion}, on
  en déduit que $Γ ⊢ e_1 : \tInt$ et $Γ ⊢ e_2 : \tInt$.

  Appliquons l'hypothèse de récurrence sur $e_1$. Trois cas peuvent se produire.

\begin{itemize}

  \item $e_1 = v_1$. On a alors $\mm{m}{e_1}{m'}{v_1}$ avec $m' = m$.

    On applique l'hypothèse de récurrence à $e_2$.

      \begin{itemize}

        \item $e_2 = v_2$: alors $\mm{m'}{e_2}{m''}{v_2}$ avec $m'' = m$. On
          peut alors appliquer \textsc{Exp-BinOp}, sauf dans le cas d'une
          division par zéro ($ \opbin ∈ \{ / ; \% ; /. \} $ et
          $ v_2 = 0 $) où alors $v_1~\widehat{\opbin}~v_2 = \serr{div}$. Dans ce cas, on a
          alors par \textsc{Exp-Err} $\msi{m}{e} → \serr{div}$.
          Notons que comme les opérandes sont bien typées, $\serr{typ}$ ne peut pas être levée.

        \item $∃ (e'_2, m''), \mm{m'}{e_2}{m''}{e'_2}$.

          En appliquant \textsc{Ctx} avec $C = \ctxOp{v_1}{\ctxEmpty}$, on
          en déduit $\mm{m'}{v_1~\opbin~e_2}{m''}{v_1~\opbin~e'_2}$ soit
          $\mm{m}{e}{m''}{v_1~\opbin~e'_2}$.

        \item $\msi{m'}{e_2} → Ω$.
          De \textsc{Eval-Err} avec $C = \ctxOp{v_1}{\ctxEmpty}$
          vient alors $\msi{m}{e} → Ω$.

      \end{itemize}

  \item $∃(e_1', m'), \mm{m}{e_1}{m'}{e'_1}$.
    En appliquant \textsc{Ctx} avec $C = \ctxOp{\ctxEmpty}{e_2}$, on obtient
    \linebreak
    $\mm{m}{e_1~\opbin~e_2}{m'}{e'_1~\opbin~e_2}$, ou
    $\mm{m}{e}{m'}{e'_1~\opbin~e_2}$.

  \item $\msi{m}{e_1} → Ω$.
    D'après \textsc{Eval-Err} avec $C = \ctxOp{\ctxEmpty}{e_2}$, on a
    $\msi{m}{e} → Ω$.

\end{itemize}

% }}}
\paragraph{\textsc{Op-Float}:} % {{{
Ce cas est similaire à \textsc{Op-Int}.
%}}}
\paragraph{\textsc{Op-Eq}:} %{{{
Ce cas est similaire à \textsc{Op-Int}.
%}}}
\paragraph{\textsc{Unop-Plus-Int}:} % {{{

Alors $e = +~e_1$. En appliquant l'hypothèse d'induction sur $e_1$:

\begin{itemize}
\item
  soit $e_1 = v_1$. Alors en appliquant \textsc{Exp-UnOp},
  $\mm{m}{+~v_1}{m}{\widehat{+}~v_1}$, c'est-à-dire en posant $v =
  \widehat{+}~v_1$, $\mm{m}{e}{m}{v}$.
\item
  soit $∃ e'_1, m', \mm{m}{e_1}{m'}{e'_1}$. Alors en appliquant \textsc{Ctx}
avec $C = +~\ctxEmpty$, on obtient \linebreak $\mm{m}{e}{m'}{e'_1}$.
\item
  soit $\msi{m}{e_1} → Ω$.
  De \textsc{Eval-Err} avec $C = +~\ctxEmpty$ il vient$\msi{m}{e} → Ω$.
\end{itemize}

% }}}
\paragraph{\textsc{Unop-Plus-Float}:} % {{{
Ce cas est similaire à \textsc{Unop-Plus-Int}.
% }}}
\paragraph{\textsc{Unop-Minus-Int}:} % {{{
Ce cas est similaire à \textsc{Unop-Plus-Int}.
% }}}
\paragraph{\textsc{Unop-Minus-Float}:} % {{{
Ce cas est similaire à \textsc{Unop-Plus-Int}.
% }}}
\paragraph{\textsc{Unop-Not}:}%{{{
Ce cas est similaire à \textsc{Unop-Plus-Int}.
%}}}
\paragraph{\textsc{Addr}:} % {{{

On applique l'hypothèse de récurrence à $lv$.

Les cas d'évaluation et d'erreur sont traités en appliquant respectivement
\textsc{Ctx} et \textsc{Eval-Err} avec $C = \&\ctxEmpty$. Dans le cas où $lv =
φ$, on peut appliquer \textsc{Exp-Addr}.

%}}}
\paragraph{\textsc{Set}:} % {{{

On applique l'hypothèse de récurrence à $lv$.

\begin{itemize}

\item $lv = φ$. On applique l'hypothèse de récurrence à $e$.

    \begin{itemize}
    \item $e = v$. Alors on peut appliquer \textsc{Exp-Set}.

    \item $\mm{m}{e}{m'}{e'}$. On conclut avec $C = φ ← \ctxEmpty$.
    \item $\msi{m}{e} → Ω$. Idem.
    \end{itemize}

\item $\mm{m}{lv}{m'}{lv'}$. On conclut avec $C = \ctxEmpty ← e$.
\item $\msi{m}{lv} → Ω$. Idem.

\end{itemize}

%}}}
\paragraph{\textsc{Array}:} % {{{

On va appliquer l'hypothèse de récurrence à $e_1$, puis si $e_1 = v_1$, on
l'applique à $e_2$, etc. Alors on se retrouve dans un des cas suivants:

\begin{itemize}
\item $∃ p ∈ [1;n],e'_p,m: e_1 = v_1, …, e_{p-1} = v_{p-1}, \mm{m}{e_p}{m'}{e'_p}$.
  Alors on peut appliquer \textsc{Ctx} avec
  $C = [v_1, …, v_{p-1}, \ctxEmpty, e_{p+1}, …, e_n]$.
\item $∃ p ∈ [1;n],Ω :     e_1 = v_1, …, e_{p-1} = v_{p-1}, \msi{m}{e_p} → Ω$.
  Dans ce cas \textsc{Eval-Err} est applicable avec ce même $C$.
\item $e_1 = v_1, …, e_n = v_n$.
  Alors on peut appliquer \textsc{Exp-Array} en construisant un tableau.
\end{itemize}

%}}}
\paragraph{\textsc{Struct}:} % {{{

Le schéma de preuve est similaire au cas \textsc{Array}.
En cas de pas d'évaluation ou d'erreur, on utilise le contexte
$C = \eStruct{l_1: v_1, …, l_{p-1}: v_{p-1}, \ctxEmpty, l_{p+1}: e_{p+1}, …, l_n:
e_n}$; et dans le cas où toutes les expressions sont évaluées, on applique
\textsc{Exp-Struct}.

%}}}
\paragraph{\textsc{Call}:} % {{{

On commence par appliquer l'hypothèse de récurrence à $e$. Dans le cas d'un pas
d'évaluation ou d'erreur, on applique respectivement \textsc{Ctx} ou
\textsc{Eval-Err} avec $C = \ctxEmpty (e_1, …, e_n)$.
Reste le cas où $e$ est une valeur: d'après le lemme~\ref{lemma:canon}, $e$ est
de la forme $f = \eFun{\vec{a}}{i}$.

Ensuite, appliquons le même schéma que pour \textsc{Array}.
En cas de pas d'évaluation ou d'erreur, on utilise
\textsc{Ctx} ou \textsc{Eval-Err} avec
$C = f (v_1, …, v_{p-1}, \ctxEmpty, e_{p+1}, …, e_n)$.
Le seul cas restant est celui où l'expression considérée a pour forme
$f (v_1, …, v_n)$
avec
$f = \eFun{\vec{a}}{i}$.

Soient $Γ' = (Γ_G, [a_1 : t_1, …, a_n : t_n, \vRet : t])$
et
$m_1 = \cPush{m_0}{(a_1 ↦ v_1, … a_n ↦ v_n)}$
où $Γ = (Γ_G, Γ_L)$.

On applique alors l'hypothèse de récurrence à $Γ'$, $m_1$ et $i$ (le lemme
d'inversion garantit que $Γ' ⊢ i$).

\begin{itemize}
\item $i = \iPass$. Ce cas est impossible puisqu'on prend l'hypothèse que les
fonctions se terminent par une instruction $\iReturn{\cdot}$
(page~\pageref{page:return-fonction}).

\item $\mm{m_1}{i}{m_2}{i'}$.

    Alors on peut appliquer \textsc{Exp-Call-Ctx}.

\item $\msi{m}{i} → Ω$. On peut alors appliquer \textsc{Exp-Call-Err}.

\end{itemize}

%}}}

\subsection*{Valeurs gauches}

\paragraph{\textsc{Lv-Var}:}%{{{

Le but est d'appliquer \textsc{Phi-Var}. La seule condition pour que cela soit
possible est que $\cLookup{x}{m}$ renvoie une adresse et non $\serr{var}$.

Puisque $Γ ⊢ x : t$, on peut appliquer le lemme~\ref{lemma:var-typ}: $x$ est
soit une variable locale, soit une globale. Dans ces deux cas, $\cLookup{x}{m}$
renvoie une adresse correcte.

%}}}
\paragraph{\textsc{Lv-Deref}:}%{{{

  Appliquons l'hypothèse de récurrence à $e$ vue en tant qu'expression.

\begin{itemize}
\item
  $e = v$. Puisque $Γ ⊢ v : t*$, on déduit du
  lemme~\ref{lemma:canon} que
  $v = \eNull$
  ou
  $v = \widehat{\&}~φ$.

  Dans le premier cas, puisque $\msi{m}{*\eNull} → \serr{ptr}$, on a
  $\msi{m}{e} → \serr{ptr}$.

  Dans le second cas, \textsc{Phi-Deref} s'applique.

\item
  $\mm{m}{e}{m'}{e'}$.
  De \textsc{Ctx} avec $C = *\ctxEmpty$, on obtient
  $\mm{m}{e}{m'}{*e'}$.

\item
  $\msi{m}{e} → Ω$.
  En appliquant \textsc{Eval-Err} avec $C = *\ctxEmpty$, on obtient
  $\msi{m}{e} → Ω$.

\end{itemize}

% }}}
\paragraph{\textsc{Lv-Index}:} % %{{{

De même, on applique l'hypothèse de récurrence à $lv$.

\begin{itemize}
\item $lv = v$.

Comme $Γ ⊢ v : t[]$, on déduit du lemme~\ref{lemma:canon} que
$v = \eArray{v_1; …; v_p}$.
Appliquons l'hypothèse de récurrence à $e$.

\begin{itemize}
\item $e = v'$. Puisque $Γ ⊢ e : \tInt$, on réapplique le
lemme~\ref{lemma:canon} et $v' = n$.
D'après \textsc{Phi-Array}, $ \mms{lv[e]}{\eArray{v_1; …; v_p}\widehat{[n]}} $.
Deux cas sont à distinguer:
si $n ∈ [0;p-1]$, la partie droite vaut $v_{n+1}$ et donc
$\mm{m}{lv[e]}{m}{v_{n+1}}$.
Sinon elle vaut $\serr{array}$ et $\msi{m}{lv[e]} → \serr{array}$ par \textsc{Exp-Err}.

\item $\mm{m}{e}{m'}{e'}$.
En appliquant \textsc{Ctx} avec $C = v[\ctxEmpty]$, on en déduit
\item $\mm{m}{lv[e]}{m'}{lv[e']}$.

\item $\msi{m}{e} → Ω$.
Avec \textsc{Eval-Err} sous ce même contexte,
$\msi{m}{lv[e]} → Ω$
\end{itemize}

\item $\mm{m}{lv}{m'}{e'}$.
On applique alors \textsc{Ctx} avec $C = \ctxEmpty[e]$, et
$\mm{m}{lv[e]}{m'}{e'[e]}$.

\item $\msi{m}{lv} → Ω$.
Toujours avec $C = \ctxEmpty[e]$, de \textsc{Eval-Err} il vient
$\msi{m}{lv[e]} → Ω$.

\end{itemize}
%}}}
\paragraph{\textsc{Lv-Field}:}%{{{

On applique l'hypothèse de récurrence à $lv$.

\begin{itemize}

\item $lv = φ$
Alors \textsc{Phi-Struct} s'applique. Puisque $(l, t) ∈ S$, l'accès au champ $l$
ne provoque pas d'erreur $\serr{field}$. Donc $\mm{m}{e}{m}{φ[l]}$.

\item $\mm{m}{lv}{m'}{lv'}$
En appliquant \textsc{Ctx} avec $C = \ctxEmpty.l_S$, il vient
$\mm{m}{lv}{m'}{lv'}$.

\item $\msi{m}{lv} → Ω$
En appliquant \textsc{Eval-Err} avec $C = \ctxEmpty.l_S$, on a
$\msi{m}{lv} → Ω$.

\end{itemize}

%}}}
\paragraph{\textsc{Ptr-Arith}:} % %{{{

Le schéma est similaire au cas \textsc{Op-Int}. Le seul cas intéressant arrive
lorsque $e_1$ et $e_2$ sont des valeurs. D'après le lemme~\ref{lemma:canon}:

\begin{itemize}
\item $e_1 = \eNull$ ou $e_1 = φ$
\item $e_2 = n$
\end{itemize}

D'après \textsc{Exp-Binop}, $\mms{e}{e_1~\widehat{\opbin}~n}$.

On se réfère ensuite à la définition de $\widehat{\opbin}$
(page~\pageref{page:def-arith-ptr-error}): si $e_1$ est de la forme $φ[m]$,
alors $e_1~\widehat{\opbin}~n = φ[m+n]$. Donc $\mms{e}{φ[m+n]}$.

Dans les autres cas ($e_1 = \eNull$ ou $e_1 = φ$ avec $φ$ pas de la forme
$φ'[m]$), on a $e_1~\widehat{\opbin}~n = \serr{ptr}$. Donc d'après
\textsc{Exp-Err}, $\msi{m}{e} → \serr{ptr}$.

%}}}
\end{proof}

\section{Préservation}
\label{proof:preservation}

On rappelle l'énoncé du théorème~\ref{thm:preservation}.

\begin{theorem}[Préservation]

Soient $Γ$ un environnement de typage, et $m$ un état mémoire tels que
$\mcomp{Γ}{m}$.

Alors:

\begin{itemize}
\item
    Si $Γ ⊢ lv : t$ et $\mm{m}{lv}{m'}{φ}$,
    alors $\mcomp{Γ}{\cCleanup{m'}}$ et $\semtypphi{m'}{φ}{τ}$ où $\tComp{τ}{t}$.

\item
    Si $Γ ⊢ lv : t$ et $\mm{m}{lv}{m'}{lv'}$,
    alors $\mcomp{Γ}{\cCleanup{m'}}$ et $Γ ⊢ lv' : t$.

\item
    Si $Γ ⊢ e : t$ et $\mm{m}{e}{m'}{v}$,
    alors $\mcomp{Γ}{\cCleanup{m'}}$ et $m' ⊧ v : τ$ où $\tComp{τ}{t}$.

\item
    Si $Γ ⊢ e : t$ et $\mm{m}{e}{m'}{e'}$,
    alors $\mcomp{Γ}{\cCleanup{m'}}$ et $Γ ⊢ e' : t$.

\item
    Si $Γ ⊢ i$ et $\mm{m}{i}{m'}{i'}$,
    alors $\mcomp{Γ}{\cCleanup{m'}}$ et $Γ ⊢ i'$.
\end{itemize}

  Autrement dit, si une construction est typable, alors un pas d'évaluation ne
  modifie pas son type et préserve le typage de la mémoire.

\end{theorem}

\begin{proof}

On procède par induction sur la dérivation de $\mm{m}{\cdot}{m'}{\cdot}$.
Plusieurs remarques sont à faire: d'abord, en ce qui concerne le typage de la
mémoire, il suffit de montrer que $\mcomp{Γ}{m'}$ car cela implique que
$\mcomp{Γ}{\cCleanup{m'}}$. Ensuite, la règle \textsc{Ctx} est traitée à part,
car elle peut être appliquée en contexte d'expression, d'instruction ou de
valeur gauche. Enfin la règle \textsc{Trans} ne pose pas de problème, il suffit
d'appliquer l'hypothèse de récurrence à ses prémisses.

\subsection*{Cas $Γ ⊢ lv : t$ et $\mm{m}{lv}{m'}{φ}$}

\paragraph{\textsc{Phi-Deref}:} % {{{

On sait que $Γ ⊢ *~v : t$ où $v = \widehat{\&}~φ$.
Par inversion, $Γ ⊢ v : t~*$.
Alors d'après le lemme~\ref{lemma:repr}, il existe
$τ'$ tel que $m ⊧ v : τ'$ et $\tComp{τ'}{t~*}$.
Par inversion de la relation de typage sémantique,
$τ' = τ~*$ où $\tComp{τ}{t}$.
Alors par inversion de \textsc{S-Ptr}, on obtient que
$\semtypphi{m}{φ}{τ}$.

% }}}
\paragraph{\textsc{Phi-Var}, \textsc{Phi-Struct} et \textsc{Phi-Array}:} % {{{

Tout d'abord, la mémoire n'est pas modifiée donc il n'est pas nécessaire de
montrer la compatibilité de $m'$.
Ensuite, les prémisses de ces règles ont la forme $φ$ , donc le
lemme~\ref{lemma:mem-typ} s'applique avec la conclusion correcte.

% }}}

\subsection*{Cas $Γ ⊢ e : t$ et $\mm{m}{e}{m'}{v}$}

\paragraph{\textsc{Exp-Cst}:} % {{{

Toutes les constantes sont des valeurs, donc le lemme~\ref{lemma:repr} peut
s'appliquer : $τ = \cRepr{t}$ convient.

% }}}
\paragraph{\textsc{Exp-Fun}:} % {{{

Idem: le lemme de réprésentabilité nous donne un candidat $τ = \cRepr{t}$
qui convient.

% }}}
\paragraph{\textsc{Exp-Lv}:} % {{{

Puisque $Γ ⊢ φ : t$ et $\mcomp{Γ}{m}$, on a d'après le
lemme~\ref{lemma:mem-typ}: $m ⊧ v : τ$ où $v = m[φ]$ avec $\tComp{τ}{t}$.

% }}}
\paragraph{\textsc{Exp-UnOp}:} % {{{

Il vient des définitions des différents opérateurs $\widehat{\opun}$ que
$Γ ⊢ \widehat{\opun}~v : τ$ avec $\tComp{τ}{t}$.

% }}}
\paragraph{\textsc{Exp-BinOp}:} % {{{

Idem avec les définitions des opérateurs $\widehat{\opbin}$.

% }}}
\paragraph{\textsc{Exp-Addr}:} % {{{

On peut appliquer le lemme~\ref{lemma:repr}, qui nous donne un $τ$ qui convient.

% }}}
\paragraph{\textsc{Exp-Set}:} % {{{

Deux propriétés sont à prouver. D'une part, $Γ ⊢ v : t$, et d'autre part,
$\mcomp{Γ}{m'}$ où $m' = m[φ←v]$.
Tout d'abord, le lemme d'inversion appliqué à $Γ ⊢ φ ← v : t$ nous donne que
$Γ ⊢ φ : t$ et $Γ ⊢ v : t$.
Ensuite, comme $Γ ⊢ φ : t$ et $\mcomp{Γ}{m}$, on peut appliquer le
lemme~\ref{lemma:repr}: il existe $τ$ tel que $m ⊧ v : τ$ et $\tComp{τ}{t}$.
On peut donc appliquer le lemme~\ref{lemma:mem-typ}, qui nous permet de conclure
que $\mcomp{Γ}{m'}$.

% }}}
\paragraph{\textsc{Exp-Struct}:} % {{{

Le lemme~\ref{lemma:repr} s'applique à ce cas.

% }}}
\paragraph{\textsc{Exp-Array}:} % {{{

Idem, on conclut grâce au lemme de représentabilité.

% }}}
\paragraph{\textsc{Exp-Call-Return}:} % {{{

Par inversion, il vient que
$Γ ⊢ \eFun{a_1, …, a_n}{i} : (t_1, …, t_n) → t'$
et
$∀ i ∈ [1;n], Γ ⊢ v_i : t_i$.

Posons $Γ' = (Γ_G, [a_1:t_1, …, a_n:t_n, \vRet:t])$ où $Γ = (Γ_G, Γ_L)$. Alors
par inversions successives on obtient que $Γ' ⊢ \iReturn{v}$ et $Γ' ⊢ v : t$.

Si on définit $m'' = \cPush{m}{a_1↦v_1, …, a_n↦v_n}$, alors par \textsc{M-Push}
on obtient que $\mcomp{Γ'}{m''}$. Donc, par le lemme~\ref{lemma:repr}, il existe
$τ$ tel que $m'' ⊧ v : τ$ où $\tComp{τ}{t}$.

Il reste à montrer que $m' ⊧ v' : τ$.

On distingue selon la forme de $v$. On applique un raisonnement similaire à
celui de la preuve du lemme~\ref{lemma:mem-typ}: soit $v$ est une référence au
cadre nettoyé, et dans ce cas $v' = \eNull$ et $τ$ est un type pointeur, soit
$v' = v$. Dans tous les cas on conclut car $m' ⊧ v' : τ$.

% }}}

\subsection*{Cas $Γ ⊢ i$ et $\mm{m}{i}{m'}{i'}$}

\paragraph{\textsc{Seq}:}%{{{
D'après le lemme d'inversion, $Γ ⊢ i$.
%}}}
\paragraph{\textsc{Exp}:}%{{{
D'après \textsc{Pass}, $Γ ⊢ \iPass$.
%}}}
\paragraph{\textsc{Decl-Pass}:}%{{{

$Γ ⊢ \iPass$ est immédiat, et $\mcomp{Γ}{m'}$ est établi par \textsc{M-Decl}
suivie de \linebreak \textsc{M-DeclClean}. On a bien $x ∉ Γ$ car les
déclarations de variable ne peuvent pas masquer de variables visibles existantes
(page~\pageref{page:decl-masquage}).

%}}}
\paragraph{\textsc{Decl-Return}:}%{{{

La compatibilité mémoire se démontre de la même manière que pour
\textsc{Decl-Pass}. Il reste à montrer que $Γ ⊢ \iReturn{v''}$, ce qui fait de
manière analogue au cas \textsc{Exp-Call-Return}.

%}}}
\paragraph{\textsc{Decl-Ctx}:}%{{{

On part de $Γ ⊢ \iDecl{x}{v}{i}$.
Par inversion, il existe $t$ tel que
$Γ ⊢ v : t$ et $Γ' ⊢ i$ où $Γ' = \extendLocal{Γ}{x}{t}$.

Comme $\mcomp{Γ}{m}$, le lemme~\ref{lemma:repr} s'applique : il existe
$τ$ tel que $m ⊧ v : τ$ où $\tComp{τ}{t}$.
De plus $x ∉ Γ$ car il n'y a pas de masquage (page~\pageref{page:decl-masquage}).

En appliquant \textsc{M-Decl}, on obtient donc que $\mcomp{Γ'}{m'}$.

On applique alors l'hypothèse d'induction à $\mm{m'}{i}{m''}{i'}$. Il vient que
$Γ' ⊢ i'$ et $\mcomp{Γ'}{m''}$.

On a donc $Γ' ⊢ \iDecl{x}{v'}{i'}$ par \textsc{Decl} et
$\mcomp{Γ}{\cCleanup{m'''}}$ par \textsc{M-DeclClean}.

%}}}
\paragraph{\textsc{If-False}:}%{{{
D'après le lemme d'inversion, $Γ ⊢ i_f$.
%}}}
\paragraph{\textsc{If-True}:}%{{{
D'après le lemme d'inversion, $Γ ⊢ i_t$.
%}}}
\paragraph{\textsc{While}:}%{{{

D'après le lemme d'inversion, $Γ ⊢ e : t$ et $Γ ⊢ i$.
Par \textsc{Seq}, on a $Γ ⊢ i ; \iWhile{e}{i}$.
Enfin par \textsc{If} il vient $Γ ⊢ \iThen{e}{i ; \iWhile{e}{i}}$.

%}}}
\paragraph{\textsc{Return}:}%{{{

Par le lemme d'inversion, $Γ ⊢ \iReturn{v}$.
%}}}
\paragraph{\textsc{Exp-Call-Ctx}:}%{{{

On sait que $Γ ⊢ \eFun{a_1, …, a_n}{i} (v_1, …, v_n)$ et
$\mcomp{Γ}{m_0}$. D'après le lemme d'inversion, il existe
$t_1, …, t_n$ tels que $∀i ∈ [1;n], Γ ⊢ v_i : t_i$,
$Γ ⊢ \eFun{a_1, …, a_n}{i} : (t_1, …, t_n) → t$, donc qu'en posant
$Γ' = (Γ_G, [a_1:t_1, …, a_n:t_n, \vRet:t])$ où $Γ = (Γ_G, Γ_L)$ on a $Γ' ⊢ i$.

D'un autre côté, il existe par le lemme~\ref{lemma:repr} des types $τ_i$ tels
que $∀ i ∈ [1;n], m_0 ⊧ v_i : τ_i$ avec $\tComp{τ_i}{t_i}$. En appliquant
\textsc{M-Push}, on a donc $\mcomp{Γ'}{m_1}$.

On peut alors appliquer l'hypothèse d'induction à $\mm{m_1}{i}{m_2}{i'}$: la
conclusion est que $Γ' ⊢ i'$ et $\mcomp{Γ'}{m_2}$. Comme $Γ' ⊢ a_i : t_i$, on
a $∀ i ∈ [1;n], Γ' ⊢ v'_i : t_i$. Donc on a bien
$Γ ⊢ \eFun{a_1, …, a_n}{i'} (v'_1, …, v'_n) : t$.

D'autre part, en appliquant \textsc{M-Pop}, on obtient que
$\mcomp{Γ}{\cCleanup{m_3}}$.

%}}}

\added{Nouveau paragraphe}

\subsection*{À propos de la règle \textsc{Ctx}}

L'application de la règle \textsc{Ctx} nécessite une explication particulière.
En effet, ce cas repose sur un lemme d'inversion des constructions typées sous
un contexte, qui est admis ici.

Par exemple, traitons le cas où le contexte $C$ est tel que son \enquote{trou}
soit une valeur gauche $lv$ et $\ctxSub{C}{lv}$ est une instruction (les autres
cas sont similaires). La règle appliquée est alors de la forme:

\[
    \irule{Ctx}
         { \mm{m}{lv}{m'}{lv'} }
         { \mm{m}{\ctxSub{C}{lv}}{m'}{\ctxSub{C}{lv'}} }
\]

Si $Γ ⊢ \ctxSub{C}{lv}$, on admet qu'il existe $Γ'$ et $t$ tels que:

\begin{itemize}
    \item $Γ' ⊢ lv : t$;
    \item Quelque soit $lv'$,
        si
        $Γ' ⊢ lv' : t$,
        alors
        $Γ ⊢ \ctxSub{C}{lv'}$.
\end{itemize}

Par exemple, pour $C = \ctxEmpty{}[2] = 1$, $Γ' = Γ$ et $t = \tInt[~]$
conviennent.
Pour $C = \iDecl{x}{0}{\ctxEmpty = 3.0}$, on prendra
$Γ' = \extendLocal{Γ}{x}{\tInt} $ et $t = \tFloat$.
Le fait de passer \enquote{sous} une déclaration ajoute une variable locale à
$Γ$, et ainsi l'ensemble des variables de $Γ'$ contient celui de $Γ$.

Pour prouver la préservation dans ce cas, on commence par appliquer l'hypothèse
de récurrence à la prémisse de \textsc{Ctx}, c'est-à-dire $\mm{m}{lv}{m'}{lv'}$.
Il vient que $Γ' ⊢ lv' : t$ et $\mcomp{Γ'}{\cCleanup{m'}}$.

D'après le précédent lemme d'inversion on en déduit que $Γ ⊢ \ctxSub{C}{lv'}$.
De plus $Γ'$ contient plus de variables que $Γ$ donc $\mcomp{Γ}{\cCleanup{m'}}$.

\end{proof}

\section{Progrès pour les extensions noyau}
\label{proof:progres-qualif}

(Théorème~\ref{thm:progres-qual})

\begin{proof}

On procède de la même manière que pour le théorème~\ref{thm:progres} (prouvé en
annexe~\ref{proof:progres}). En fait, puisque le schéma de preuve porte sur les
règles de typage, il suffit de traiter les cas supplémentaires.

\paragraph{\textsc{Addr-User} :} % {{{

Alors $e = \eTaint{lv}$. On applique l'hypothèse de récurrence à $lv$.

\begin{itemize}
\item $lv = φ$. Alors on peut appliquer \textsc{Phi-User}.
\item $\mm{m}{lv}{m'}{lv'}$. On conclut en utilisant \textsc{Ctx} avec $C =
\eTaint{\ctxEmpty}$.
\item $\msi{m}{lv} → Ω$. On applique \textsc{Eval-Err} avec ce même $C$.
\end{itemize}

% }}}
\paragraph{\textsc{GetU} :} % {{{

On applique l'hypothèse de récurrence à $e_d$.

\begin{itemize}
\item $e_d = v_d$. On applique l'hypothèse de récurrence à $e_s$.

\begin{itemize}
\item $e_s = v_s$.

    D'après le lemme~\ref{lemma:canon} adapté aux extensions noyau, $v_s$ a pour
    forme $φ_s$.

    On distingue la forme de $φ_s$:

    \begin{itemize}
        \item $φ_s = \vTainted{φ}$. Alors on applique
        \textsc{User-Get-OK}. Le lemme~\ref{lemma:mem-typ} adapté aux extensions
        noyau assure que les prémisses sont correctes.

        \item $∄~φ, φ_s = \vTainted{φ}$. Alors on applique
        \textsc{User-Get-Err}.
    \end{itemize}

\item $\mm{m}{e_s}{m'}{e_s'}$. Posons $C = \uGet{v_d}{\ctxEmpty}$.
    On conclut avec \textsc{Ctx}.
\item $\msi{m}{e_s} → Ω$. Idem avec \textsc{Eval-Err}.
\end{itemize}

\item $\mm{m}{e_d}{m'}{e_d'}$. On applique \textsc{Ctx} avec
    $C = \uGet{\ctxEmpty}{e_s}$.

\item $\msi{m}{e_d} → Ω$. On utilise \textsc{Eval-Err} avec ce même contexte.

\end{itemize}
% }}}
\paragraph{\textsc{PutU} :} % {{{
Ce cas est similaire au cas \textsc{GetU}, en appliquant les règles
\textsc{User-Put-OK} et \textsc{User-Put-Err}.

% }}}

\deleted{Suppression du cas addr-kernel inutile}

\end{proof}

\section{Préservation pour les extensions noyau}
\label{proof:preservation-qualif}

(Théorème~\ref{thm:preservation-qualif})

De même, il suffit de prouver les cas correspondant aux nouvelles règles.

\paragraph{\textsc{Phi-User}:} % {{{

On applique le lemme de représentation, qu'on a étend avec le cas
\linebreak
$\cRepr{t~@} = \cRepr{\cRepr{t}~@}$.

% }}}
\paragraph{\textsc{User-Get-OK}:} % {{{
Tout d'abord, d'après le lemme~\ref{lemma:inv-qualif}, $t = \tInt$, donc la
préservation du type est établie car $m' ⊧ 0 : \tInt$
La compatibilité mémoire est obtenue en appliquant \textsc{M-Write}.
% }}}
\paragraph{\textsc{User-Get-Err}:} % {{{
La seule partie à prouver est la préservation, qui se fait de la même manière
que dans le cas précédent.
% }}}
\paragraph{\textsc{User-Put-OK}:} % {{{
Idem que dans le cas {\textsc{User-Put-OK}.
% }}}
\paragraph{\textsc{User-Put-Err}:} % {{{
Idem que pour \textsc{User-Get-Err}.
% }}}

% vim: spelllang=fr


\chapter{Analyse de bitmasks}

\label{anx:bitmasks}.
Par exemple, en ouvrant un fichier on précise le mode du fichier (lecture,
écriture ou les deux) par les bits 1 et 2, s'il faut créer le fichier ou non
s'il n'existe pas par le bit 7, s'il dans ce cas il doit être effacé par le bit
8), etc. On obtient un mode en réalisant un ``ou'' bit à bit entre des
constantes. Ces deux utilisations du type \texttt{int} n'ont rien à voir ; il
faudrait donc empêcher d'utiliser un descripteur de fichier comme un mode, et
vice-versa. De même, aucun opérateur n'a de sens sur les descripteurs de
fichier, mais l'opérateur \texttt{\textbar{}} du ``ou'' bit à bit doit rester
possible pour les modes.


\backmatter{}

%\chapterstar{Liste des figures}

\clearpage{}

\renewcommand\listfigurename{Liste des figures}
\listoffigures{}

\clearpage{}

\renewcommand{\listtheoremname}{Liste des définitions}
\listoftheorems[ignoreall,show={definition}]

\renewcommand{\listtheoremname}{Liste des théorèmes et lemmes}
\listoftheorems[ignoreall,show={theorem,lemma}]

\chapter{Références web}

\insertlinks{}

\bibliographystyle{alpha}
\bibliography{typ}

\begin{headingpage}
\thispagestyle{empty}

\begin{center}\Large \textbf{Résumé}\end{center}

Les noyaux de système d'exploitation manipulent des données fournies par
les programmes utilisateur via certaines interfaces, comme les appels système.
Si celles-ci sont manipulées sans prendre une attention particulière, des
failles de sécurité peuvent apparaître, comme des fuites de données
confidentielles ou l'elévation de privilège d'un attaquant.

Le but de cette thèse est d'utiliser des techniques de typage statique afin de
détecter les manipulations dangereuses de pointeurs contrôlés par l'espace
utilisateur.

De nombreux systèmes d'exploitation, dont le noyau Linux sur lequel nous nous
basons, sont écrits dans le langage C. Afin de pouvoir appliquer les techniques
de typage statique à ce langage, on commence par isoler un sous-ensemble sûr
nommé \langname, pour lequel on définit une sémantique opérationnelle et un
système de types. Les propriétés classiques de sûreté du typage sont établies.

Dans un second temps on étend \langname pour intégrer la notion de valeur
provenant de l'espace utilisateur. On montre que l'extension triviale du système
de types brise la sûreté du typage, et que l'on peut réétablir celle-ci en
donnant un type particulier aux pointeurs dont la valeur est contrôlée par
l'espace utilisateur, forçant leur déférencement à se faire au sein d'un
ensemble limités de fonction sûres. Un cas d'étude est déroulé pour montrer que
cette technique permet de détecter un bug qui a frappé un pilote de carte
graphique AMD dans le noyau Linux.

Ces travaux ouvrent des perspectives sur deux points. D'une part, il serait
possible d'appliquer à de plus grandes parties du noyau Linux. Les diverses
interfaces qui exposent des données contrôlées par l'espace utilisateur peuvent
bénéficier de cette vérification statique. Ceci est d'autant plus vrai que les
pilotes de périphériques dont elles sont issues sont en général moins testées
que le reste du code du noyau. D'autre part, cette technique revient à
introduire des types abstraits opaques, séparant leur représentation de leur
interface. Ce concept étant absent de C, de nombreuses erreurs de programmation
notamment liées aux manipulations d'entiers peuvent également être vérifiées à
l'aide de cette méthode.

\begin{center}\Large \textbf{Abstract}\end{center}

Lorem ipsum dolor sit amet, consectetuer adipiscing elit, sed diam nonummy nibh
euismod tincidunt ut laoreet dolore magna aliquam erat volutpat. Ut wisi enim ad
minim veniam, quis nostrud exerci tation ullamcorper suscipit lobortis nisl ut
aliquip ex ea commodo consequat. Duis autem vel eum iriure dolor in hendrerit in
vulputate velit esse molestie consequat, vel illum dolore eu feugiat nulla
facilisis at vero eros et accumsan et iusto odio dignissim qui blandit praesent
luptatum zzril delenit augue duis dolore te feugait nulla facilisi. Nam liber
tempor cum soluta nobis eleifend option congue nihil imperdiet doming id quod
mazim placerat facer possim assum. Typi non habent claritatem insitam; est usus
legentis in iis qui facit eorum claritatem. Investigationes demonstraverunt
lectores legere me lius quod ii legunt saepius. Claritas est etiam processus
dynamicus, qui sequitur mutationem consuetudium lectorum. Mirum est notare quam
littera gothica, quam nunc putamus parum claram, anteposuerit litterarum formas
humanitatis per seacula quarta decima et quinta decima. Eodem modo typi, qui
nunc nobis videntur parum clari, fiant sollemnes in futurum

%\begin{center}\Large \textbf{Abstract}\begin{center}

% TODO (PLAS suit)

%Manipulating user-provided pointers in the kernel of an operating system can
%lead to security flaws if done in an incautious manner. We present an efficient
%system to detect and prevent this class of erroneous memory manipulation.

%At the core of our approach is \langname, an imperative language that we equip
%with a qualified type system, where two kinds of pointers are distinguished:
%\emph{safe} pointers, whose value is statically proved to be controlled by the
%kernel, and \emph{unsafe} ones, whose value comes from userspace through
%run-time system calls. Dereferencing unsafe pointers is forbidden in a static
%manner by the means of a strong type system.

%A concrete case study is described based on a bug that affected a video driver
%in the Linux kernel. We also explain a technique to automatically translate GNU
%C code to our core language, which will enable us to analyze larger fractions of
%the kernel in order to find similar vulnerabilities.

\end{headingpage}


\end{document}
